% BEGIN EMBEDDED BIBLIOGRAPHY
% Bibliography data are part of this source; the compiler writes its temporary input.
\begin{filecontents*}[overwrite]{\jobname.bib}

@article{FILS1978,
  author = {Fr\"ohlich, J. and Israel, R. and Lieb, E. H. and Simon, B.},
  title = {Phase transitions and reflection positivity.
    {I}. General theory and long range lattice models},
  journaltitle = {Communications in Mathematical Physics},
  volume = {62}, date = {1978}, pages = {1--34},
  doi = {10.1007/BF01940327}
}
@online{PongeSchatten2025,
  author = {Ponge, R.},
  title = {Singular values and {Schatten} classes},
  organization = {Introduction to Noncommutative Geometry},
  date = {2025}, note = {Chapter 4},
  url = {https://raphaelponge.org/wp-content/uploads/2025/07/summer-chap4-1.pdf}
}
@online{Sorensen,
  author = {Sorensen, E.},
  title = {{TA} session: tightness of probability measures},
  organization = {SLMath Stochastic Quantization Summer School},
  date = {2024-07-01}, note = {Theorem 1.1},
  url = {https://www.mit.edu/~skycao/SLMath/Session_1_notes_(tightness).pdf}
}
@online{Peterson2017,
  author = {Peterson, J.}, title = {Real analysis},
  organization = {Vanderbilt University}, date = {2017},
  note = {Theorem 6.2.9},
  url = {https://math.vanderbilt.edu/peters10/teaching/spring2017/RealAnalysis.pdf}
}
@online{Potzelberger,
  author = {P\"otzelberger, K.}, title = {Probability: handout},
  organization = {Vienna University of Economics and Business},
  note = {Theorem 6.2},
  url = {https://statmath.wu.ac.at/courses/qfin_prob/Handout.pdf}
}
\end{filecontents*}
% END EMBEDDED BIBLIOGRAPHY
% Document LPM1 (formerly LP1M1). Canonical source; see codenames.json.
\documentclass[11pt,letterpaper,reqno]{amsart}
% BEGIN INLINE manuscript-style
\makeatletter
\NeedsTeXFormat{LaTeX2e}[2020/10/01]
% Use with \documentclass[11pt,letterpaper,reqno]{amsart}.
% This file controls presentation only. It supplies no mathematical content.

\RequirePackage[T1]{fontenc}
\RequirePackage{lmodern}
\RequirePackage{microtype}
\RequirePackage[letterpaper,margin=1in,headheight=15pt]{geometry}
\RequirePackage{amsmath,amsthm,amssymb,mathtools}
\RequirePackage{enumitem}
\RequirePackage{perpage}
\RequirePackage{xcolor}
\definecolor{manuscriptred}{rgb}{0.6,0.1,0.1}
\definecolor{manuscriptblue}{rgb}{0.1,0.3,0.5}
\definecolor{manuscriptmagenta}{HTML}{60445C}
\RequirePackage[breaklinks=true,colorlinks=true,
  linkcolor=manuscriptred,citecolor=manuscriptblue,urlcolor=manuscriptmagenta]{hyperref}
\RequirePackage{cleveref}
\hypersetup{pdfauthor={}}

% Retain the AMS paragraph, title, heading and running-head conventions.
% Equations retain one global Arabic sequence, as in the reference.
\setcounter{tocdepth}{1}
\renewcommand{\thefigure}{\Roman{figure}}
\renewcommand{\thetable}{\Roman{table}}
\MakePerPage{footnote}
\renewcommand{\thefootnote}{%
  \ifnum\value{footnote}>9\arabic{footnote}\else\fnsymbol{footnote}\fi}

% Upright statements. Main results are A, B, C; supporting statements
% share a section-local sequence. These end marks indicate typography,
% never proof completeness, validation, or a deferred-proof status.
\theoremstyle{definition}
\newtheorem{ms@theorem}{Theorem}
\renewcommand{\thems@theorem}{\Alph{ms@theorem}}
\newtheorem*{ms@theoremstar}{Theorem}
\newtheorem{ms@proposition}{Proposition}[section]
\newtheorem{ms@lemma}[ms@proposition]{Lemma}
\newtheorem{ms@definition}[ms@proposition]{Definition}
\newtheorem{ms@corollary}[ms@proposition]{Corollary}
\newtheorem*{ms@corollarystar}{Corollary}

\newcommand{\ms@majorstart}{%
  \pushQED{\qed}\renewcommand{\qedsymbol}{$\blacksquare\square$}}
\newcommand{\ms@minorstart}{%
  \pushQED{\qed}\renewcommand{\qedsymbol}{$\blacksquare$}}
\newenvironment{theorem}[1][]{\ms@majorstart\ms@theorem[#1]}
  {\popQED\endms@theorem}
\newenvironment{theorem*}[1][]{\ms@majorstart\ms@theoremstar[#1]}
  {\popQED\endms@theoremstar}
\newenvironment{proposition}[1][]{\ms@minorstart\ms@proposition[#1]}
  {\popQED\endms@proposition}
\newenvironment{lemma}[1][]{\ms@minorstart\ms@lemma[#1]}
  {\popQED\endms@lemma}
\newenvironment{definition}[1][]{\ms@minorstart\ms@definition[#1]}
  {\popQED\endms@definition}
\newenvironment{corollary}[1][]{\ms@majorstart\ms@corollary[#1]}
  {\popQED\endms@corollary}
\newenvironment{corollary*}[1][]{\ms@majorstart\ms@corollarystar[#1]}
  {\popQED\endms@corollarystar}
% The source's lemma/proposition variants are deliberately not carried over:
% choose one marker for each statement type rather than encode proof status.

\theoremstyle{remark}
\newtheorem{remark}[ms@proposition]{Remark}
\newtheorem*{remark*}{Remark}
% The ordinary amsthm proof environment retains its open square.

\crefname{ms@theorem}{theorem}{theorems}
\Crefname{ms@theorem}{Theorem}{Theorems}
\crefname{ms@proposition}{proposition}{propositions}
\Crefname{ms@proposition}{Proposition}{Propositions}
\crefname{ms@lemma}{lemma}{lemmas}
\Crefname{ms@lemma}{Lemma}{Lemmas}
\crefname{ms@definition}{definition}{definitions}
\Crefname{ms@definition}{Definition}{Definitions}
\crefname{ms@corollary}{corollary}{corollaries}
\Crefname{ms@corollary}{Corollary}{Corollaries}
\crefname{remark}{remark}{remarks}
\Crefname{remark}{Remark}{Remarks}

% Bibliography configuration for this manuscript.
  \RequirePackage[backend=biber,style=alphabetic,backref=true,
    giveninits=true,isbn=false,date=year]{biblatex}
  \addbibresource{\jobname.bib}
  \setcounter{biburlnumpenalty}{100}
  \setcounter{biburllcpenalty}{100}
  \setcounter{biburlucpenalty}{100}
  \renewcommand*{\bibfont}{\footnotesize}
  \renewbibmacro{in:}{}
  \DeclareFieldFormat{pages}{#1}
  % Use the public alphabetic style. No private author-hash overrides;
  % repeated author names are printed in full.
\makeatother
% END INLINE manuscript-style
\newcommand{\R}{\mathbb R}
\newcommand{\C}{\mathbb C}
\newcommand{\Z}{\mathbb Z}
\newcommand{\N}{\mathbb N}
\newcommand{\E}{\mathbb E}
\newcommand{\dd}{\mathrm d}
\DeclareMathOperator{\supp}{supp}
\DeclareMathOperator{\tr}{tr}
\DeclareMathOperator{\diag}{diag}
\allowdisplaybreaks[2]
\title[Reflection positivity on the infinite lattice]{Reflection positivity for chamber QED on the infinite lattice}
\date{}
\hypersetup{pdftitle={Reflection positivity for chamber QED on the infinite lattice},pdfauthor={}}
\begin{document}
\begin{abstract}
We consider the real curvature field of three-dimensional lattice QED
with chamber fermions and all flux sectors included. At fixed positive
charge, Wilson parameter in $(0,1)$, positive hopping weights and
arbitrary real mass parameter, the cubic finite-volume laws are tight.
Every limit along arbitrary integer side lengths is invariant under
the lattice symmetries and satisfies sharp Gaussian domination and
reflection positivity at all three coordinate and six diagonal
families of lattice mirrors. The reflection form is nonnegative on
the full $L^2$ space generated by plaquettes in either open side.
The proof compares cubic laws with reflected perturbations on adapted
tori. Local operator estimates and finite layer elimination recover
translation symmetry after the perturbation is removed; unequal slab
folds and convex separation then give positivity for the original law.
\end{abstract}
\maketitle
\tableofcontents
\section{Introduction}\label{sec:introduction}

We fix the lattice spacing and use unit lattice coordinates. The photon
field is a real one-cochain on $\Z^3$. Write $d$ for the lattice
coboundary, orient a plaquette by the ordered pair $(e_j,e_k)$, $j<k$,
and let $\mathcal P_2$ contain the plaquettes with these orientations.
The space in which the volume limit is taken is
\[
 \Omega=\{F\in\R^{\mathcal P_2}:dF=0\},
\]
with the product topology. Reversing a plaquette's orientation negates
its field value. In particular, pullback by a lattice isometry acts on
$F$ as a two-cochain, including the orientation sign.

Here is the finite model. Let
$W=\Z^3\rtimes\bigl(O(3)\cap\operatorname{GL}(3,\Z)\bigr)$ and
$C_0=\{0<x_3<x_2<x_1<1/2\}$. Fix $v_0\in C_0$, and let
$s_0,s_1,s_2,s_3$ be the reflections in its walls
$x_3=0$, $x_2=x_3$, $x_1=x_2$, $x_1=1/2$, respectively.
The fermion sites are $\mathcal X=Wv_0$, and the four neighbors
of $x=gv_0$ are $x^{(i)}=gs_iv_0$. Write $\pi(x)$ for the nearest
point of $\Z^3$ to $x$. With $\gamma_j=\operatorname{diag}(\sigma_j,-\sigma_j)$,
where $\sigma_j$ are the Pauli matrices, put $\gamma_v=\sum_jv_j\gamma_j$.
For charge $q$, Wilson parameter $\rho$, weights $w_i$, and mass
parameter $M$, the chamber operator is
\[
 (\mathcal D_A\psi)(x)
 =M\psi(x)-\sum_{i=0}^3\frac{w_i}{2}
   \left(\rho-\gamma_{(x^{(i)}-x)/|x^{(i)}-x|}\right)
   e^{iqA(\pi(x),\pi(x^{(i)}))}\psi(x^{(i)}),
\]
where $A(y,x)=-A(x,y)$ and $A(x,x)=0$.

On the cubic torus $(\Z/n\Z)^3$, the photons are periodic and
$\mathcal D_A^{C(n)}$ acts on spinors satisfying
$\psi(x+ne_j)=-\psi(x)$ for $j=1,2,3$. Put $r_q=2\pi/q$.
For a bounded Borel function $G$ on $\Omega$, define
\[
 \mu_n(G)=\frac1{Z_{C(n)}[1]}
 \sum_{\substack{z\in C^2((\Z/n\Z)^3;\Z)\\dz=0}}
 \int_{[0,r_q)^{3n^3}}
 G(\widetilde F)\,e^{-\|F\|^2/2}
       \det\mathcal D_A^{C(n)}\,\dd A,
 \qquad F=dA+r_qz.
\]
Here $\widetilde F$ is the periodic extension of $F$, the norm sums
over the torus plaquettes, and $Z_{C(n)}[1]$ is the same expression
without $G$. Thus every closed integer cochain is summed, with no
restriction on its flux. The determinant is nonnegative and
$0<Z_{C(n)}[1]<\infty$; these facts are proved in
\cref{prop:reg-determinant-sign,prop:reg-partition}.

\begin{theorem}[Reflection positivity in infinite volume]\label{thm:main}\label{thm:cubic-limits}\label{thm:diag-lattice-limits}
Fix $q>0$, $0<\rho<1$, $M\in\R$ and $w_0,\ldots,w_3>0$.
The family $(\mu_n)_{n\ge4}$ is tight on $\Omega$. Let
$n_i\ge4$ be any sequence of integers tending to infinity, and let
$\nu$ be any subsequential weak limit of $(\mu_{n_i})$. Then:
\begin{enumerate}
\item[\textup{(i)}] For every finitely supported real two-cochain $\beta$,
\[
 \int_\Omega e^{\langle\beta,F\rangle}\,\dd\nu(F)
       \le e^{\|\beta\|^2/2},
 \qquad
 \langle\beta,F\rangle=\sum_{\sigma\in\mathcal P_2}\beta_\sigma F_\sigma,
 \qquad \|\beta\|^2=\sum_{\sigma\in\mathcal P_2}\beta_\sigma^2.
\]
\item[\textup{(ii)}] The law $\nu$ is invariant under lattice translations,
signed coordinate permutations, and $F\mapsto-F$.
\item[\textup{(iii)}] The law $\nu$ is invariant and reflection positive
at every mirror
\[
 x_j=c\quad(c\in\tfrac12\Z),\qquad
 x_j\pm x_k=c\quad(c\in\Z,\ j<k).
\]
More precisely, write a mirror as $\Pi=\{\ell(x)=c\}$ and let
\[
 \theta_\Pi x
 =x-\frac{2(\ell(x)-c)}{|\nabla\ell|^2}\nabla\ell.
\]
Choose either open side $H$ of $\Pi$, and let $\Sigma_H$ be the
sigma-algebra generated by the values $F_\sigma$ for which every
corner of $\sigma$ lies in $H$. Then
\[
 \int_\Omega\overline{A(\theta_\Pi^*F)}A(F)\,\dd\nu(F)\ge0
 \qquad\bigl(A\in L^2(\Sigma_H,\nu)\bigr).
\]
\end{enumerate}
\end{theorem}

There are three coordinate and six diagonal mirror families in
\textup{(iii)}. The support convention excludes plaquettes on or
crossing the mirror. For example, reflection in $x_1=0$ changes the
signs of $F_{12}$ and $F_{13}$ and preserves the sign of $F_{23}$.
These signs occur both in invariance of the law and in its reflected
form.

\subsection{The comparison argument}

The coordinate mirrors are compatible with every cubic torus. Their
finite-volume reflected forms, together with a Gaussian bound,
give tightness and coordinate reflection positivity of every weak
limit. A diagonal mirror requires a different period lattice.
The task is therefore to obtain diagonal positivity for the limit
of the original cubic laws.

Reflected slab operators give a comparison with a diagonal torus
of the same volume. Insert a real local trigonometric polynomial
$O$ in one cell, reflect it into all cells, and tilt the diagonal
law by the exponential of $t$ times their sum. The resulting law
has reflection positivity at the cell walls. Its expectation of
$O$ bounds the cubic expectation, up to a pressure error divided
by $t$. The pressure error tends to zero as the volume grows, so
the volume limit must precede $t\downarrow0$.

The perturbation occupies the entire torus. Its pointwise bound
does not control local expectations uniformly in volume. We obtain
that control by amplifying local insertions in two cell structures
and applying a coordinate operator bound. Heat regularization then
proves approximate reflection symmetry on two adjacent fermion
layers. Fermion elimination and an imaginary contour shift reduce
each curvature character to these two layers, with an error of
order $t$ independent of the volume. After the ordered limits,
the comparison law is invariant under coordinate link reflections
and hence under translations by $2\Z^3$.

These translations allow a given observable to be placed inside a
cell with any prescribed diagonal wall. For each real local
trigonometric polynomial $O$ we obtain a comparison law $\nu_O$
with positivity at that wall and
$\nu(O)\le\nu_O(O)$. The comparison law may depend on $O$.
The laws satisfying the Gaussian bound and positivity at the wall
form a compact convex set, and separation places $\nu$ itself in
that set.

The repeated reflections require divisible side lengths, whereas
the theorem allows arbitrary integers $n_i$. An unequal-height
slab estimate compares each original cube to a nearby divisible
companion, whose side $m_i$ satisfies $m_i/n_i\to1$. The H\"older-weighted
volume of the exceptional slabs is a vanishing fraction of the original
volume, and a uniform pressure bound controls their contribution. All auxiliary states
are constructed on the companions; the final separation argument
still applies to the prescribed limit of the original cubes.


We give the finite-volume identities and limiting arguments needed
for the fixed-lattice theorem, including the unequal-slab comparison
for arbitrary cubic side lengths and the extension of the reflection
form to the full open-side $L^2$ space. Repeated reflection and the
chessboard method belong to the general theory of Fr\"ohlich, Israel,
Lieb, and Simon~\cite[Section~4]{FILS1978}. Here the slab trace must also
retain the photon flux sum and the fermion boundary character when
the period lattice changes. The local operator and layer-elimination
arguments establish the additional symmetry needed for that comparison.
The theorem concerns fixed lattice parameters and subsequential volume
limits; it asserts neither uniqueness nor a continuum limit.
\section{The chamber regularization}\label{sec:regularization}

The photon field is a real one-cochain on the cubic lattice. The fermion
has one site in each chamber of the lattice mirrors. This choice has a
specific consequence: a fermion bond crossing a mirror joins a site to
its reflected site and is perpendicular to the mirror. The hopping term
then admits a positive reflected expansion. We prove this expansion and
the resulting finite-volume reflection positivity below. The construction
also requires a nonzero partition function when the bare mass has
undetermined sign. For this purpose we identify a nonzero coefficient
of the fermion determinant, independent of the mass. The determinant
sign and positive normalization are proved in
\cref{prop:reg-determinant-sign,prop:reg-partition}.

Throughout this section, $0<\rho<1$, $q>0$, and the mass parameter $M$
is real. The hopping weights $w_0,\ldots,w_3$ are positive. All constants may depend on these fixed parameters unless a uniform
bound is stated.

\subsection{The photon field and the flux sum}\label{subsec:reg-photon}

We orient each link $(x,\mu)$ from $x$ to $x+e_\mu$ and each plaquette
$(x;\mu\nu)$, $\mu<\nu$, by the ordered pair $(e_\mu,e_\nu)$.
The real $k$-cochains are denoted by $C^k$, and the integer cochains by
$C^k(\Z)$. The coboundary $d$ and its adjoint $d^*$ use these
orientations. Thus
\begin{equation}\label{eq:reg-coboundary}
 (dA)_{\mu\nu}(x)=A_\mu(x)+A_\nu(x+e_\mu)
                     -A_\mu(x+e_\nu)-A_\nu(x).
\end{equation}
An oriented link may also be written as a pair of its endpoints;
$A(y,x)=-A(x,y)$ and $A(x,x)=0$. Pullback by a lattice isometry includes
the orientation of the image cell. In particular, pullback commutes
with $d$. For a coordinate reflection, this gives a minus sign on
the curvature components containing the normal direction and a
plus sign on the tangential component. These are the two-form
reflection signs used throughout.

Let $\Gamma\subset\Z^3$ be a rank-three subgroup with
\[
 \min\{|g|_\infty:g\in\Gamma\setminus\{0\}\}\ge4.
\]
We write $C^k(\Gamma)$ for the cochains on $\Z^3/\Gamma$ and identify
them with their periodic extensions. The torus has
$V=[\Z^3:\Gamma]$ photon sites, $3V$ links and $3V$ plaquettes. The
inner product on real cochains is the sum over positively oriented
cells, and $\|\cdot\|$ denotes its norm.

A photon configuration is a class in
\begin{equation}\label{eq:reg-configuration}
 \operatorname{Conf}(\Gamma)
 =\{(A,z):A\in C^1(\Gamma),\ z\in C^2(\Gamma;\Z),\ dz=0\}/\!\sim,
 \qquad
 (A,z)\sim\left(A+\frac{2\pi}{q}k,z-dk\right),
\end{equation}
where $k\in C^1(\Gamma;\Z)$. Its curvature and link phases are
\begin{equation}\label{eq:reg-curvature}
 F=dA+\frac{2\pi}{q}z,
 \qquad u_b=\exp(iqA(b)).
\end{equation}
Both are independent of the representative, and $dF=0$. The integer
cochain is required to be closed; there is no restriction on its flux.
The curvature therefore remains a real, unbounded field.

Put $r_q=2\pi/q$. Each class has a unique representative with
$A(b)\in[0,r_q)$ at every link. The measure on the quotient is defined
by
\begin{equation}\label{eq:reg-configuration-integral}
 \int_{\operatorname{Conf}(\Gamma)}\Phi(A,z)\,\dd(A,z)
 =\sum_{\substack{z\in C^2(\Gamma;\Z)\\dz=0}}
       \int_{[0,r_q)^{3V}}\Phi(A,z)\,\dd A.
\end{equation}
This formula holds for nonnegative measurable functions and for
absolutely integrable complex functions. Any measurable fundamental
domain may be used instead. Indeed, the changes of representative
preserve Lebesgue measure in $A$ and counting measure in $z$.

We will repeatedly use the finite number
\begin{equation}\label{eq:reg-gaussian-sum}
 C_q(\varepsilon)=\sup_{y\in\R}\sum_{j\in\Z}
                  e^{-\varepsilon(y+r_qj)^2},\qquad \varepsilon>0.
\end{equation}
It satisfies
\[
 C_q(\varepsilon)\le2+\frac{q}{2\pi}\sqrt{\frac{\pi}{\varepsilon}}.
\]
To see this, retain the two terms nearest the origin and bound each
remaining term by $r_q^{-1}$ times the Gaussian integral over the
adjacent interval of length $r_q$ closer to the origin. Those intervals
have disjoint interiors. Write $C_q=C_q(1/2)$. Dropping the constraint
$dz=0$ in the flux sum gives
\begin{equation}\label{eq:reg-maxwell-integral}
 0<\int_{\operatorname{Conf}(\Gamma)}e^{-\|F\|^2/2}\,\dd(A,z)
 \le(r_qC_q)^{3V}.
\end{equation}
The lower bound follows already from $z=0$; for the upper bound, sum
each $z_\sigma$ separately at fixed $A$.

The periodic extensions of the curvatures lie in
\[
 \Omega=\{F\in C^2:dF=0\},
\]
with the topology of pointwise convergence on plaquettes. This is a
closed subspace of a countable product of copies of $\R$. Local
observables depend on finitely many plaquette values. For a finitely
supported real two-cochain $\beta$, the character
$\exp(i\langle\beta,F\rangle)$ is bounded and continuous on $\Omega$.

\subsection{Chambers and the fermion graph}\label{subsec:reg-chambers}

The lattice mirrors are
\begin{equation}\label{eq:reg-mirrors}
 \{x_\mu=c\},\quad c\in\tfrac12\Z,
 \qquad
 \{x_\mu\pm x_\nu=c\},\quad c\in\Z,\quad\mu<\nu.
\end{equation}
They form three coordinate families and six diagonal families. Let
$\mathcal L$ be the corresponding linear forms $\ell$ and write
$\ell(x)=n_\ell\cdot x$. The admissible positions are
$\operatorname{Adm}(x_\mu)=\frac12\Z$ and
$\operatorname{Adm}(x_\mu\pm x_\nu)=\Z$. The reflection in a mirror
$\Pi=(\ell,c)$ is
\begin{equation}\label{eq:reg-reflection}
 \theta_\Pi x=x-\frac{2(\ell(x)-c)}{|n_\ell|^2}n_\ell.
\end{equation}
Its upward normal is $n_\ell/|n_\ell|$. Coordinate mirrors at integer
positions are called site planes, and those at half-integer positions
are called link planes.

The group generated by these reflections is
\[
 W=\Z^3\rtimes O_{\mathrm{cub}},
 \qquad O_{\mathrm{cub}}=O(3)\cap\operatorname{GL}(3,\Z).
\]
Here $O_{\mathrm{cub}}$ consists of the $48$ signed permutation matrices.
A fundamental chamber is
\[
 C_0=\{0<x_3<x_2<x_1<\tfrac12\}.
\]
Its wall reflections are
\begin{align*}
 s_0(x)&=(x_1,x_2,-x_3),&
 s_1(x)&=(x_1,x_3,x_2),\\
 s_2(x)&=(x_2,x_1,x_3),&
 s_3(x)&=(1-x_1,x_2,x_3).
\end{align*}
Denote the walls by $\Pi_i$ and their inward unit normals by
\begin{equation}\label{eq:reg-wall-normals}
 n_0=e_3,\quad n_1=(e_2-e_3)/\sqrt2,
 \quad n_2=(e_1-e_2)/\sqrt2,\quad n_3=-e_1.
\end{equation}

Fix $v_0=(v_1,v_2,v_3)\in C_0$. The fermion sites and their neighbors
are
\begin{equation}\label{eq:reg-fermion-sites}
 \mathcal X=Wv_0,\qquad (gv_0)^{(i)}=gs_iv_0.
\end{equation}
The unordered bond $\{x,x^{(i)}\}$ has type $i$. A site has the unique
form $x=k+rv_0$, with $k\in\Z^3$ and $r\in O_{\mathrm{cub}}$; put
$\pi(x)=k$ and $y(x)=rv_0$. Thus $\pi(x)$ is its nearest photon site.

\begin{proposition}\label{prop:reg-geometry}
The chambers are precisely $gC_0$, $g\in W$, and $W$ acts on them
simply transitively. No fermion site lies on a mirror. Every bond
crosses exactly one mirror, its chamber wall, and crosses it
orthogonally at its midpoint. Consequently a bond crossing $\Pi$ is
$\{x,\theta_\Pi x\}$.

The projection $\pi$ is $W$-equivariant. Bonds of types $0,1,2$ have
the same photon site at both ends; a type-$3$ bond joins two nearest
photon sites. Exactly eight fermion bonds lie over each photon link.
If a mirror does not cut the link and its open segment lies on one
side, all eight bonds lie on that side.
The graph is connected and bipartite, and each site has one neighbor
of each type.
\end{proposition}

\begin{proof}
Coordinate reflections and coordinate transpositions generate
$O_{\mathrm{cub}}$, and
$\theta_{x_\mu,1/2}\theta_{x_\mu,0}=\tau_{e_\mu}$ gives the
translations. Conversely, \eqref{eq:reg-reflection} is an affine signed
permutation with integer translation, so the mirror group is $W$.

For a point on no mirror, subtract the unique nearest lattice point.
The resulting three coordinates are nonzero, have distinct absolute
values, and lie in $(-1/2,1/2)$. There is exactly one signed
permutation making them positive and decreasing. This places the
point in exactly one translate of $C_0$. These translates are disjoint
open convex sets, meet no mirror, and cover its complement. They are
therefore the chambers, with the asserted free transitive action.

For the four bonds from $v_0$, the midpoint is, respectively,
\[
 (v_1,v_2,0),\quad
 (v_1,(v_2+v_3)/2,(v_2+v_3)/2),\quad
 ((v_1+v_2)/2,(v_1+v_2)/2,v_3),\quad
 (1/2,v_2,v_3).
\]
Each belongs to the relative interior of its wall and to no other
mirror. Each half of the bond, except its midpoint, lies in one of
the two adjacent chambers. Transport by $g\in W$ proves the bond
assertion at every site.

In the coordinates $x=k+rv_0$, the first three neighbors are
$k+rs_iv_0$, whereas the fourth is
$k+re_1+rr_{e_1}v_0$, where $r_{e_1}$ changes the sign of the first
coordinate. This proves the statements about $\pi$ and its
equivariance. Over $(k,\mu)$ the lower ends are exactly the eight
sites for which $y_\mu=v_1$: the remaining two absolute values can
be ordered in two ways, with four choices of their signs. Their
type-$3$ neighbors are the upper ends of these eight bonds.

The bond over $b=(k,\mu)$ has the segment
$k+te_\mu+y_\perp$, $v_1\le t\le1-v_1$, where the other two
coordinates of $y_\perp$ have absolute values $v_2,v_3$. For a
diagonal height involving $x_\mu$, its varying part is
$t\pm y_\nu\in(0,1)$, the same interval between integers as on the
open photon link. For heights not involving $x_\mu$, the changes
are $y_\nu$ or $y_\nu\pm y_{\nu'}$, whose absolute values are
smaller than $1/2$ or $1$, respectively. Away from a mirror
containing $b$ these preserve its side. The coordinate height
$x_\mu$ gives the same side except at the unique link plane through
the midpoint, which cuts $b$. This proves the assertion about
photon links.

The generators $s_0,s_1,s_2$ generate the signed permutations, and
$s_3r_{e_1}=\tau_{e_1}$; conjugation gives all coordinate translations.
Thus the four wall reflections generate $W$, proving connectivity.
The two parts are the sites $gv_0$ with $\det r=1$ and with
$\det r=-1$, since each neighbor changes this determinant.
\end{proof}

The geometric location records reflection, while $(k,r)$ labels the
fermion sites by a photon site and a signed permutation. If
$d_i=\operatorname{dist}(v_0,\Pi_i)$, then
\[
 x^{(i)}-x=-2d_i\,rn_i,
 \qquad \frac{x^{(i)}-x}{|x^{(i)}-x|}=-rn_i.
\]
The directions, bond types and projected endpoints in these
coordinates do not depend on $v_0$. Hence neither the operator below
nor its determinant depends on $v_0$ after this identification.
\subsection{The chamber operator and the finite-volume law}
\label{subsec:reg-law}

Let $\sigma_\mu$ be the Pauli matrices and
\[
 \gamma_\mu=\diag(\sigma_\mu,-\sigma_\mu),\qquad
 \gamma_v=\sum_\mu v_\mu\gamma_\mu.
\]
They form a Hermitian Clifford triple. For a unit vector $n$, put
$P_n^\pm=(1\pm\gamma_n)/2$. Both projections have rank two, and
\begin{equation}\label{eq:reg-clifford-reflection}
 \gamma_n\gamma_v\gamma_n=-\gamma_{r_nv},\qquad
 r_nv=v-2(v\cdot n)n.
\end{equation}
This follows by multiplying
$\gamma_n\gamma_v+\gamma_v\gamma_n=2n\cdot v$ by $\gamma_n$.

For an oriented bond $(x,x')$ of type $i$ define
\begin{equation}\label{eq:reg-hop}
 U_A(x,x')=\exp(iqA(\pi(x),\pi(x'))),\qquad
 w(x,x')=\frac{w_i}{2}\left(\rho-\gamma_{(x'-x)/|x'-x|}\right).
\end{equation}
The chamber operator is
\begin{equation}\label{eq:reg-dirac}
 (\mathcal D_A\psi)(x)=M\psi(x)-\sum_{i=0}^3
                   w(x,x^{(i)})U_A(x,x^{(i)})\psi(x^{(i)}).
\end{equation}
Its lattice mass is
\begin{equation}\label{eq:reg-lattice-mass}
 m=M-\frac\rho2\sum_iw_i.
\end{equation}
A fermion boundary condition is a character
$\chi:\Gamma\to\{1,-1\}$. Write $\mathbb T=(\Gamma,\chi)$.
Its spinors satisfy $\psi(x+g)=\chi(g)\psi(x)$ for $g\in\Gamma$.
The restriction $\mathcal D_A^{\mathbb T}$ acts on a space of dimension $192V$,
with the inner product obtained by summing over representatives of
$\mathcal X/\Gamma$. For each site introduce generators
$\psi_\alpha(x),\bar\psi_\alpha(x)$, $1\le\alpha\le4$, obeying the same
boundary conditions. Their Grassmann algebra is denoted by
$\operatorname{Gr}_{\mathbb T}$. Its Berezin integral is normalized by
\begin{equation}\label{eq:reg-berezin}
 \int_{\mathbb T}\prod_{x\in\mathcal X/\Gamma}
                  \prod_{\alpha=1}^4\psi_\alpha(x)\bar\psi_\alpha(x)=1.
\end{equation}
Products of pairs are even, so their ordering is immaterial. With
$S_{\mathbb T}(A)=\sum_x\bar\psi(x)(\mathcal D_A\psi)(x)$, this normalization
gives $\int_{\mathbb T}e^{-S_{\mathbb T}(A)}=\det \mathcal D_A^{\mathbb T}$.

The unnormalized functional is
\begin{equation}\label{eq:reg-functional}
 Z_{\mathbb T}[G]=\int_{\operatorname{Conf}(\Gamma)}
      e^{-\|F\|^2/2}\int_{\mathbb T}e^{-S_{\mathbb T}(A)}G(A,z)
      \,\dd(A,z).
\end{equation}
For a scalar $G$ the inner integral is
$G(A,z)\det \mathcal D_A^{\mathbb T}$. For a Grassmann-valued integrand,
absolute convergence is understood coefficient by coefficient.
Bounded measurable coefficients suffice: $e^{-S_{\mathbb T}(A)}$
has finitely many coefficients, each a polynomial in the link phases,
and \eqref{eq:reg-maxwell-integral} supplies an integrable majorant.

\begin{proposition}\label{prop:reg-determinant-sign}
For every real $A$ and $M$,
\[
 \det \mathcal D_A^{\mathbb T}\ge0,
 \qquad
 \|\mathcal D_A^{\mathbb T}\|\le |M|+\frac{1+\rho}{2}\sum_iw_i.
\]
Moreover $\Re \mathcal D_A^{\mathbb T}=m+\frac\rho2\nabla_A^*\nabla_A$, where
\[
 (\nabla_A\psi)(x,i)=\sqrt{w_i/2}\,
           \bigl(\psi(x)-U_A(x,x^{(i)})\psi(x^{(i)})\bigr).
\]
The determinant is not identically zero as a function of the link
phases, for every real mass parameter $M$.
\end{proposition}

\begin{proof}
Write
\[
 H_A\psi(x)=\sum_iw_iU_A(x,x^{(i)})\psi(x^{(i)}),\qquad
 J_A\psi(x)=\sum_iw_i\gamma_{(x^{(i)}-x)/|x^{(i)}-x|}
                         U_A(x,x^{(i)})\psi(x^{(i)}).
\]
Reversing a bond conjugates its phase and reverses its direction.
Reindexing by the involution $x\mapsto x^{(i)}$ therefore gives
$H_A^*=H_A$ and $J_A^*=-J_A$. Hence
\[
 \mathcal D_A=M-\frac\rho2H_A+\frac12J_A,\qquad
 \nabla_A^*\nabla_A=\sum_iw_i-H_A.
\]
This proves the real-part identity. In the splitting
$\C^4=\C^2\oplus\C^2$, the two blocks of $H_A$ coincide, and the two
blocks of $J_A$ have opposite signs. It follows that
\begin{equation}\label{eq:reg-two-blocks}
 \mathcal D_A^{\mathbb T}=\diag(\check{\mathcal D}_A^{\mathbb T},
                                      (\check{\mathcal D}_A^{\mathbb T})^*),
 \qquad \det \mathcal D_A^{\mathbb T}=|\det\check{\mathcal D}_A^{\mathbb T}|^2.
\end{equation}
Each hopping matrix has norm $w_i(1+\rho)/2$, while the associated
phase-weighted permutation of sites is unitary. Their sum gives the norm bound.

For the last assertion we view the determinant as a Laurent
polynomial $P_{\mathbb T}$ in the independent variables $u_b$.
Only type-$3$ bonds contain a variable, with $u_b$ at the lower end
and $u_b^{-1}$ at the upper end. Fix a coordinate direction $\mu$.
The coefficient of
\begin{equation}\label{eq:reg-extreme-monomial}
 \prod_{b\parallel e_\mu}u_b^{32}
\end{equation}
will be computed without using $M$.

Partition the sites modulo $\Gamma$ into six layers according to
the increasing values
\[
 y_\mu=-v_1,-v_2,-v_3,v_3,v_2,v_1.
\]
Each layer has $8V$ sites. Each site has one upward neighbor,
one downward neighbor, and two neighbors in its own layer. The
upward map is a permutation of the sites and visits successive layers
modulo six. Its bond types are
\begin{equation}\label{eq:reg-upward-types}
 (i_1,i_2,i_3,i_4,i_5,i_6)=(2,1,0,1,2,3).
\end{equation}
For layers $1,\ldots,5$ the upward phase is one; in layer six it is
$u_b$ for the link $(\pi(x),\mu)$. These facts follow by changing the
smallest sign or exchanging successive absolute values in the
neighbor formulas above. The sixth step changes $v_1$ to $-v_1$
and increases the photon coordinate by one.

There are eight lower and eight upper fermion ends over each photon
link, each carrying four matrix rows. An exponent $32$ at a link
therefore requires every lower row to take the upward type-$3$
block and every upper row to avoid the reverse block. Imposing this
at all links in direction $\mu$ forces every row in every layer to
take its upward block. Indeed, layer-six rows fill all columns in
layer one. A layer-one row cannot then use a column in its own layer,
and the reverse block is excluded by the exponent condition; it must
go upward to layer two. Inductively the columns in the current and
preceding layers are already occupied, so layers two through five
must also go upward. The selected terms use no other link variables.

The remaining freedom is a permutation of the four spin indices at
each site. The site permutation, repeated four times, has sign one.
Summing over the spin permutations therefore gives the product of
the $4\times4$ determinants of the selected hopping blocks. The
boundary signs have fourth power one, and
\[
 \det w(x,x^\uparrow)
   =(w_{i_k}/2)^4(1-\rho^2)^2
\]
at a site in layer $k$. Consequently the coefficient of
\eqref{eq:reg-extreme-monomial} is $C_*^V$, where
\begin{equation}\label{eq:reg-extreme-coefficient}
 C_*=(w_0/2)^{32}(w_1/2)^{64}(w_2/2)^{64}(w_3/2)^{32}
                                           (1-\rho^2)^{96}>0.
\end{equation}
This coefficient is independent of $M$. In particular the determinant
polynomial, and its restriction to the phase torus, are nonzero.
\end{proof}

The real-part estimate alone would require control of the sign of
$m$ to prove invertibility. The coefficient computation instead
proves that the determinant has positive integral at every real $M$,
which is all that is needed to define the law.

\begin{proposition}\label{prop:reg-partition}
The partition function satisfies $0<Z_{\mathbb T}[1]<\infty$.
More precisely,
\begin{equation}\label{eq:reg-partition-bounds}
 r_q^{3V}e^{-6r_q^2V}C_*^V
 \le Z_{\mathbb T}[1]
 \le(r_qC_q)^{3V}\left(|M|+\frac{1+\rho}{2}\sum_iw_i\right)^{192V}.
\end{equation}
Thus $\E_{\mathbb T}[G]=Z_{\mathbb T}[G]/Z_{\mathbb T}[1]$ defines a
probability law on scalar observables. Its curvature law
$\mu_{\mathbb T}$ is the law of the periodic extension $\widetilde F$
on $\Omega$.
\end{proposition}

\begin{proof}
For a nonnegative Laurent polynomial $P$ on the phase torus, every
Fourier coefficient satisfies
\[
 |\widehat P(\lambda)|
 =\left|\int P(u)\overline{u^\lambda}\,\dd u\right|
 \le\int P(u)\,\dd u,
\]
where Haar measure is normalized. By the preceding proposition and
\eqref{eq:reg-extreme-coefficient},
\[
 \int_{[0,r_q)^{3V}}\det \mathcal D_A^{\mathbb T}\,\dd A
 \ge r_q^{3V}C_*^V.
\]
On this box $|(dA)_\sigma|\le2r_q$ and
$\|dA\|^2\le12Vr_q^2$. The $z=0$ term in
\eqref{eq:reg-functional} gives the lower bound in
\eqref{eq:reg-partition-bounds}. For the upper bound use
$|\det \mathcal D_A^{\mathbb T}|\le\|\mathcal D_A^{\mathbb T}\|^{192V}$ and
\eqref{eq:reg-maxwell-integral}. The density
$e^{-\|F\|^2/2}\det \mathcal D_A^{\mathbb T}$ is nonnegative. After division
by its strictly positive integral it is a probability density, whose
image under $(A,z)\mapsto\widetilde F$ is the asserted law.
\end{proof}

\section{The reflected form on an adapted torus}\label{sec:finite-reflection}

The chamber bonds crossing a mirror join reflected endpoints. Their
Clifford projections give positive fermion coefficients. We combine
this expansion with a Fourier representation of the integer closedness
constraints to obtain a positive form, including all flux sectors.
The positive representation keeps the contributions from both
reflection seams explicitly.

\subsection{Symmetries and reflected Grassmann products}
\label{subsec:reg-symmetries}

The field strength transforms as a two-cochain, with the orientation
of each plaquette included. For an antiperiodic cubic torus the law is invariant
under $W$ and under $F\mapsto-F$. We record the corresponding
fermionic identities, since reflection positivity uses their signs.

Gauge covariance follows directly from
\[
 U_{A-d\lambda}(x,x')
 =e^{iq\lambda(\pi(x))}U_A(x,x')e^{-iq\lambda(\pi(x'))}.
\]
For lattice covariance set
$\gamma_5=\left(\begin{smallmatrix}0&1_2\\1_2&0\end{smallmatrix}\right)$.
It anticommutes with the Clifford triple, and
\[
 (\gamma_n\gamma_5)\gamma_v(\gamma_n\gamma_5)^{-1}
 =\gamma_{r_nv}.
\]
If the linear part of $g\in W$ is $r_{n_1}\cdots r_{n_k}$, the
corresponding product of these spin matrices, followed by the
permutation $x\mapsto g^{-1}x$, intertwines $\mathcal D_A^{\mathbb T}$ and
$\mathcal D_{(g^{-1})^*A}^{g\mathbb T}$, where
$g\mathbb T=(r\Gamma,\chi\circ r^{-1})$ and $r$ is the linear part
of $g$. Taking determinants and changing
variables in \eqref{eq:reg-functional} proves
\[
 \mu_{g\mathbb T}=((g^{-1})^*)_*\mu_{\mathbb T}.
\]
The configuration measure is preserved because pullback permutes
cochain coordinates up to signs and maps one fundamental domain
to another. All translations preserve $\mathbb T$, and all elements
of $W$ preserve an antiperiodic cubic torus, proving the stated
invariance.

For charge conjugation let $C=\sigma_2\oplus\sigma_2$. The Pauli
identity $C\gamma_v^\top C^{-1}=-\gamma_v$, together with reversal
of the oriented bond, gives
\[
 C(\mathcal D_A^{\mathbb T})^\top C^{-1}=\mathcal D_{-A}^{\mathbb T}.
\]
Thus the determinant is unchanged by $A\mapsto-A$. The substitution
$(A,z)\mapsto(-A,-z)$ preserves the configuration measure and sends
$F$ to $-F$, proving charge symmetry of the curvature law.

At a mirror $\Pi$ with upward normal $n$, the Grassmann reflection
$\Theta_\Pi$ is antilinear and reverses products. In components it is
\begin{equation}\label{eq:reg-grassmann-reflection}
 \Theta_\Pi\psi_\alpha(x)
  =-\sum_\beta\bar\psi_\beta(\theta_\Pi x)(\gamma_n)_{\beta\alpha},
 \qquad
 \Theta_\Pi\bar\psi_\alpha(x)
  =-\sum_\beta(\gamma_n)_{\alpha\beta}\psi_\beta(\theta_\Pi x).
\end{equation}
On scalar coefficient functions it acts as
$\Phi(A,z)\mapsto\overline{\Phi(\theta_\Pi^*A,\theta_\Pi^*z)}$.
The Clifford identities give $\Theta_\Pi^2=1$ and
\begin{equation}\label{eq:reg-reflected-bilinear}
 \Theta_\Pi(\bar\psi(x)Y\psi(x'))
 =\bar\psi(\theta_\Pi x')\gamma_nY^*\gamma_n\psi(\theta_\Pi x).
\end{equation}
For a hopping term,
$\gamma_nw(x,x')\gamma_n=w(\theta_\Pi x',\theta_\Pi x)$, by
\eqref{eq:reg-clifford-reflection}. Reflection therefore permutes the
oriented terms of the action, including the real mass terms.

If $Y$ is a finite set of sites, then
\begin{equation}\label{eq:reg-reflected-berezin}
 \int_{\theta_\Pi Y}\Theta_\Pi\xi=\overline{\int_Y\xi}.
\end{equation}
Only the top-degree coefficient needs to be checked. At one site
reflection maps $\prod_\alpha\psi_\alpha\bar\psi_\alpha$ to the same
product at the reflected site, multiplied by $(\det\gamma_n)^2=1$.
Products from different sites are even. This proves
\eqref{eq:reg-reflected-berezin}, and a measure-preserving change of
photon variables gives
\begin{equation}\label{eq:reg-reflected-functional}
 Z_{\theta_\Pi\mathbb T}[\Theta_\Pi G]
 =\overline{Z_{\mathbb T}[G]}
\end{equation}
for bounded coefficient integrands, and for the integrable
Gaussian-growth classes considered below.

On even integrands, compositions of these reflections define an
action of $W$. The ambiguity of a reflection word acts only by the
common sign on all fermion generators. To verify this, compose two
words representing the same lattice isometry. The resulting even
word has trivial geometric action, and its spin matrix is
$\diag(u,u)$ with $u$ a product of an even number of Pauli matrices
$\sigma_n$. It commutes with all Pauli matrices, hence is scalar;
its determinant is one, so $u=\pm1_2$. Its action on an even monomial
is consequently the identity. We write $G^g$ for this action when
folding even integrands. For a real scalar observable it is
$G^g(F)=G(g^*F)$.

\subsection{The fermion crossing at a mirror}\label{subsec:reg-crossing}

Let a bond of type $j$ cross $\Pi$, with lower endpoint $x$ and upper
endpoint $x'=\theta_\Pi x$. Its direction is the upward normal $n$.
The two hopping matrices are
\begin{equation}\label{eq:reg-crossing-projectors}
 \begin{split}
 w(x,x')&=\tfrac{w_j}{2}\bigl((1+\rho)P_n^--(1-\rho)P_n^+\bigr),\\
 w(x',x)&=\tfrac{w_j}{2}\bigl((1+\rho)P_n^+-(1-\rho)P_n^-\bigr).
 \end{split}
\end{equation}
Choose orthonormal bases $e_{n,1}^\pm,e_{n,2}^\pm$ of the two
projected spaces. For each $k=1,2$ introduce four crossing modes.
Their coefficients and phases are displayed together:
\begin{equation}\label{eq:reg-crossing-modes}
\begin{array}{c|c|c|c}
 &\zeta&\omega&u\\ \hline
 1&(e_{n,k}^-)^*\psi(x')&\tfrac{w_j}{2}(1+\rho)&U_A(x,x')\\
 2&(e_{n,k}^+)^*\psi(x')&\tfrac{w_j}{2}(1-\rho)&U_A(x,x')\\
 3&\bar\psi(x')e_{n,k}^+&\tfrac{w_j}{2}(1+\rho)&\overline{U_A(x,x')}\\
 4&\bar\psi(x')e_{n,k}^-&\tfrac{w_j}{2}(1-\rho)&\overline{U_A(x,x')}
\end{array}
\end{equation}
All coefficients are positive in the range $0<\rho<1$.

\begin{proposition}\label{prop:reg-crossing-expansion}
Let $J$ be a finite collection of bonds crossing $\Pi$, and let
$S_\Pi^\times(J)$ be their part of the action. Index all their crossing
modes by a finite set $\mathcal I$. Then
\begin{equation}\label{eq:reg-crossing-positive}
 -S_\Pi^\times(J)=\sum_{i\in\mathcal I}\omega_i u_i\Theta_\Pi(\zeta_i)\zeta_i.
\end{equation}
Choose an order on $\mathcal I$ and, for $I\subset\mathcal I$, put
$\zeta_I=\prod_{i\in I}\zeta_i$, $\omega_I=\prod_{i\in I}\omega_i$
and $u_I=\prod_{i\in I}u_i$. Then
\begin{equation}\label{eq:reg-crossing-exponential}
 e^{-S_\Pi^\times(J)}=\sum_{I\subset\mathcal I}\omega_Iu_I\Theta_\Pi(\zeta_I)\zeta_I.
\end{equation}
At site and diagonal mirrors all $u_i$ are one. At a link plane,
groups one and two have charge $+1$, groups three and four charge
$-1$, and their phases are $e^{iq\mathfrak q(i)A(b)}$ at the cut
photon link $b$.
\end{proposition}

\begin{proof}
The reflections of the four modes in \eqref{eq:reg-crossing-modes}
are, respectively,
\[
 \bar\psi(x)e_{n,k}^-,\quad-\bar\psi(x)e_{n,k}^+,
 \quad-(e_{n,k}^+)^*\psi(x),\quad(e_{n,k}^-)^*\psi(x).
\]
Multiplying each reflected mode by the original mode and using
anticommutation gives the four projected terms in
\eqref{eq:reg-crossing-projectors}, with the indicated phases.
Summation over $k$ uses
$P_n^\pm=\sum_ke_{n,k}^\pm(e_{n,k}^\pm)^*$ and gives the negative of
the two bond terms. Summation over bonds proves
\eqref{eq:reg-crossing-positive}.

Each $\Theta_\Pi(\zeta_i)\zeta_i$ is even and has square zero.
These elements commute. Reversal of the order under $\Theta_\Pi$
gives
\[
 \prod_{i\in I}\Theta_\Pi(\zeta_i)\zeta_i
 =\Theta_\Pi(\zeta_I)\zeta_I,
\]
with no remaining sign. Expanding the product of the finite
exponentials proves \eqref{eq:reg-crossing-exponential}.
Only type-$3$ bonds carry a photon phase, and they cross link
planes. Their projected endpoints are the oriented endpoints of
the cut link, giving the last assertion.
\end{proof}

\subsection{Adapted tori and the second mirror}\label{subsec:reg-adapted}

An adapted torus for $\ell\in\mathcal L$ has the form
\begin{equation}\label{eq:reg-adapted-torus}
 \Gamma=\Gamma_\parallel\oplus\Z v,
 \quad \Gamma_\parallel=\Gamma\cap\{\ell=0\},
 \quad v=\frac{p}{|n_\ell|^2}n_\ell\in\Z^3,
 \quad \chi(v)=-1,
\end{equation}
where $p=\ell(v)>0$ is the height period. For an admissible $c$ set
\[
 c'=c+p/2,\quad \Pi=(\ell,c),\quad\Pi'=(\ell,c'),\quad
 U_+=\{c<\ell<c'\},\quad U_-=\{c-p/2<\ell<c\}.
\]
The two mirrors are admissible: for coordinate families $p\ge4$,
while for diagonal families $p$ is even and $p\ge8$. Their
reflections satisfy $\theta_{\Pi'}=\tau_v\theta_\Pi$ and induce the
same geometric reflection on the torus, exchanging its two halves.
An antiperiodic coordinate torus is adapted to each coordinate
family. The torus with periods $ae_1-ae_2$, $be_1+be_2$ and $ce_3$,
antiperiodic on these periods, is adapted to the corresponding two
diagonal families and the third coordinate family. Cubic images
give adapted tori for all nine families.

Choose representatives $\mathcal X_+$ in $U_+$ and
$\mathcal X_-=\theta_\Pi\mathcal X_+$ in $U_-$. They partition the
fermion sites of the torus. Let $S_+$ contain the mass terms and
oriented bonds entirely in $U_+$, and set
$\theta=\theta_\Pi$, $\Theta=\Theta_\Pi$; the second boundary mirror
induces the same geometric reflection on the torus.
The coefficient functions of $S_+$ involve only links in the closed
upper slab; the reflected part is $\Theta S_+$.

The crossing modes at $\Pi$ lie in
$\operatorname{Gr}(\mathcal X_+)$. At $\Pi'$ their upper endpoints
are outside $U_+$; write
$\check\zeta_{i'}=\Theta_{\Pi'}\zeta_{i'}$, so the checked modes
lie in $\operatorname{Gr}(\mathcal X_+)$. The boundary condition gives
\begin{equation}\label{eq:reg-second-reflection}
 \Theta_{\Pi'}=\chi(v)^{\deg}\Theta_\Pi
                 \quad\hbox{on }\operatorname{Gr}_{\mathbb T}.
\end{equation}
In the crossing term at $\Pi'$ we therefore have
\[
 \Theta_{\Pi'}(\zeta_{i'})\zeta_{i'}
 =-\chi(v)\Theta(\check\zeta_{i'})\check\zeta_{i'}.
\]
The minus sign is from exchanging the two odd factors. Since
$\chi(v)=-1$, it is canceled by antiperiodicity. Both boundary
crossings thus have the same positive coefficients, and
\begin{equation}\label{eq:reg-two-crossings}
 e^{-S_{\mathbb T}}=e^{-S_+}\Theta(e^{-S_+})
 \sum_{I,I'}\omega_I\omega_{I'}u_Iu_{I'}
       \Theta(\zeta_I)\zeta_I
       \Theta(\check\zeta_{I'})\check\zeta_{I'}.
\end{equation}
Here $I$ and $I'$ range over the finite crossing-mode sets at the
two mirrors.

The reflected Berezin integral separates the two halves:
\begin{equation}\label{eq:reg-half-berezin}
 \int_{\mathbb T}\Theta(\xi')\xi
 =\overline{\int_{\mathcal X_+}\xi'}
                  \int_{\mathcal X_+}\xi,
 \qquad \xi,\xi'\in\operatorname{Gr}(\mathcal X_+).
\end{equation}
Indeed, the top monomial is the product of the two half-top
monomials, and reflection maps one to the other by
\eqref{eq:reg-reflected-berezin}. We must still separate the photon
flux constraints; these couple the halves at the cut cubes.

\subsection{The photon crossing and Gauss' law}\label{subsec:reg-gauss}

Let $\mathcal M$ be the two boundary mirrors, understood periodically.
Classify a cell by its relative interior: it is upper or lower if
that interior is in the respective half, lies on the mirrors if the
whole cell does, and is cut if its relative interior meets a mirror
without being contained in it. Write $B_+,B_-,B_0,B_\times$ for the
four sets of links, $P_+,P_-,P_0,P_\times$ for the plaquettes, and
$Q_+,Q_-,Q_\times$ for the cubes.

Reflection exchanges the upper and lower cells. It fixes a cell
lying on the mirrors with orientation sign $+1$, and a cut cell
with orientation sign $-1$. Each cell meets at most one boundary
mirror: its height variation is at most $|n_\ell|^2$, smaller than
$p/2$. To justify the orientation assertion, a point in the
intersection of the relative interior with the mirror is fixed
and lies in the relative interior of the reflected cell. Relative
interiors of different cells are disjoint, so the cell is fixed.
Reflection is the identity on its affine hull in the first case
and a reflection of that affine hull in the cut case.

An upper cell has only upper faces or faces lying on the mirrors;
a cut cell has only upper, lower or cut faces. Moreover a plaquette
containing a cut link is cut, and a cube containing a cut plaquette
is cut. An upper face can touch the boundary at a relative interior
point only when its affine height is constant. A face of a cut cell
cannot lie in the mirror: reflection in the supporting plane of a
face would not preserve that cell. Finally, the height takes both
signs on a cut face, hence on the relative interior of its coface.
In particular, the incidences at a cut plaquette and a cut cube give
\begin{equation}\label{eq:reg-cut-differentials}
 (dA)_\sigma=(d_+A)_\sigma-(d_+\theta^*A)_\sigma
                                      +(d_\times A)_\sigma,
 \qquad
 (dz)(Q)=(d_+z)(Q)-(d_+\theta^*z)(Q)+(d_\times z)(Q),
\end{equation}
where $d_+$ and $d_\times$ retain the incidences from upper and cut
faces, respectively. For example,
$(d_+A)_\sigma=\sum_{b\in B_+}[\sigma:b]A(b)$.
The minus sign in the lower term follows from orientation reversal
of the cut cell and the identity $d\theta^*=\theta^*d$.

Use representatives
\[
 A=(A_+,A_-,A_0,a),\qquad z=(z_+,z_-,z_0,z_\times).
\]
Choose the interval $[0,r_q)$ for upper and mirror links, its signed
reflection for lower links, and $[-r_q/2,r_q/2)$ for cut links.
This is a fundamental domain. Reflection maps the lower variables
measure-preservingly to an independent copy of the upper variables,
fixes $(A_0,z_0)$, and changes the sign of $a$.

The upper and mirror photon weights are
\begin{equation}\label{eq:reg-half-photon-weights}
 W_+=\prod_{\sigma\in P_+}e^{-F_\sigma^2/2}
               \prod_{Q\in Q_+}1_{\{(dz)(Q)=0\}},
 \qquad W_0=\prod_{\sigma\in P_0}e^{-F_\sigma^2/2}.
\end{equation}
The full Maxwell weight with the closedness constraint is
$W_0W_+\Theta(W_+)$ times the Gaussian and constraint factors at
the cut plaquettes and cubes. The half weights contain only upper
and mirror variables. Their integrals are bounded uniformly in the
mirror variables by
\[
 \int W_+\,\dd(A_+,z_+)
       \le r_q^{|B_+|}C_q^{|P_+|},\qquad
 \int W_0\,\dd(A_0,z_0)
       \le r_q^{|B_0|}C_q^{|P_0|}.
\]
The proof drops the cube constraints and uses
\eqref{eq:reg-gaussian-sum} on each plaquette.

At a cut cube impose integer closedness by the Fourier multiplier
\[
 1_{\{(dz)(Q)=0\}}
   =\int_0^{2\pi}e^{i\alpha_Q(dz)(Q)}\frac{\dd\alpha_Q}{2\pi}.
\]
For $\alpha\in[0,2\pi)^{Q_\times}$ and
$k\in\Z^{P_\times}$ define
\begin{equation}\label{eq:reg-dual-crossing-variables}
 \begin{aligned}
 \alpha'_\sigma=\sum_{Q\in Q_\times}[Q:\sigma]\alpha_Q,
 \quad \beta_\sigma=\frac{q}{2\pi}(\alpha'_\sigma-2\pi k_\sigma),
 \\
 u_{\alpha,k}
 =\exp\left(i\sum_Q\alpha_Q(d_+z)(Q)
                  -i\sum_\sigma\beta_\sigma(d_+A)_\sigma\right).
 \end{aligned}
\end{equation}
The function $u_{\alpha,k}$ depends only on the upper variables and
has modulus one.

Poisson summation, in the normalization needed here, reads
\begin{equation}\label{eq:reg-poisson-crossing}
 \sum_{j\in\Z}e^{-(y+r_qj)^2/2+ij\alpha}
 =\frac{q}{\sqrt{2\pi}}\sum_{k\in\Z}
          e^{-\beta_k^2/2-i\beta_ky},
 \qquad\beta_k=\frac{q}{2\pi}(\alpha-2\pi k).
\end{equation}
It follows by Fourier transforming the Gaussian under
$t=y+r_qs$. Insert the cube multipliers and apply this formula at
each cut plaquette. Splitting $dA$ and $dz$ by
\eqref{eq:reg-cut-differentials} gives the exact identity
\begin{align}\label{eq:reg-photon-crossing}
 &\sum_{z_\times}
       \prod_{\sigma\in P_\times}e^{-F_\sigma^2/2}
       \prod_{Q\in Q_\times}1_{\{(dz)(Q)=0\}}
 \nonumber\\
 &\quad=\int\prod_{Q\in Q_\times}\frac{\dd\alpha_Q}{2\pi}
       \sum_{k\in\Z^{P_\times}}
       \prod_{\sigma\in P_\times}
               \frac{q}{\sqrt{2\pi}}e^{-\beta_\sigma^2/2}\,
       u_{\alpha,k}\Theta(u_{\alpha,k})
       \exp\left(iq\sum_{b\in B_\times}a_b(d^*k)(b)\right).
\end{align}
Here $k$ is extended by zero outside $P_\times$. To obtain the final
factor, observe that at a cut link $b$,
\[
 \sum_{\sigma\in P_\times}[\sigma:b]\alpha'_\sigma
 =\sum_{Q\in Q_\times}\alpha_Q
                    \sum_\sigma[Q:\sigma][\sigma:b]=0.
\]
The restriction to cut plaquettes may be dropped because any
plaquette containing a cut link is cut; the last sum is zero by
$d^2=0$. Thus the part of $\beta$ involving $\alpha'$ cancels from
the cut-link factor.

All the weights before the two reflected factors on the right of
\eqref{eq:reg-photon-crossing} are positive. Their total integral is
at most $C_q^{|P_\times|}$. Indeed, with $y=0$ in
\eqref{eq:reg-poisson-crossing},
\[
 \frac{q}{\sqrt{2\pi}}\sum_k e^{-\beta_k^2/2}
 \le\sum_j e^{-(r_qj)^2/2}\le C_q
\]
at each cut plaquette. The original Gaussian sum also has an
absolute bound $C_q^{|P_\times|}$, justifying insertion of the
Fourier multipliers and the exchanges of their sums and integrals.

Cut links occur only at coordinate link planes. A link fixed by a
mirror reflection but not contained in its plane has its endpoints
exchanged; hence it is perpendicular to the mirror, whose normal
is a coordinate vector. A diagonal mirror has cut plaquettes and
cubes, but no cut links.

For crossing-mode subsets $I,I'$ let
$\mathfrak q(I\cup I')(b)$ be the sum of their charges at modes lying
over $b$. The fermionic phases in \eqref{eq:reg-two-crossings} are
\[
 u_Iu_{I'}=\exp\left(iq\sum_{b\in B_\times}
                              a_b\mathfrak q(I\cup I')(b)\right).
\]
Integration of the cut links over one period gives
\begin{equation}\label{eq:reg-gauss-indicator}
 \int_{[-r_q/2,r_q/2)^{B_\times}}
 e^{iq\sum_ba_b((d^*k)(b)+\mathfrak q(I\cup I')(b))}\,\dd a
 =r_q^{|B_\times|}
   1_{\{d^*k+\mathfrak q(I\cup I')=0\text{ on }B_\times\}}.
\end{equation}
The coefficients $(d^*k)(b)+\mathfrak q(I\cup I')(b)$ are integers.
This is precisely
Gauss' law at the cut links, with a nonnegative multiplier. When
there are no cut links the condition is empty.

\subsection{The reflected form on an adapted torus}
\label{subsec:reg-reflected-form}

Let $\mathcal B^+$ be the Grassmann-valued integrands whose
generators are in $\mathcal X_+$ and whose bounded measurable
coefficient functions depend only on the upper and mirror
curvatures and link phases. For the later operator estimates we
also use $\mathcal G^+$, with the same generator support and
coefficient-variable dependence as $\mathcal B^+$, but with boundedness
replaced by the following condition: for some $C<\infty$ and
$\varepsilon>0$, each coefficient $\Phi$ satisfies
\begin{equation}\label{eq:reg-upper-growth}
 W_+|\Phi|\le C\prod_{\sigma\in P_+}e^{-\varepsilon F_\sigma^2}.
\end{equation}
This condition is imposed on the upper and mirror variables and
places no growth factor on mirror plaquettes. The bound is relevant
only where the upper cube constraints hold. Such half-configurations
can be extended to full closed configurations by reflecting their
upper $z$, keeping $z_0$, and taking $z_\times=0$; the cut-cube
identity \eqref{eq:reg-cut-differentials} proves closedness. Thus
these coefficient functions have an unambiguous reading wherever
$W_+\ne0$. In particular $\mathcal B^+\subset\mathcal G^+$.

\begin{proposition}\label{thm:reg-adapted-positivity}
For a torus adapted to $\ell$ and $G,G'\in\mathcal G^+$, the integral
$Z_{\mathbb T}[\Theta(G')G]$ converges absolutely, coefficient by
coefficient, and admits a representation
\begin{equation}\label{eq:reg-positive-representation}
 Z_{\mathbb T}[\Theta(G')G]
   =\int\overline{\widehat{G'}(s)}\widehat G(s)\,\dd\mathfrak m(s),
\end{equation}
where $\mathfrak m$ is a finite positive measure and
$G\mapsto\widehat G$ is linear with bounded measurable values.
This holds for even, odd and mixed-parity integrands. In particular
$Z_{\mathbb T}[\Theta(G)G]\ge0$.

For every bounded scalar observable $\Phi$ depending only on curvatures
of plaquettes wholly contained in $U_+$, the curvature law satisfies
\begin{equation}\label{eq:reg-law-positivity}
 \int_\Omega\overline{\Phi(\theta^*F)}\Phi(F)\,\dd\mu_{\mathbb T}(F)
 \ge0.
\end{equation}
It is invariant under this reflection.
\end{proposition}

\begin{proof}
The measure space has coordinates
\[
 s=(A_0,z_0,\alpha,k,I,I'),
\]
with mirror variables $(A_0,z_0)$, cut-cube variables
$\alpha\in[0,2\pi)^{Q_\times}$, integers $k\in\Z^{P_\times}$, and
the two finite mode subsets. Define
\begin{align}\label{eq:reg-positive-measure}
 \dd\mathfrak m(s)={}&W_0\,\dd(A_0,z_0)
     \prod_{Q\in Q_\times}\frac{\dd\alpha_Q}{2\pi}
     \prod_{\sigma\in P_\times}
                         \frac{q}{\sqrt{2\pi}}e^{-\beta_\sigma^2/2}
 \nonumber\\
 &\quad\cdot\omega_I\omega_{I'}r_q^{|B_\times|}
        1_{\{d^*k+\mathfrak q(I\cup I')=0\text{ on }B_\times\}},
\end{align}
where $z_0,k,I,I'$ carry counting measure. Put
\begin{equation}\label{eq:reg-half-transform}
 \widehat G(s)=\int W_+u_{\alpha,k}
         \int_{\mathcal X_+}G e^{-S_+}\zeta_I\check\zeta_{I'}
                                      \,\dd(A_+,z_+).
\end{equation}
The action $S_+$ has bounded coefficients depending only on upper
and mirror links. Thus the integrand in this formula is a
half-integrand of the stated kind.

We first prove the identity for bounded coefficients and definite
parities. Inserting \eqref{eq:reg-two-crossings} into the reflected
functional combines each term with the half factors
\[
 \xi=G e^{-S_+}\zeta_I\check\zeta_{I'},\qquad
 \xi'=G'e^{-S_+}\zeta_I\check\zeta_{I'}.
\]
The sign from this rearrangement is
$(-1)^{(|G|+|G'|)(|I|+|I'|)}$. It is one when the parities agree.
When they disagree, one of $\xi,\xi'$ is odd, so its half Berezin
integral vanishes; the corresponding full integral also vanishes
by \eqref{eq:reg-half-berezin}. The sign consequently disappears
from the integrated formula in all cases.

Now insert \eqref{eq:reg-photon-crossing} and integrate the cut
links. Formula \eqref{eq:reg-gauss-indicator} supplies the last
factor in \eqref{eq:reg-positive-measure}. At fixed $s$, the
remaining factors are
\[
 \Theta\bigl(W_+u_{\alpha,k}\xi'\bigr)
                    W_+u_{\alpha,k}\xi.
\]
The lower photon variables map measure-preservingly onto an
independent set of upper variables. Combining this change of
variables with \eqref{eq:reg-half-berezin} gives
$\overline{\widehat{G'}(s)}\widehat G(s)$, proving
\eqref{eq:reg-positive-representation} once the exchanges of
integration have been justified.

For that justification let $C$ bound all the finitely many half
Berezin coefficients of $G e^{-S_+}\zeta_I\check\zeta_{I'}$ and of
the corresponding expression with $G'$. The scalar phase factors
have modulus one. Before the cut-link integration, an absolute
majorant therefore has total integral at most
\begin{equation}\label{eq:reg-representation-majorant}
 C^2(r_qC_q)^{3V}
       \prod_{i\text{ at }\Pi}(1+\omega_i)
       \prod_{i'\text{ at }\Pi'}(1+\omega_{i'}).
\end{equation}
The upper, lower and mirror variables are bounded by their
Gaussian sums; the cut-link intervals contribute their lengths;
the dual photon weights are bounded by $C_q^{|P_\times|}$; and the
finite mode sums give the two products. This proves absolute
integrability and all required uses of Fubini's theorem. The
same bounds without the half factors show that $\mathfrak m$ is
finite. It is positive by construction. Formula
\eqref{eq:reg-half-transform} is measurable and bounded uniformly
in $s$ by $Cr_q^{|B_+|}C_q^{|P_+|}$. Mixed parities follow by
sesquilinearity.

For $G,G'\in\mathcal G^+$, use
\eqref{eq:reg-upper-growth} with exponents $\varepsilon$ and
$\varepsilon'$. The algebraic identities above are unchanged.
The absolute majorant now replaces the upper and lower Gaussian
sums by $C_q(\varepsilon)^{|P_+|}$ and
$C_q(\varepsilon')^{|P_-|}$, respectively, while the mirror and
cut sums retain exponent $1/2$. These quantities are finite by
\eqref{eq:reg-gaussian-sum}. This proves coefficientwise
integrability, the same representation, and the boundedness of
$\widehat G$. Taking $G'=G$ proves positivity.

For a scalar $\Phi$ with this plaquette dependence, divide this positivity
by $Z_{\mathbb T}[1]>0$ and use the definition of the curvature
law. Reflection invariance follows from
\eqref{eq:reg-reflected-functional}, since the adapted torus is
preserved by reflection. This proves
\eqref{eq:reg-law-positivity} and completes the theorem.
\end{proof}

The same form remains positive after multiplication of the weight
by $\Phi_+\Theta(\Phi_+)$, with $\Phi_+$ bounded and scalar on the
upper half: apply the theorem to $G\Phi_+$. This observation will
be used for the reflected arrays of observables. It uses positivity
of the joint reflected form before integrating out the fermion.
The determinant's pointwise nonnegativity by itself does not imply
reflection positivity.


\section{Slab operators and finite-volume comparison}\label{sec:slab-comparison}

A cubic torus admits the coordinate reflection-positive representations
of the preceding section. A diagonal mirror requires a different period
lattice. To compare the two geometries, we express the functional as a
trace of operators associated with slabs. Reflection gives the adjoint
of a slab operator; repeated reflected copies give its Schatten norm.
The resulting comparison changes the closing vector, and, when needed,
the fermion character. Three such changes connect the cubic torus with
the diagonal torus and with a second coordinate torus of the same volume.
Their pressures differ by a quantity tending to zero. This is the
finite-volume input for diagonal reflection positivity. The slab
kernel rules and trace are proved in \cref{thm:slab-operators};
the folds and pressure estimates are proved in
\cref{prop:slab-fold-holder,thm:slab-chessboard,prop:slab-pressure}.

The same boundary construction will give the coordinate trace used in
\cref{sec:coordinate} to prove Gaussian domination. There the real
external two-cochain enters through unitary electric translations and
magnetic multiplication operators; two choices of time direction
remove its electric and magnetic parts.

All parameters remain fixed. We use the counting Berezin convention
\eqref{eq:reg-berezin}.

\subsection{Local observables and their lifts}\label{subsec:slab-local}

We need local observables that may contain fermion variables and whose
coefficients grow with the curvature. Their definition must make sense
in every flux sector. Let $B=y+[-r,r]^3$, and write $B^{(j)}$ for the
oriented lattice $j$-cells contained in $B$. The fermion sites assigned
to $B$ are
\[
 \mathcal X(B)=\{x\in\mathcal X:
          \pi(x)+[-\tfrac12,\tfrac12]^3\subset B\}.
\]
This support condition includes the photon site at which the fermion
transforms under gauge transformations.

Let $\mathcal G(B)$ consist of the Borel maps
$G:\R^{B^{(1)}}\to\operatorname{Gr}(\mathcal X(B))$ with gauge-invariant
values and with the following growth bound. If
$G(\widehat A)=\sum_\Psi a_\Psi(\widehat A)\Psi$, then, for some
$C<\infty$ and $\varepsilon>0$,
\begin{equation}\label{eq:slab-box-growth}
 |a_\Psi(\widehat A)|\le C
   \exp\left((\tfrac12-\varepsilon)
                    \sum_{\sigma\in B^{(2)}}(d\widehat A)_\sigma^2\right).
\end{equation}
Here gauge invariance means
\begin{equation}\label{eq:slab-box-gauge}
 a_\Psi(\widehat A-d\lambda)
        e^{iq\langle\lambda,\mathfrak q(\Psi)\rangle}
       =a_\Psi(\widehat A)
\end{equation}
for every real function $\lambda$ on the photon sites of $B$;
$\mathfrak q(\Psi)$ counts the $\psi$ generators minus the
$\bar\psi$ generators at fermion sites projecting to each photon site.
We denote the even subspace by $\mathcal G^{\mathrm{ev}}(B)$.

A local lift of a configuration $(A,z)$ on $B$ is a real one-cochain
$\widehat A$ on $B^{(1)}$ satisfying
\begin{equation}\label{eq:slab-local-lift}
 d\widehat A=F\quad\text{on }B^{(2)},\qquad
 e^{iq\widehat A(b)}=e^{iqA(b)}\quad(b\in B^{(1)}).
\end{equation}
We evaluate $G$ at this lift and then impose the fermion identifications
of the torus or cylinder. We now verify that this procedure is Borel
and independent of the lift.

Replace $B$ by its integer box
$\prod_\mu[a_\mu,b_\mu]$, where $a_\mu=\lceil y_\mu-r\rceil$ and
$b_\mu=\lfloor y_\mu+r\rfloor$. If this box has no sites, use the
unique empty lift. Otherwise its axial spanning tree consists of all
links in direction $3$, the links in direction $2$ on $x_3=a_3$, and the
links in direction $1$ on $x_2=a_2$, $x_3=a_3$. Prescribe a one-cochain
on this tree. Its values on the remaining links are determined by
\begin{align}\label{eq:slab-lift-recursion}
 \widehat A_1(x+e_2)&=\widehat A_1(x)+\widehat A_2(x+e_1)
                                  -\widehat A_2(x)-F_{12}(x),
                                  &&x_3=a_3,\notag\\
 \widehat A_1(x+e_3)&=\widehat A_1(x)+\widehat A_3(x+e_1)
                                  -\widehat A_3(x)-F_{13}(x),\\
 \widehat A_2(x+e_3)&=\widehat A_2(x)+\widehat A_3(x+e_2)
                                  -\widehat A_3(x)-F_{23}(x).\notag
\end{align}
First use the first line in increasing $x_2$, then the last two lines
in increasing $x_3$. Each link outside the tree occurs as an unknown
once. The construction is linear in $F$ and the tree values. It gives
the prescribed $13$- and $23$-curvatures and the $12$-curvature on the
bottom face. If $dF=0$, induction over the cubes gives the remaining
$12$-curvatures: on each cube the two closed two-cochains $F$ and
$d\widehat A$ agree on five faces, hence on the sixth.

On the tree choose $\widehat A(b)=\arg_q(e^{iqA(b)})$, where
$\arg_q$ takes values in $[0,r_q)$. This is a Borel choice. Subtract
the recursion for $A$ from \eqref{eq:slab-lift-recursion}. The tree
differences lie in $r_q\Z$, and every subsequent difference is a signed
sum of earlier ones and of a value of $r_qz$. Thus
$\widehat A-A$ is integer-valued in units of $r_q$ on every link. This
proves \eqref{eq:slab-local-lift}, with Borel dependence on $F$ and the
tree phases.

Two local lifts differ by $r_qk$, where $k$ is an integer closed
one-cochain. Sum $k$ along the tree paths from its bottom corner to
define an integer function $\lambda$. Then $d\lambda=k$ on the tree.
The zero-curvature recursion shows that $k-d\lambda=0$ on every link.
Equation \eqref{eq:slab-box-gauge}, applied with $-r_q\lambda$, gives
$a_\Psi(\widehat A+r_qd\lambda)=a_\Psi(\widehat A)$, since
$\langle\lambda,\mathfrak q(\Psi)\rangle$ is integer. The value of $G$
is therefore independent of the lift.

For an open set $U$, let $\mathcal A(U)$ be the linear span of the maps
\begin{equation}\label{eq:slab-admissible-product}
 \widehat A\longmapsto\Phi(d\widehat A)
                   \prod_{j=1}^rG_j(\widehat A|_{B_j^{(1)}}),
\end{equation}
where the boxes $B_j\subset U$ are pairwise disjoint,
$G_j\in\mathcal G^{\mathrm{ev}}(B_j)$, and $\Phi$ is a bounded Borel
function of finitely many curvature values at plaquettes contained in
$U$. A box $B$ embeds in a period group $\Gamma$ when the sets
$B+[-1,1]^3+\gamma$, $\gamma\in\Gamma$, are pairwise disjoint.
Boxes $B_1,\ldots,B_r$ embed disjointly in $\Gamma$ when each
embeds and all sets $B_j+\gamma$, indexed by distinct pairs
$(j,\gamma)$, are pairwise disjoint. A product is read on the
quotient only when its boxes embed disjointly.

It remains to check that equal maps \eqref{eq:slab-admissible-product}
have equal values on configurations. This requires one lift for every
box in a representation. On $\Z^3$, a closed integer two-cochain $z$
has the integer primitive
\begin{align*}
 k_1(x)&=0,\\
 k_2(x)&=\sum_{s=0}^{x_1-1}z_{12}(s,x_2,x_3),\\
 k_3(x)&=\sum_{s=0}^{x_1-1}z_{13}(s,x_2,x_3)
                    +\sum_{s=0}^{x_2-1}z_{23}(0,s,x_3).
\end{align*}
The sums are oriented for negative endpoints. The first two components
of $dk=z$ follow by taking differences. In the third, $dz=0$ changes
the summand into a difference of consecutive $z_{23}$ values, so the
sum telescopes. Therefore $\widehat A=A+r_qk$ is a lift on the whole
lattice. It need not be periodic, even when $(A,z)$ is periodic.

The same conclusion holds on a closed slab. Choose a coordinate
$\mu_0$ absent from the normal $n_\ell$. The slab cell complex is a
planar strip times $\Z$ in direction $\mu_0$.
Let $\mu<\nu$ be the other coordinates, denoted by $(a,r)$.
Write $Q(a,r)$ for the unit square with lower corner $(a,r)$ and
$I_r$ for the interval of $a$ such that $[(a,r),(a,r+1)]$ lies
in the strip. When $I_r$ is nonempty choose $a_r\in I_r$.
On the base sheet set the $\mu$-link values of $k$ to zero and put
\[
 k([(a,r),(a,r+1)]\times\{0\})
       =\sum_{a_r\le s<a}z(Q(s,r)\times\{0\}).
\]
All squares in this oriented sum lie in the strip because $I_r$
is an interval. Set $k=0$ on the $\mu_0$-links. For a horizontal
link $b$ in direction $\nu'\in\{\mu,\nu\}$, extend by
\[
 k(b\times\{j\})=k(b\times\{0\})
       +s_{\nu'}\sum_{0\le j'<j}z(b\times[j',j'+1]),
 \qquad
 s_{\nu'}=\begin{cases}-1,&\nu'<\mu_0,\\+1,&\nu'>\mu_0.\end{cases}
\]
The sums with reversed endpoints have the opposite sign.
The resulting primitive has $dk=z$ on all vertical
plaquettes and on the base sheet. The closed two-cochain $z-dk$ then
has vanishing vertical components; closedness makes its horizontal
component constant between sheets. It vanishes on the base sheet and
hence everywhere. Again the primitive is taken on the lifted strip,
not on its periodic quotient.

Evaluation at this common lift, followed by the quotient homomorphism
on fermion generators, agrees with the separate local evaluations.
It is linear in the map $G$, so it is independent of any representation
of $G$ as a sum of products. Moreover, different box plaquettes have
different representatives on the quotient. Multiplying
\eqref{eq:slab-box-growth} by the Maxwell density therefore leaves
\begin{equation}\label{eq:slab-local-majorant}
 C\exp(-\varepsilon_G\|F\|^2),\qquad
 \varepsilon_G=\min(\tfrac12,\varepsilon_1,\ldots,\varepsilon_r)>0.
\end{equation}
The finite Grassmann expansion of the action has bounded coefficients
for real $A$. The Gaussian sum \eqref{eq:reg-gaussian-sum} now proves
absolute integrability, coefficient by coefficient. This also proves
the required Borel measurability on the quotient.

\subsection{Boundary spaces and slab kernels}\label{subsec:slab-kernels}

Fix a mirror family $\ell$ and a cross-section
$\Xi=(\Gamma_\parallel,\chi_\parallel)$, where
$\Gamma_\parallel\subset\Z^3\cap\{\ell=0\}$ has rank two and systole
at least four. Here and below, the systole of a period lattice
$\Gamma$ is $\min_{\gamma\in\Gamma\setminus\{0\}}|\gamma|_\infty$.
Work on the cylinder $\Z^3/\Gamma_\parallel$. The
admissible mirror positions are
$c\in\frac12|n_\ell|^2\Z$. Write $\mathcal M^j(c)$ for the $j$-cells
contained in $\{\ell=c\}$, $\mathcal X^j(c)$ for those cut through
their relative interiors by this plane, and $\mathcal B^j(c,c')$ for
the $j$-cells contained in $\{c\le\ell\le c'\}$ and in neither
bounding plane. In particular a slab cell may touch either boundary.
All these sets are finite on a bounded slab.

For a chain of admissible mirrors, the cylinder cells partition into
the cells in its mirrors, those cut by them, and those lying between
consecutive mirrors. A cell cut by $c$ has its height range
$[c-|n_\ell|^2/2,c+|n_\ell|^2/2]$. Its non-cut faces lie on one
side of the plane; the other faces are cut by the same plane. This
partition applies also when two mirrors are consecutive. A fermion
bond crosses at most one mirror of such a chain, and no fermion site
lies in a mirror.

A closed-slab configuration $\varpi=(A,z)$ assigns real link values
and integer plaquette values to all cells of the closed slab, with
$dz=0$ on its cubes. We integrate its interior values using
$A(b)\in[0,r_q)$ and counting measure for $z$. Its interior photon
weight is
\begin{equation}\label{eq:slab-interior-weight}
 W_{c,c'}(\varpi)=
    \prod_{\sigma\in\mathcal B^2(c,c')}e^{-F_\sigma^2/2}
    \prod_{Q\in\mathcal B^3(c,c')}\mathbf1[(dz)(Q)=0].
\end{equation}
We use the larger class $\mathcal I_\Xi[c,c']$ of Borel
Grassmann-valued functions of closed-slab configurations, with fermion
support $c<\ell<c'$, invariant under
$(A,z)\mapsto(A+r_qk,z-dk)$, whose coefficients obey
\begin{equation}\label{eq:slab-integrand-growth}
 |a_\Psi(\varpi)|\le C_G
   \exp\left((\tfrac12-\varepsilon_G)
          \sum_{\sigma\in\mathcal B^2(c,c')}F_\sigma^2\right),
 \qquad 0<\varepsilon_G\le\tfrac12.
\end{equation}
The local products supported in the open slab give elements of this
class by the lift construction above. Products in consecutive slabs
remain in the class, with the smaller of the two positive
$\varepsilon$'s. Reflection and translation also preserve the class.

At the mirror $c$, let $\mathcal J(c)$ be the ordered set of crossing
modes from \eqref{eq:reg-crossing-modes}, taken modulo
$\Gamma_\parallel$. Let $\mathcal F(c)$ be the exterior space with
orthonormal basis $e_I$, $I\subset\mathcal J(c)$. Put
$\omega_c e_I=\omega_Ie_I$ and $U_{\mathrm f}e_I=(-1)^{|I|}e_I$.
Translation of the modes includes both the character
$\chi_\parallel$ and the sign required to reorder their product:
\[
 U_ve_I=s_v(I)e_{r_v(I)},\qquad
 \tau_v\zeta_I^{(c)}=s_v(I)\zeta_{r_v(I)}^{(c+\ell(v))}.
\]
The $U_v$ are unitary, satisfy the translation law, and commute with
parity. For an integer translation parallel to the normal, choose the
orders compatibly so that $s_v(I)=1$.

For a Grassmann element $H$ supported in the open slab define
\begin{equation}\label{eq:slab-fermion-matrix}
 N_{c,c'}[H]_{I',I}
    =\int_{(c,c')}\zeta_I^{(c)}H
                        \Theta_{c'}(\zeta_{I'}^{(c')}).
\end{equation}
An operator $\Upsilon$ on $\mathcal F(c)$ is represented at the
mirror by
\[
 \operatorname{Ins}_c[\Upsilon]
       =\sum_{I,J}\Upsilon_{IJ}\Theta_c(\zeta_J^{(c)})\zeta_I^{(c)}.
\]
The particular insertion
$\operatorname{Ins}_c[\omega_c\diag(u_I(A))]$ is the crossing
exponential \eqref{eq:reg-crossing-exponential}. The ordering in these
definitions gives
\begin{equation}\label{eq:slab-matrix-composition}
 N_{c,c''}[H\operatorname{Ins}_{c'}[\Upsilon]H']
             =N_{c',c''}[H']\Upsilon N_{c,c'}[H].
\end{equation}
Indeed, an entry of the left side is a sum of products
$[\zeta_IH\Theta_{c'}\zeta_J]
 [\zeta_KH'\Theta_{c''}\zeta_{I''}]$ in disjoint sets of slab
generators. Berezin factorization gives the entries on the right
without moving any of these factors. Reflecting the integrand reverses
their order and conjugates the Berezin coefficient; with
$v=2(c'-c)n_\ell/|n_\ell|^2$ this gives
\begin{equation}\label{eq:slab-matrix-adjoint}
 N_{c',2c'-c}[\Theta_{c'}H]=U_vN_{c,c'}[H]^*.
\end{equation}

We next include the photon crossing in the boundary data. The cut
incidence maps $d_c$ and $d_c^+$ are the restrictions of the coboundary
to cut faces and to upper faces, as in
\eqref{eq:reg-cut-differentials}. Introduce one angle $\vartheta_Q$
for each cut cube and one integer $k_\sigma$ for each cut plaquette,
and set
\begin{align}\label{eq:slab-dual-measure}
 \varkappa_\sigma&=r_q^{-1}
       ((d_c^*\vartheta)_\sigma-2\pi k_\sigma),\notag\\
 \dd\mu_c^\times&=
      \prod_{Q\in\mathcal X^3(c)}\frac{\dd\vartheta_Q}{2\pi}
      \prod_{\sigma\in\mathcal X^2(c)}
                     \frac{q}{\sqrt{2\pi}}e^{-\varkappa_\sigma^2/2},\\
 e_c((\vartheta,k);\varpi)&=
       e^{i\langle\vartheta,d_c^+z\rangle
                             -i\langle\varkappa,d_c^+A\rangle},\notag\\
 D_c(k,I)&=r_q^{|\mathcal X^1(c)|}
       \mathbf1[d_c^*k+\mathfrak q(I)=0].\notag
\end{align}
Here $\mathfrak q(I)(b)$ is the number of modes of groups $1,2$ over
the cut link $b$, minus the number of groups $3,4$ over it. There are
no cut links at diagonal mirrors. At a coordinate site plane there
are no cut cells at all; its Fourier space is a point.

The one-mirror photon identity proved by the Fourier and Poisson
calculation of the preceding section is
\begin{multline}\label{eq:slab-photon-gluing}
 \sum_{z^\times}\int\dd A^\times\,
       \prod_{\sigma\in\mathcal X^2(c)}e^{-F_\sigma^2/2}
       \prod_{Q\in\mathcal X^3(c)}\mathbf1[(dz)(Q)=0]\,u_I(A^\times)
       \\
 =\int\dd\mu_c^\times\,D_c(k,I)
           e_c((\vartheta,k);\varpi_+)
            \overline{e_c((\vartheta,k);\theta_c^*\varpi_-)}.
\end{multline}
The upper and lower data on the right have the same values on the
mirror. Their upper and reflected lower faces provide the non-cut
faces on the left. Cube Fourier multipliers enforce closedness,
Poisson summation gives the factors $e^{-\varkappa^2/2}$, and compact
cut-link integration gives $D_c(k,I)$. This identity is used in the
direction shown when constructing a kernel and in the opposite
direction when gluing two kernels.

Mirror variables are
$\xi=(\mathfrak m,\vartheta,k,I)$, with
$\mathfrak m=(A_0,z_0)\in[0,r_q)^{\mathcal M^1(c)}
\times\Z^{\mathcal M^2(c)}$. On their space $Y_c$ put
\begin{equation}\label{eq:slab-boundary-measure}
 \dd\mu_c(\xi)=\dd A_0
      \prod_{\sigma\in\mathcal M^2(c)}e^{-F_\sigma(\mathfrak m)^2/2}
                  \dd\mu_c^\times(\vartheta,k)D_c(k,I),
 \qquad \mathcal H(c)=L^2(Y_c,\mu_c).
\end{equation}
The integer and mode variables carry counting measure. This is a
finite positive measure. Summing the mirror plaquette Gaussians gives
$C_q^{|\mathcal M^2(c)|}$, while Poisson summation at fixed angles
gives
\[
 \frac q{\sqrt{2\pi}}\sum_{k_\sigma}e^{-\varkappa_\sigma^2/2}
 \le\sum_{j\in\Z}e^{-r_q^2j^2/2}.
\]
The remaining factors are bounded by
$r_q^{|\mathcal M^1(c)|+|\mathcal X^1(c)|}2^{|\mathcal J(c)|}$.
The induced signed translations $V_v:\mathcal H(c)\to
\mathcal H(c+\ell(v))$ are unitary and satisfy the group law. The
parity operator multiplies the component $I$ by $(-1)^{|I|}$.

Let $S_{(c,c')}(A)$ be the fermion action containing precisely the mass
and bond terms with fermion endpoints in the open slab. For
$G\in\mathcal I_\Xi[c,c']$, define the kernel
\begin{multline}\label{eq:slab-kernel}
 \mathfrak S[G](\xi',\xi)=\sqrt{\omega_{I'}\omega_I}
      \int\dd\varpi^\circ\,W_{c,c'}(\varpi)
      e_c((\vartheta,k);\varpi)
       \overline{e_{c'}((\vartheta',k');\theta_{c'}^*\varpi)}
       \\
       \cdot N_{c,c'}[G(\varpi)e^{-S_{(c,c')}(A)}]_{I',I},
       \qquad \varpi=\mathfrak m\vee\varpi^\circ\vee\mathfrak m'.
\end{multline}
Set the whole integrand to zero where $W_{c,c'}(\varpi)=0$;
$G(\varpi)$ is evaluated only on the closed configurations in its
domain.
The associated integral operator maps $\mathcal H(c)$ to
$\mathcal H(c')$. The square roots of the mode weights belong in the
kernel; the measure \eqref{eq:slab-boundary-measure} contains none of
these weights. This places one full crossing weight at a glued mirror.

\begin{proposition}\label{thm:slab-operators}
Every $\mathfrak S[G]$ is Hilbert--Schmidt and depends linearly on $G$.
If $G$ has constants $C_G,\varepsilon_G$ in
\eqref{eq:slab-integrand-growth}, then
\begin{align}\label{eq:slab-kernel-bound}
 |\mathfrak S[G](\xi',\xi)|
    &\le C_{c,c'}C_GC_q(\varepsilon_G)^{|\mathcal B^2(c,c')|},\notag\\
 \|\mathfrak S[G]\|_2
    &\le\sqrt{\mu_c(Y_c)\mu_{c'}(Y_{c'})}
                C_{c,c'}C_GC_q(\varepsilon_G)^{|\mathcal B^2(c,c')|}.
\end{align}
The constant depends on the fixed cross-section, mirror positions,
$q,M,\rho,w_i$, and chamber base point. For consecutive slabs,
\begin{align}\label{eq:slab-operator-rules}
 \mathfrak S[G']\mathfrak S[G]&=\mathfrak S[GG'],\notag\\
 \mathfrak S[\tau_vG]&=V_v\mathfrak S[G]V_v^{-1},\\
 \mathfrak S[\Theta_{c'}G]&=V_{2(c'-c)n_\ell/|n_\ell|^2}
                                         \mathfrak S[G]^*.\notag
\end{align}
If $\mathbb T=(\Gamma,\chi)$ has cross-section $\Xi$, $g\in\Gamma$
is a closing vector with minimal positive height $p=\ell(g)$, and
$c_0<\cdots<c_r=c_0+p$ are admissible, $r\ge2$, and
$G_j\in\mathcal I_\Xi[c_{j-1},c_j]$ for every $j$, then
\begin{equation}\label{eq:slab-trace}
 Z_{\mathbb T}[G_1\cdots G_r]
    =\operatorname{Tr}_{\mathcal H(c_0)}
       \left(V_g^{-1}U_{\mathrm f}^{(1+\chi(g))/2}
               \mathfrak S[G_r]\cdots\mathfrak S[G_1]\right).
\end{equation}
These identities apply to arbitrary parities. For homogeneous $G$,
$U_{\mathrm f}\mathfrak S[G]=(-1)^{\deg G}\mathfrak S[G]U_{\mathrm f}$.
\end{proposition}

\begin{proof}
Each coefficient of the finite Grassmann polynomial
$e^{-S_{(c,c')}(A)}$ is bounded uniformly in real $A$. Each coefficient
in the matrix in \eqref{eq:slab-kernel} is consequently bounded by a
fixed multiple of the majorant for $G$. The phases have modulus one.
After multiplication by the photon weight, the absolute integrand is
bounded by
\[
 C C_G\prod_{\sigma\in\mathcal B^2(c,c')}
                           e^{-\varepsilon_GF_\sigma^2}.
\]
Drop the cube constraints and sum the integers separately. Each sum
is at most $C_q(\varepsilon_G)$; each compact link integral contributes
$r_q$. The mode factors and finite matrix dimension enter $C_{c,c'}$.
This proves the pointwise bound in \eqref{eq:slab-kernel-bound}.
The boundary measures are finite, so the kernel is in
$L^2(Y_{c'}\times Y_c)$, with the stated norm. These bounds also
justify all the gluing integrals below before any cancellation.

For composition, multiply the two kernels and integrate over the common
boundary. Their two square roots give $\omega_{I'}$ for the middle
mode. The two phases and its boundary measure are precisely the right
side of \eqref{eq:slab-photon-gluing}. Applying that identity restores
the photon variables on the cut cells, their Maxwell weights, their
cube constraints, and the phase $u_{I'}$ of the crossing modes. Sum
the middle modes by \eqref{eq:slab-matrix-composition}, with
$\Upsilon=\omega_{c'}\diag(u_{I'})$. This restores the crossing
fermion action. The plaquettes in the middle mirror come from its
boundary measure, and the cell partition restores the full weight on
the combined slab. No cube lies in a mirror, so every closedness
constraint is present once. This is the kernel of $\mathfrak S[GG']$.

Translation acts by a signed relabeling of the modes and an
orientation-preserving relabeling of cells, proving covariance.
For reflection use \eqref{eq:slab-matrix-adjoint}; the two boundary
phases exchange their roles and the interior weight is preserved.
Reflection may reverse a compact integration interval, but integer
shift invariance makes any interval of length $r_q$ equivalent to
$[0,r_q)$. The reflected upper boundary is the translate of the
original lower one by $2(c'-c)n_\ell/|n_\ell|^2$. This gives the last
identity of \eqref{eq:slab-operator-rules}. Parity follows because a
Berezin integral of odd total degree vanishes.

For the trace, composition reduces the claim to two slabs. A product
of two Hilbert--Schmidt kernels is trace class, and its trace is the
absolutely convergent integral of their cyclic product. The closing
operator identifies the upper mirror data with the translate by $g$
of the lower data. Its contribution to this cyclic kernel is
\[
 (-\chi(g))^{|I_0|}s_g(I_0).
\]
Here $s_g$ is the mode relabeling sign; $(-\chi(g))^{|I_0|}$ is the
remaining Grassmann closing sign. It is $1$ for antiperiodic closure
and the fermion parity for periodic closure. To see this sign directly,
translate the lower crossing product to the upper boundary. Each of
its $|I_0|$ generators contributes $\chi(g)$, while moving it across
the adjacent upper product contributes $(-1)^{|I_0|}$ in every term
of even total degree. Their product is $(-\chi(g))^{|I_0|}$.

Apply \eqref{eq:slab-photon-gluing} at the two mirrors. For fixed
photon variables the remaining mode sum is
\[
 \tr_{\mathcal F(c_0)}
    (U_g^{-1}U_{\mathrm f}^{(1+\chi(g))/2}
                       N_2\Upsilon_1N_1\Upsilon_0),
 \qquad \Upsilon_j=\omega_{c_j}\diag(u_I(A)).
\]
Berezin composition and closure replace this trace by the full torus
Berezin integral. The two open-slab actions and the two crossing
actions include each mass and hopping term once. The photon variables
likewise include every torus link and plaquette once, and the cut and
interior cube constraints impose $dz=0$ everywhere. Their measure is
therefore the measure in \eqref{eq:reg-functional}. This proves
\eqref{eq:slab-trace} for even total degree. With odd total degree both
sides vanish, the operator under the trace anticommuting with parity.
Finally, expand nonhomogeneous entries into their even and odd parts.
\end{proof}

\subsection{Reflected folds and Schatten estimates}\label{subsec:slab-folds}

For an operator $Y$ between Hilbert spaces, write
$\|Y\|_s=(\operatorname{Tr}(Y^*Y)^{s/2})^{1/s}$. We use Schatten's
inequality in the following form; see \cite[slides~55--57]{PongeSchatten2025}. If $1\le s_j<\infty$,
$Y_j:H_{j-1}\to H_j$ belongs to the Schatten class of exponent $s_j$,
$H_r=H_0$, and $\sum_js_j^{-1}=1$, then the cyclic product is trace
class and
\begin{equation}\label{eq:slab-schatten-holder}
 |\operatorname{Tr}(Y_r\cdots Y_1)|
                \le\prod_{j=1}^r\|Y_j\|_{s_j}.
\end{equation}
The statement for different spaces follows from the usual inequality
on their orthogonal direct sum: extend each $Y_j$ by zero on all other
summands. The nonzero block of the product is the displayed cyclic
product, and the singular values of each extension are those of $Y_j$.
A unitary closing identification can be absorbed into $Y_r$ without
changing its Schatten norm.

Let $\mathbb T$ have height period $p$ in the family $\ell$. Its
straightened torus $\mathbb T^0$ has the same cross-section and the
closing vector
\begin{equation}\label{eq:slab-straight-vector}
 g^0=\frac p{|n_\ell|^2}n_\ell,
 \qquad \Gamma^0=\Gamma_\parallel\oplus\Z g^0,
 \qquad\chi^0(g^0)=-1.
\end{equation}
We use this torus only when $g^0\in\Z^3$ and its lattice has systole
at least four. Its volume is the volume of $\mathbb T$: both lattices
have the same covolume in the mirror and the same normal height period.

Divide one period into $2k$ slabs of equal height $h'$, with all
$c_j=c_0+jh'$ admissible and $p=2kh'$. For an integrand $G$ in one of
these slabs, form its fold by reflecting it successively in the
bounding planes, taking all $2k$ copies. This is an integrand on
$\mathbb T^0$. Its fermion boxes are disjoint: they lie strictly
between different consecutive mirrors; within each slab the
cross-section embedding separates their translates.

\begin{proposition}\label{prop:slab-fold-holder}
For integrands $G_j\in\mathcal I_\Xi[c_{j-1},c_j]$,
\begin{align}\label{eq:slab-fold-holder}
 Z_{\mathbb T^0}[\operatorname{fold}_\ell G_j]
              &=\|\mathfrak S[G_j]\|_{2k}^{2k},\notag\\
 |Z_{\mathbb T}[G_1\cdots G_{2k}]|
              &\le\prod_{j=1}^{2k}
               Z_{\mathbb T^0}[\operatorname{fold}_\ell G_j]^{1/(2k)}.
\end{align}
In particular every folded factor on the right is nonnegative.
\end{proposition}

\begin{proof}
Fix $j$ and write $a=c_{j-1}$. Put
$Y=\mathfrak S[G_j]:\mathcal H(a)\to\mathcal H(a+h')$, and form
the fold on $[a,a+p]$, with initial boundary $a$ in the trace formula.
Let
$v=2h'n_\ell/|n_\ell|^2$, which is integer because both mirror
positions are admissible. The first reflected pair has product
$V_vY^*Y$ by \eqref{eq:slab-operator-rules}. Its translate by $sv$
has product $V_{(s+1)v}Y^*YV_{sv}^{-1}$. Multiplying all $k$ pairs
telescopes to
\[
 V_{kv}(Y^*Y)^k.
\]
Since $kv=g^0$ and $\chi^0(g^0)=-1$, the closing translation in
\eqref{eq:slab-trace} cancels $V_{g^0}$ and introduces no parity.
The trace is
$\operatorname{Tr}_{\mathcal H(a)}(Y^*Y)^k=\|Y\|_{2k}^{2k}$.
This calculation takes place at the starting boundary of each slab
and proves the first identity for every $G_j$.

For the original torus, the operator
$V_g^{-1}U_{\mathrm f}^{(1+\chi(g))/2}$ is unitary. Apply
\eqref{eq:slab-schatten-holder} to its trace formula, with every
exponent equal to $2k$. The first identity identifies the resulting
norms with folded partition functions. The Hilbert--Schmidt estimate
ensures all operators belong to the required Schatten classes.
\end{proof}

We next apply this inequality to cells. This is the repeated-reflection
chessboard argument of \cite[Section~4]{FILS1978}, with the period
lattice and its fermion character carried through each fold. A frame is a triple of the
mirror families with mutually orthogonal normals $n_1,n_2,n_3$.
The two frames needed below are $(x_1,x_2,x_3)$ and
$(x_1-x_2,x_1+x_2,x_3)$. Choose admissible offsets $o_a$ and an integer
height width $h>0$. The open cells are
\begin{equation}\label{eq:slab-cells}
 Q_\kappa=\{x:o_a+h\kappa_a<\ell^a(x)
                         <o_a+h(\kappa_a+1),\ a=1,2,3\},
 \qquad |Q|=\frac{h^3}{|n_1||n_2||n_3|}.
\end{equation}
The reflections in the planes $\ell^a=o_a+hj$ generate a product of
three infinite dihedral groups. Exactly one element of this group
carries a given open cell to any specified open cell. Its action on
even Grassmann-valued functions is generated by the reflections
$\Theta$ and translations used above. Denote the resulting copy of
$G\in\mathcal A(Q_\kappa)$ in $Q_{\kappa'}$ by $G^{[\kappa']}$.

A rectangular torus in this frame has periods
$\gamma_a=p_an_a/|n_a|^2\in\Z^3$, with $\chi(\gamma_a)=-1$.
If each $p_a\in2h\Z$, the copies fit its period lattice. Let
$N=|\mathbb T|/|Q|$ be the number of cells, and define
\begin{equation}\label{eq:slab-chessboard-norm}
 \operatorname{Fold}_{\mathbb T}G
          =\prod_{\kappa'\bmod\Gamma}G^{[\kappa']},
 \qquad \|G\|_{\mathbb T}
          =Z_{\mathbb T}[\operatorname{Fold}_{\mathbb T}G]^{1/N}.
\end{equation}
This is an unnormalized seminorm. Its defining partition function is
nonnegative, by the first identity of
\eqref{eq:slab-fold-holder} applied to one layer of the folded array.
The value is independent of the chosen representatives of the cells.

\begin{proposition}\label{thm:slab-chessboard}
If a rectangular torus $\mathbb T$ carries the cells
\eqref{eq:slab-cells}, then
\begin{equation}\label{eq:slab-chessboard-rectangular}
 \left|Z_{\mathbb T}\left[\prod_{\kappa}G_\kappa\right]\right|
                       \le\prod_\kappa\|G_\kappa\|_{\mathbb T},
 \qquad G_\kappa\in\mathcal A(Q_\kappa).
\end{equation}
Let $C(a,b,c)$ have coordinate periods $ae_1,be_2,ce_3$,
all with character $-1$.
Let $e_-=e_1-e_2$, $e_+=e_1+e_2$, and let $D(a,b,c)$ have periods
$ae_-,be_+,ce_3$, all with character $-1$. Let $D'(a,b,c)$ have the
same periods and characters, except for $\chi(be_+)=+1$.
For the diagonal frame and $n\in2h\Z$, $n\ge8$, the cells have
representatives
\[
 0\le\kappa_1<n/h,\qquad 0\le\kappa_2<2n/h,
                     \qquad0\le\kappa_3<n/h.
\]
They represent the cells of $C(n,n,n)$, $D'(n/2,n,n)$, and
$D(n/2,n,n)$, and
\begin{equation}\label{eq:slab-chessboard-diagonal}
 \left|Z_{\mathbb T}\left[\prod_\kappa G_\kappa\right]\right|
           \le\prod_\kappa\|G_\kappa\|_{D(n/2,n,n)},
 \quad \mathbb T=C(n,n,n)\ \text{or }D'(n/2,n,n).
\end{equation}
For coordinate cells of width $h$, $a,c\in2h\Z$, $a,c\ge4$, one has
\begin{equation}\label{eq:slab-chessboard-coordinate}
 \left|Z_{D(a,2a,c)}\left[\prod_\kappa G_\kappa\right]\right|
                \le\prod_\kappa\|G_\kappa\|_{C(a,4a,c)}.\qedhere
\end{equation}
\end{proposition}

\begin{proof}
Group the cell integrands into layers perpendicular to the first normal,
and apply \eqref{eq:slab-fold-holder}. Its right side consists of the
arrays obtained by retaining one layer and reflecting it throughout
the straightened torus. Next apply the same inequality to each such
array in direction two, and then in direction three. The reflected
array after these three steps is the full fold of the integrand from
one original cell. The three exponents multiply to
\[
 \frac h{p_1}\frac h{p_2}\frac h{p_3}
       =\frac1N.
\]
Thus the last product consists of one factor
$\|G_\kappa\|_{\mathbb T_3}$ for each original cell, where
$\mathbb T_3$ is the final rectangular torus.

We verify that all three arrays have consistent periods and that the
same cells represent them. If $\Gamma_0$ is the initial lattice, then
straightening successively in the frame gives
\begin{align*}
 \Gamma_1&=(\Gamma_0\cap\{\ell^1=0\})\oplus\Z\gamma_1,\\
 \Gamma_2&=(\Gamma_0\cap\{\ell^1=\ell^2=0\})
                                  \oplus\Z\gamma_1\oplus\Z\gamma_2,\\
 \Gamma_3&=\Z\gamma_1\oplus\Z\gamma_2\oplus\Z\gamma_3,
 \qquad \gamma_a=p_an_a/|n_a|^2.
\end{align*}
The kernel in the second line is a rank-one group along $n_3$; its
positive generator is $\gamma_3$, by the definition of the third
height period. Characters already fixed at $\gamma_1,\gamma_2$ are
retained in subsequent cross-sections, so all three final characters
are $-1$. The brick $0\le\kappa_a<p_a/h$ represents $\Gamma_3$.
It represents $\Gamma_2$ as well; then its section $\kappa_2=0$
represents the second cross-section, and its section $\kappa_1=0$
represents the first one. Restoring the normal coordinate in each
case gives representatives for the preceding lattice. These groups
all have the same covolume.

We also check the stronger box embedding condition used in the local
classes. If $B=y+[-r,r]^3$ lies in an open cell, then
$2|n_a|^2r<h$ for every $a$, by taking the range of $\ell^a$ on $B$.
For a nonzero period $\gamma$, at least one height is nonzero and
$|\ell^a(\gamma)|\ge2h$. Consequently
$|\gamma|_\infty>4r$, and the systole also gives
$|\gamma|_\infty\ge4$. These inequalities imply
$|\gamma|_\infty>2r+2$: use the second if $r<1$, the first if
$r\ge1$. Thus translates of $B+[-1,1]^3$ are disjoint. Boxes in
different cell orbits cannot meet after period translations, since
the open cells themselves are disjoint. It follows that every layer
belongs to its slab class. This justifies the three applications.

For a rectangular torus, every straightening leaves the torus unchanged,
which proves \eqref{eq:slab-chessboard-rectangular}. For a cubic torus
with even side $n$, the diagonal operations are
\begin{equation}\label{eq:slab-three-untwistings}
 C(n,n,n)\xrightarrow{\ell_-}D'(n/2,n,n)
          \xrightarrow{\ell_+}D(n/2,n,n),
 \qquad D(n/2,n,n)\xrightarrow{x_1}C(n/2,2n,n).
\end{equation}
For the first operation the cross-section is generated by $ne_+,ne_3$,
with characters $+1,-1$; its height period is $n$, and straightening
replaces the closing vector $ne_1$ by $(n/2)e_-$. For the second, the
cross-section is generated by $(n/2)e_-,ne_3$, both with character
$-1$. Its closing vector $ne_+$ already points along the normal,
and the operation changes its character from $+1$ to $-1$. For the
third, the cross-section is generated by $2ne_2,ne_3$, both with
character $-1$; the height period is $n/2$ and the straight closing
vector is $(n/2)e_1$. The lattices have systole at least four.

The first two diagonal operations, followed by the unchanged $x_3$
operation, have height periods $(n,2n,n)$. Hence they give
\eqref{eq:slab-chessboard-diagonal}; starting at $D'$ instead, use
the order $(\ell_+,x_3,\ell_-)$ to obtain the same final fold.
Likewise, starting at $D(a,2a,c)$ and using coordinate directions
gives periods $(a,4a,c)$ and ends at $C(a,4a,c)$, proving
\eqref{eq:slab-chessboard-coordinate}.
\end{proof}

\subsection{Pressure and the comparison error}\label{subsec:slab-pressure}

Write $C(n)=C(n,n,n)$ and
\[
 f(n_1,n_2,n_3)=\frac1{n_1n_2n_3}\log Z_{C(n_1,n_2,n_3)}[1].
\]
The mass and charge in this function are held fixed. The following
error will compare local expectations in the cubic and diagonal
geometries.

\begin{proposition}\label{prop:slab-pressure}
The pressure $f$ is nonincreasing in each side, and
\[
 f(n_1,n_2,n_3)\longrightarrow f_\infty:=\inf f\in\R
                     \quad\text{as }\min_\mu n_\mu\longrightarrow\infty.
\]
For every even $n\ge8$,
\begin{equation}\label{eq:slab-partition-chain}
 Z_{C(n)}[1]\le Z_{D'(n/2,n,n)}[1]
       \le Z_{D(n/2,n,n)}[1]\le Z_{C(n/2,2n,n)}[1].
\end{equation}
With $\Delta(n)=f(n/2,2n,n)-f(n,n,n)$,
\begin{equation}\label{eq:slab-pressure-error}
 0\le\Delta(n)\le f(n/2,n/2,n/2)-f_\infty\longrightarrow0.
\end{equation}
The logarithms of any two partition functions in
\eqref{eq:slab-partition-chain} differ in absolute value by at most
$n^3\Delta(n)$.
\end{proposition}

\begin{proof}
Fix the cross-section with sides $n_2,n_3$ and let
$Y=\mathfrak S[1]:\mathcal H(0)\to\mathcal H(1/2)$ be the coordinate
half-step. The first identity of \eqref{eq:slab-fold-holder}, applied
to $2n_1$ half-slabs, gives
\begin{equation}\label{eq:slab-pressure-trace}
 Z_{C(n_1,n_2,n_3)}[1]=\operatorname{Tr}(Y^*Y)^{n_1}.
\end{equation}
The positive trace-class operator $Y^*Y$ is independent of $n_1$ and
is nonzero, since the partition function is positive. Consequently
\[
 f(n_1,n_2,n_3)=\frac1{n_2n_3}\log\|Y^*Y\|_{n_1}.
\]
Schatten norms decrease as their exponent increases, so $f$ decreases
in $n_1$. Permuting the time direction proves the other two claims.
The same trace identity shows separate convexity of $\log Z$ in the
side lengths: H\"older's inequality gives
$\operatorname{Tr}A^{(1-t)s+ts'}\le
 (\operatorname{Tr}A^s)^{1-t}(\operatorname{Tr}A^{s'})^t$ for $A=Y^*Y$.

The lower partition bound \eqref{eq:reg-partition-bounds} gives the
volume-independent floor
\[
 f\ge\log C_*+3\log r_q-6r_q^2.
\]
Together with the upper bound in the same equation, it makes
$f_\infty$ finite. Monotonicity now gives convergence in any manner
with all sides tending to infinity: choose a finite triple at which
$f<f_\infty+\delta$ and then exceed its three sides.

Apply \eqref{eq:slab-fold-holder} with constant integrands to each of
the three operations \eqref{eq:slab-three-untwistings}. A fold of
the constant is the constant, so $Z_{\mathbb T}[1]\le Z_{\mathbb T^0}[1]$.
This proves the chain. Each torus in it has volume $n^3$, so its first
and last logarithms differ by $n^3\Delta(n)$. Side monotonicity gives
$f(n/2,2n,n)\le f(n/2,n/2,n/2)$, while
$f(n,n,n)\ge f_\infty$, proving \eqref{eq:slab-pressure-error}.
All intermediate logarithms lie between the endpoints, which proves
the last assertion.
\end{proof}

Thus the estimate concerns the logarithm per unit volume. It does not
give a bounded ratio of partition functions: the retained bound is
$e^{n^3\Delta(n)}$, with $\Delta(n)\to0$. The local comparison
divides the logarithmic error bound $n^3\Delta(n)$ by the number of
copies of the observable and by the perturbation strength.

\section{Gaussian domination on coordinate tori}\label{sec:coordinate}

The compactness estimate comes from the coordinate trace. We first
retain the real slice connections and all magnetic fluxes in that
trace, then remove the external curvature field by two applications
of Schatten H\"older in different time directions. The all-flux trace
\eqref{eq:slab-external-trace} gives the Gaussian bound in
\cref{thm:slab-gaussian}.

\subsection{Coordinate slices and the all-flux trace}\label{subsec:slab-coordinate-trace}

For $s>0$ write
\[
 g_s(x)=(2\pi s)^{-1/2}e^{-x^2/(2s)}.
\]
Tensor powers of $g_s$ denote product densities. Thus $e^{-a|E|^2}$
has convolution kernel $g_{2a}$ in each real link variable.

For Gaussian domination we use a slice representation with real
spatial connections. Take direction $1$ as time on
$C_{\mathbf n}=C(n_1,n_2,n_3)$, $\mathbf n=(n_1,n_2,n_3)$,
$n_\mu\ge4$, and let
$\Sigma=\Z/n_2\Z\times\Z/n_3\Z$. Its links and plaquettes are
$\Sigma^{(1)}$ and $\Sigma^{(2)}$, its coboundary is $d_s$, and
$v=|\Sigma|=n_2n_3$. A connection has spatial slices $A_j$ and
temporal links $a_j$, and its curvature has the parts
\begin{equation}\label{eq:slab-slice-curvatures}
 F_j^{\mathrm{el}}=A_{j+1}-A_j-d_sa_j+r_qz_j^{\mathrm{el}},
 \qquad F_j^{\mathrm{mag}}=d_sA_j+r_qz_j^{\mathrm{mag}}.
\end{equation}
Closedness reads
$z_{j+1}^{\mathrm{mag}}-z_j^{\mathrm{mag}}=d_sz_j^{\mathrm{el}}$.

We first unfold the integer variables. Put
$k_0=0$, $k_{j+1}=k_j+z_j^{\mathrm{el}}$. Closedness is equivalent to
\[
 z_j^{\mathrm{mag}}=z_0^{\mathrm{mag}}+d_sk_j,
                       \qquad d_sk_{n_1}=0.
\]
The integer shift with spatial values $k_j$ and zero temporal values
removes all electric integers except the last one, $k_{n_1}$, and
changes every spatial slice to $A_j+r_qk_j$. Summing $k_j$ while
integrating $A_j$ over its compact box unfolds the intermediate
slices to all of $\R^{\Sigma^{(1)}}$. The last integer remains as a
winding vector
\[
 g=r_qk_{n_1}\in\Gamma_s,
 \qquad\Gamma_s=r_q(\Z^{\Sigma^{(1)}}\cap\ker d_s).
\]

Fix a plaquette $\sigma_0$ of $\Sigma$ and write
$k=\sum_\sigma z_0^{\mathrm{mag}}(\sigma)$. Then
$z_0^{\mathrm{mag}}-k\mathbf1_{\sigma_0}=d_sk'$ for an integer
one-cochain $k'$. Integer exactness follows from the dual graph whose
vertices are the plaquettes. A link of
$\Sigma$ is a dual edge directed toward the plaquette where its
coefficient in $d_s$ is positive. The dual graph is connected. Choose a spanning tree rooted at
$\sigma_0$. Unit signed flow along the path from the
root to a vertex $\sigma$ has divergence
$\delta_\sigma-\delta_{\sigma_0}$. Summing these paths with the
integer coefficients of a zero-sum two-cochain gives precisely that
cochain as its divergence, hence as $d_sk'$. Two choices of $k'$
differ by an integer element of $\ker d_s$.

Shift every slice by $r_qk'$. The magnetic integers now equal
$k\mathbf1_{\sigma_0}$ in every slice. Summation over $k'$ modulo
closed integer cochains replaces the compact box for $A_0$ by a
fundamental domain of $\Gamma_s$. We arrive at paths
\[
 A_0\in\R^{\Sigma^{(1)}}/\Gamma_s,\qquad
 A_j\in\R^{\Sigma^{(1)}}\ (0<j<n_1),\qquad
 A_{n_1}=A_0+g,\quad g\in\Gamma_s,
\]
with $a_j\in[0,r_q)^{\Sigma}$ and flux $k\in\Z$. Their curvatures are
\begin{equation}\label{eq:slab-unfolded-curvatures}
 F_j^{\mathrm{el}}=A_{j+1}-A_j-d_sa_j,
 \qquad F_j^{\mathrm{mag}}=d_sA_j+r_qk\mathbf1_{\sigma_0}.
\end{equation}
Thus integration over the original configuration quotient equals
summation over $k,g$ and integration over these paths. For a
nonnegative integrand the identity follows by Tonelli's theorem. For
an absolutely integrable complex integrand it follows by applying
that identity to its modulus and then using Fubini's theorem. Each
change above preserves link phases and curvature, so it applies to
the determinant and to all curvature observables used here.

There are two finite Fock spaces in this representation: $\mathcal F'$
at a coordinate site plane and $\mathcal F$ at a link plane. The
matrix of the open half-slab gives a map
$N(A):\mathcal F'\to\mathcal F$, and we include both crossing weights
by writing
\[
 \widehat N(A)=\omega^{1/2}N(A)(\omega')^{1/2}.
\]
Let $u(a)$ be the diagonal unitary on $\mathcal F$ with entries
$e^{iq\langle a,\mathfrak q(I)\rangle}$. The fermion trace is
\begin{equation}\label{eq:slab-slice-determinant}
 \det \mathcal D_A^{C_{\mathbf n}}=
    \tr_{\mathcal F'}\prod_{j=n_1-1,\ldots,0}
             \bigl(\widehat N(A_{j+1})^*u(a_j)\widehat N(A_j)\bigr).
\end{equation}
This formula also holds on an unfolded path, since
$A_{n_1}-A_0\in\Gamma_s$ leaves $\widehat N$ unchanged.

To derive \eqref{eq:slab-slice-determinant}, split the torus at each
site and link plane. The half-slab above a link plane is the reflected
half-slab below it, hence its matrix is $N(A_{j+1})^*$. The crossing
at a site plane has no link phase and gives $\omega'$. At a link
plane it gives $\omega^{1/2}u(a_j)\omega^{1/2}$. Compose the matrices
and close antiperiodically; the closing parity is the identity.
Cyclicity moves half of the final $\omega'$ to the front, giving
\eqref{eq:slab-slice-determinant}. Every entry of $\widehat N(A)$ is
a polynomial in the link phases: the half-slab action is quadratic
in finitely many generators, its exponential terminates, and
Berezin integration extracts a coefficient. Consequently
\begin{equation}\label{eq:slab-half-matrix-bound}
 C_N:=\sup_{A\in\R^{\Sigma^{(1)}}}\|\widehat N(A)\|<\infty,
 \qquad\widehat N(A+r_qk')=\widehat N(A)
       \quad(k'\in\Z^{\Sigma^{(1)}}).
\end{equation}

On real slice connections the two Hilbert spaces are
\[
 H^s=\bigoplus_{k\in\Z}
      L^2(\R^{\Sigma^{(1)}}/\Gamma_s;\mathcal F'),\qquad
 H^l=\bigoplus_{k\in\Z}
      L^2(\R^{\Sigma^{(1)}}/\Gamma_s;\mathcal F).
\]
All quotient measures are induced by Lebesgue measure. Define commuting
self-adjoint electric operators $E_b$ by
$e^{-i\langle x,E\rangle}\Phi(A)=\Phi(A-x)$. On $H^l$ let
\begin{equation}\label{eq:slab-electric-operators}
 U_G(a)\Phi(A)=u(a)\Phi(A-d_sa),\qquad
 P_G=r_q^{-v}\int_{[0,r_q)^{\Sigma}}U_G(a)\,\dd a,
 \qquad K=r_q^ve^{-|E|^2/2}P_G.
\end{equation}
The integer mode charges make $U_G$ a representation of the compact
gauge group $\R^{\Sigma}/r_q\Z^{\Sigma}$. Its average $P_G$ is the
orthogonal projection onto invariant vectors. It commutes with every
function of $E$, because $u(a)$ acts on the Fock index and all spatial
translations commute. Hence
\begin{equation}\label{eq:slab-electric-square-root}
 K^{1/2}=r_q^{v/2}e^{-|E|^2/4}P_G.
\end{equation}

For a real slice two-cochain $b$, define the multiplication operator
$M_b:H^s\to H^l$ by
\begin{equation}\label{eq:slab-magnetic-multiplier}
 (M_b\Phi)_k(A)=
     e^{-\|d_sA+r_qk\mathbf1_{\sigma_0}+b\|^2/4}
                                      \widehat N(A)\Phi_k(A),
 \qquad X_b=K^{1/2}M_b.
\end{equation}
The magnetic weight is bounded and $\widehat N$ is bounded, so $M_b$
is bounded. We must prove that $X_b$ is Hilbert--Schmidt; the real
quotient is not compact and there are infinitely many flux sectors.

The periodized Gaussian
\[
 \widetilde g_{1/2}([x])
   =\sum_{\gamma\in\Gamma_s}\prod_{b\in\Sigma^{(1)}}
                       \pi^{-1/2}e^{-(x_b+\gamma_b)^2}
\]
has integral one on the quotient and is bounded. Indeed
$\Gamma_s\subset r_q\Z^{\Sigma^{(1)}}$, so the series is dominated
by the product of one-dimensional Gaussian sums on that larger lattice.
It is therefore in $L^1\cap L^\infty\subset L^2$. The kernel of
$K^{1/2}$ is
\[
 K^{1/2}(A',A)=r_q^{-v/2}\int_{[0,r_q)^{\Sigma}}
       \widetilde g_{1/2}([A'-A-d_sa])u(a)\,\dd a.
\]
Its operator norm is bounded by the nonnegative function
\[
 \widehat g([x])=r_q^{-v/2}\int_{[0,r_q)^{\Sigma}}
                  \widetilde g_{1/2}([x-d_sa])\,\dd a,
\]
which has finite $L^1$ and $L^\infty$ norms and hence finite $L^2$ norm.
Integrating the square of the kernel of $X_b$ over $A'$ now gives
\begin{equation}\label{eq:slab-shifted-hs}
 \|X_b\|_2^2\le
    (\dim\mathcal F')C_N^2\|\widehat g\|_2^2
       \sum_{k\in\Z}\int_{\R^{\Sigma^{(1)}}/\Gamma_s}
             e^{-\|d_sA+r_qk\mathbf1_{\sigma_0}+b\|^2/2}\,\dd A.
\end{equation}

To bound this last sum, split the quotient as
\[
 (\ker d_s/\Gamma_s)\times(\ker d_s)^\perp.
\]
The first factor is compact: $\ker d_s$ is the kernel of an integer
matrix, hence has an integer basis, and $\Gamma_s$ has full rank in
it. On the second factor $d_s$ is injective, so
$\|d_sA\|^2$ is positive definite. Moreover
$\ker d_s^*$ consists of constant plaquette functions. This follows
from the two equations
\[
 (d_s^*b)(y,2)=b(y)-b(y-e_3),\qquad
 (d_s^*b)(y,3)=b(y-e_2)-b(y).
\]
Thus $\operatorname{ran}d_s$ is the orthogonal complement of the
constant two-cochain. The orthogonal constant component of
$r_qk\mathbf1_{\sigma_0}+b$ has squared norm
$v^{-1}(r_qk+\sum_\sigma b_\sigma)^2$. Translate the variable in
$(\ker d_s)^\perp$ to remove the remaining component. We obtain
\begin{equation}\label{eq:slab-flux-gaussian}
 \sum_k\int e^{-\|d_sA+r_qk\mathbf1_{\sigma_0}+b\|^2/2}\,\dd A
    =C_s\sum_{k\in\Z}
               e^{-(r_qk+\sum_\sigma b_\sigma)^2/(2v)}<\infty,
\end{equation}
where
$C_s=\operatorname{vol}(\ker d_s/\Gamma_s)
\int_{(\ker d_s)^\perp}e^{-\|d_sA\|^2/2}\,\dd A<\infty$.
This constant is independent of $k$ and $b$. Equations
\eqref{eq:slab-shifted-hs}--\eqref{eq:slab-flux-gaussian} prove the
Hilbert--Schmidt property for every real magnetic perturbation. The
bounds here depend on the fixed cross-section; their purpose is to
justify the trace argument.

For a real external two-cochain $\beta$ on $C_{\mathbf n}$, put
\begin{equation}\label{eq:slab-external-partition}
 Z_{C_{\mathbf n}}(\beta)=\int_{\operatorname{Conf}(C_{\mathbf n})}
          e^{-\|F+\beta\|^2/2}\det \mathcal D_A^{C_{\mathbf n}}\,\dd(A,z).
\end{equation}
Let $\beta_j^{\mathrm{el}}$ and $\beta_j^{\mathrm{mag}}$ be its
electric and magnetic slices. Then
\begin{equation}\label{eq:slab-external-trace}
 Z_{C_{\mathbf n}}(\beta)=(2\pi)^{n_1n_2n_3}
       \operatorname{Tr}_{H^s}\prod_{j=n_1-1,\ldots,0}
            (X_{j+1}^*U_jX_j),
 \quad X_j=X_{\beta_j^{\mathrm{mag}}},\quad
 U_j=e^{i\langle\beta_j^{\mathrm{el}},E\rangle},
 \quad X_{n_1}=X_0.
\end{equation}

Here is the conversion from the unfolded path integral. Write
$g_t(x)=(2\pi t)^{-1/2}\exp(-x^2/(2t))$ for $t>0$. The electric
Maxwell factors are the kernels
\[
 e^{-|A_{j+1}-A_j-d_sa_j+\beta_j^{\mathrm{el}}|^2/2}
       =(2\pi)^v g_1^{\otimes 2v}
                       (A_{j+1}-A_j-d_sa_j+\beta_j^{\mathrm{el}}).
\]
After summing the winding vector they are periodized on $\Gamma_s$.
Gaussian convolution with a translated kernel is
$e^{i\langle\beta_j^{\mathrm{el}},E\rangle}e^{-|E|^2/2}$.
Integration over $a_j$ includes the matrix $u(a_j)$ from
\eqref{eq:slab-slice-determinant}; it is exactly the gauge average in
\eqref{eq:slab-electric-operators}. The magnetic Gaussian in slice
$j$ splits into the two quarter-Gaussians in $M_j$ and $M_j^*$.
The determinant half-slab matrices supply their other factors.
Thus the cyclic path integral is the trace on $H^l$ of
\[
 \prod_{j=n_1-1,\ldots,0}
       (K^{1/2}U_jK^{1/2}M_jM_j^*).
\]
Each time slice contributed $(2\pi)^v$, giving the constant in
\eqref{eq:slab-external-trace}. Move the final $M_0^*$ to the front
and regroup the products as $X_{j+1}^*U_jX_j$. The
Hilbert--Schmidt property above and the absolute Gaussian path
majorant justify this cyclicity and every interchange. In particular
the external-field integral is finite; alternatively this follows
directly by absorbing a linear factor into a smaller Gaussian.

We spell out the path and kernel identification used here. On a
fundamental domain write every intermediate real slice as
$A_j=A_j^\flat+g_j$, $g_j\in\Gamma_s$, and set $g_0=0$ and
$g_{n_1}=A_{n_1}-A_0$. Each unperiodized electric kernel then has
the argument
$A_{j+1}^\flat-A_j^\flat+g_{j+1}-g_j-d_sa_j$.
The map $(g_1,\ldots,g_{n_1})\mapsto
(g_{j+1}-g_j)_{0\le j<n_1}$ is a bijection of $\Gamma_s^{n_1}$.
Summing these differences independently is therefore precisely
periodization of every cyclic kernel. The magnetic factors and
$\widehat N$ are unchanged by the shifts $g_j$. The finite Fock trace
is changed from $\mathcal F'$ to $\mathcal F$ by moving its last
$\widehat N(A_0)^*$ to the front; periodicity identifies it with
$\widehat N(A_{n_1})^*$.

For absolute convergence replace every matrix factor by its norm,
and bound each shifted electric Gaussian by a fixed constant times
$e^{-|A'-A|^2/4}$ and each shifted magnetic factor by a constant
times $e^{-\|F^{\mathrm{mag}}\|^2/4}$. This defines nonnegative
scalar kernels. Their periodized convolutions are bounded and in
$L^2$, and multiplication by the magnetic Gaussian is square
integrable after summing fluxes by \eqref{eq:slab-flux-gaussian}.
Hence each scalar majorant operator is Hilbert--Schmidt. The
integral of their cyclic product is finite: group the first factor
against the product of the others and use the Hilbert--Schmidt
Cauchy--Schwarz inequality. Tonelli first gives this conclusion
for the nonnegative majorants, and Fubini then permits the kernel
expansions and periodization changes for the original integrand.
This proves the trace identity without a formal interchange of
the infinite flux sum and the operator trace.

\subsection{Gaussian domination}\label{subsec:slab-gaussian}

\begin{proposition}\label{thm:slab-gaussian}
For every coordinate torus with sides at least four and every real
two-cochain $\beta$,
\begin{equation}\label{eq:slab-gaussian-partition}
 0<Z_{C_{\mathbf n}}(\beta)\le Z_{C_{\mathbf n}}(0).
\end{equation}
Consequently its normalized curvature law obeys
\begin{equation}\label{eq:slab-gaussian-domination}
 \E_{C_{\mathbf n}}e^{\langle\beta,F\rangle}
                              \le e^{\|\beta\|^2/2}.
\end{equation}
For real two-cochains $\beta_i$ and complex numbers $z_i$,
\begin{equation}\label{eq:slab-complex-gaussian}
 \left|\E_{C_{\mathbf n}}e^{\sum_i z_i\langle\beta_i,F\rangle}\right|
       \le\exp\left(\tfrac12
                    \left\|\sum_i(\operatorname{Re}z_i)\beta_i\right\|^2\right).
\end{equation}
The constants are independent of the volume, the charge, and the mass.
\end{proposition}

\begin{proof}
The determinant is nonnegative and positive on a set of positive
configuration measure, since $Z_{C_{\mathbf n}}(0)>0$. The shifted photon
density is strictly positive everywhere, proving the strict lower
bound in \eqref{eq:slab-gaussian-partition}.

In \eqref{eq:slab-external-trace} move $X_0^*$ through the trace to
obtain a trace on $H^l$:
\[
 (2\pi)^{-n_1n_2n_3}Z_{C_{\mathbf n}}(\beta)
       =\operatorname{Tr}_{H^l}
          (U_{n_1-1}X_{n_1-1}X_{n_1-1}^*\cdots U_0X_0X_0^*).
\]
Every $U_j$ is unitary. Apply Schatten's inequality with all exponents
equal to $n_1$; the trace on the left is nonnegative. Since $X_jX_j^*$
and $X_j^*X_j$ have the same nonzero eigenvalues, this gives
\begin{equation}\label{eq:slab-magnetic-reduction}
 Z_{C_{\mathbf n}}(\beta)
       \le\prod_{j=0}^{n_1-1}
            \bigl((2\pi)^{n_1n_2n_3}
                    \operatorname{Tr}(X_j^*X_j)^{n_1}\bigr)^{1/n_1}
       =\prod_{j=0}^{n_1-1}Z_{C_{\mathbf n}}(\beta^{(j)})^{1/n_1}.
\end{equation}
Here $\beta^{(j)}$ has zero electric part and the same magnetic
part $\beta_j^{\mathrm{mag}}$ in every time slice. The equality uses
the trace formula again, with the same $X_j$ in each slice and all
unitary factors equal to the identity.

Each $\beta^{(j)}$ is supported only on $23$-plaquettes. Interchange
directions $1$ and $2$. The torus becomes $C(n_2,n_1,n_3)$ and the
perturbation is supported only on $13$-plaquettes, so its magnetic
part for the new time direction is zero. The reduction
\eqref{eq:slab-magnetic-reduction} on this torus bounds its partition
function by the unperturbed one. Lattice isometry covariance of the
determinant and photon measure identifies these partition functions
with $Z_{C_{\mathbf n}}(\beta^{(j)})$ and $Z_{C_{\mathbf n}}(0)$. Substitution in
\eqref{eq:slab-magnetic-reduction} proves the upper bound in
\eqref{eq:slab-gaussian-partition}.

Replacing $\beta$ by $-\beta$ and completing the square gives
\eqref{eq:slab-gaussian-domination}. For complex parameters, take
the absolute value inside the probability integral and apply that
bound to $\sum_i(\operatorname{Re}z_i)\beta_i$, proving
\eqref{eq:slab-complex-gaussian}.
\end{proof}


\subsection{Coordinate limits and the reflection form}\label{subsec:coordinate-limits}

\begin{proof}[The coordinate part of \Cref{thm:main}]
The determinant and normalization results make $\mu_n$ a probability
law on $\Omega$ for every integer $n\ge4$. If a finitely supported
cochain $\beta$ embeds in the torus without overlap, its
periodization $\beta_n$ has
\[
 \langle\beta_n,F\rangle=\langle\beta,\widetilde F\rangle,
 \qquad \|\beta_n\|=\|\beta\|.
\]
\Cref{thm:slab-gaussian} therefore gives the claimed Gaussian
bound for this test under $\mu_n$. Every fixed test embeds once
$n$ is large. A single plaquette embeds for every $n\ge4$: its
box of radius $1/2$, enlarged by $[-1,1]^3$, has side three.
Applying the bound with $\beta=\pm R1_\sigma$ and then Markov's
inequality gives
\[
 \mu_n\{|F_\sigma|>R\}\le2e^{-R^2/2},\qquad R>0,
\]
uniformly in $n\ge4$.

Enumerate the plaquettes as $\sigma_1,\sigma_2,\ldots$. Given
$0<\varepsilon<1$, put
\[
 R_j=\sqrt{2\log(2^{j+1}/\varepsilon)},\qquad
 K_\varepsilon=\Omega\cap\prod_{j\ge1}[-R_j,R_j].
\]
The set $\Omega$ is closed in the countable product, so
$K_\varepsilon$ is compact. The union bound gives
$\mu_n(\Omega\setminus K_\varepsilon)\le\varepsilon$ for all
$n\ge4$. Thus the family is tight. Since $\Omega$ is Polish,
Prokhorov's theorem gives subsequential weak limits; see
\cite[Theorem~1.1]{Sorensen}.

Let $\mu_{n_i}\to\nu$ along such a subsequence, and fix $\beta$.
For all sufficiently large $i$, the support of $\beta$ embeds.
For $K>0$ the function
$\min\{e^{\langle\beta,F\rangle},K\}$ is bounded and continuous.
Weak convergence followed by monotone convergence in $K$ gives
\textup{(i)}. Each $\mu_{n_i}$ is invariant under $W$ and charge
conjugation by the finite-volume symmetry identities. The cubic
period lattice and its character, equal to $-1$ on each coordinate
period, are preserved by every signed coordinate permutation.
Translations preserve them as well. Their pullbacks and $F\mapsto-F$
are continuous on $\Omega$, so these symmetries pass to $\nu$.
This proves \textup{(ii)} and reflection invariance at every mirror.

Fix a coordinate mirror $\Pi=\{x_j=c\}$, $c\in\tfrac12\Z$,
and a bounded continuous function $G$ of finitely many plaquettes
wholly in $\{x_j>c\}$. Choose $h_G>0$ so that every corner of
these plaquettes satisfies $c<x_j<c+h_G$. For $n_i>2h_G$, the
support lies in the open positive half-period slab
\[
 \{c<x_j<c+n_i/2\}.
\]
The cube is adapted to this mirror, with closing vector $n_ie_j$
and character $-1$. Its other seam is at $c+n_i/2$, which belongs
to $\tfrac12\Z$ for both even and odd $n_i$. Thus both seams are
admissible, including when one is a site plane and the other a
link plane. The adapted positivity theorem applies and gives
\[
 \mu_{n_i}\bigl(\overline{G(\theta_\Pi^*F)}G(F)\bigr)\ge0.
\]
The integrand is bounded and continuous, so the inequality passes
to $\nu$.

\label{par:side-density}
We give the extension to the whole side, which will also be used
for the diagonal mirrors. Suppose a law $\lambda$ is invariant
under $\theta_\Pi$ and its reflected form is nonnegative on
bounded continuous local tests on one open side. On each
finite-coordinate marginal, compactly supported continuous
functions are dense in $L^2$. Reflection invariance and
Cauchy--Schwarz give
\[
 \left|\lambda\bigl((\Theta A)A-(\Theta B)B\bigr)\right|
 \le\|A-B\|_{L^2(\lambda)}
       \bigl(\|A\|_{L^2(\lambda)}+\|B\|_{L^2(\lambda)}\bigr),
 \qquad \Theta A=\overline{A\circ\theta_\Pi^*}.
\]
Hence positivity extends to all finite-coordinate $L^2$ tests.
Enumerate the plaquettes in the chosen open side, and let
$\mathcal F_m$ be generated by the first $m$. For
$A\in L^2(\Sigma_H,\lambda)$, the conditional expectations
$\lambda(A\mid\mathcal F_m)$ converge to $A$ in $L^2$, because
$\Sigma_H$ is generated by $\bigcup_m\mathcal F_m$. The same
inequality extends positivity to $A$. Reflection carries one
open side to the other, giving the assertion on either side.
Applied to $\lambda=\nu$, this proves the coordinate part of
\textup{(iii)}. The diagonal part remains to be proved by comparison.
\end{proof}

\medskip\noindent
\emph{Optimality of the Gaussian constant.}
Choose a nonzero finitely supported integer two-cochain $z_0$ with
$dz_0=0$; for example, take $z_0=dk$ with $k$ equal to one on a
single oriented link and zero on the other links. Put $\beta=r_qz_0$.
For all sufficiently large $n$, these cochains embed in the cubic
torus. We use the same symbols for their periodizations, which
preserve $dz_0=0$, the source norm, and the pairing with the
periodically extended field. Completing the square gives
\[
\begin{split}
 &\sum_{dz=0}\int
 e^{-\|dA+r_qz\|^2/2+\langle\beta,dA+r_qz\rangle}
       \det\mathcal D_A^{C(n)}\,\dd A\\
 &\qquad =
 e^{\|\beta\|^2/2}\sum_{dz=0}\int
 e^{-\|dA+r_q(z-z_0)\|^2/2}
       \det\mathcal D_A^{C(n)}\,\dd A
 =e^{\|\beta\|^2/2}Z_{C(n)}[1].
\end{split}
\]
Here the sum runs over all closed integer two-cochains and the
integrals have the same domain as in the definition of $\mu_n$.
The last equality reindexes this entire sum; the determinant
depends only on $A$. Thus
$\mu_n(e^{\langle\beta,F\rangle})=e^{\|\beta\|^2/2}$.
Put $X(F)=e^{\langle\beta,F\rangle}$. The same identity for
$2\beta$ gives $\mu_n(X^2)=e^{2\|\beta\|^2}$ for large $n$,
and \textup{(i)} gives $\nu(X^2)\le e^{2\|\beta\|^2}$.
For either of these laws $\lambda$,
\[
 0\le\lambda(X)-\lambda(X\wedge K)
 \le K^{-1}\lambda(X^2)\le K^{-1}e^{2\|\beta\|^2}.
\]
The function $X\wedge K$ is bounded and continuous. First passing
to the weak limit at fixed $K$, then letting $K\to\infty$, gives
\[
 \nu(e^{\langle\beta,F\rangle})=e^{\|\beta\|^2/2}
 \qquad(\beta=r_qz_0).
\]
Since $\beta\ne0$, no constant smaller than one can replace the
constant in \textup{(i)} at the fixed parameters of the theorem.
\section{Reflected perturbations and pressure comparison}\label{sec:diagonal-positivity}

\subsection{The finite-volume fold inequality}\label{subsec:slab-finite-fold}

We give the local form of the comparison that will be used to obtain
diagonal positivity. Fix diagonal cells of even height width $h$,
and let $O$ be a real trigonometric polynomial in finitely many
curvature values at plaquettes contained in the base cell $Q_0$.
Thus $O$ is bounded and depends only on this open cell. Let $n\ge8$
be divisible by $4h$, and abbreviate $D_n=D(n/2,n,n)$. Its cell
count, and that of the cubic torus, is
\[
 N_n=\frac{n^3}{|Q|}=\frac{2n^3}{h^3}.
\]
The translations between cells form the integer lattice
\[
 \Lambda_h=\frac h2\Z e_-+\frac h2\Z e_++h\Z e_3.
\]
On $C(n)$ form the translated copies, and on $D_n$ the reflected
copies:
\begin{equation}\label{eq:slab-two-arrays}
 \widetilde O_n(F)=\sum_{v\in\Lambda_h/n\Z^3}O(\tau_v^*F),
 \qquad O_n^{\mathrm{fold}}(F)=\sum_{\kappa\bmod\Gamma_{D_n}}
                                                O(g_\kappa^*F).
\end{equation}
Each sum has $N_n$ terms. The first sum is read on cubic-periodic
curvatures, the second on $D_n$-periodic curvatures. Here $g_\kappa$
is the reflected map from the base cell to $Q_\kappa$.

Let $\mu_{C(n)}$ and $\mu_{D_n}$ be the normalized curvature laws.
For $t>0$ define the tilted diagonal law
\begin{equation}\label{eq:slab-tilted-law}
 \mu_{D_n}^t(\Phi)=
 \frac{\mu_{D_n}(e^{tO_n^{\mathrm{fold}}}\Phi)}
      {\mu_{D_n}(e^{tO_n^{\mathrm{fold}}})}.
\end{equation}
The folded array is invariant under the group generated by the cell
walls: that group permutes the cells of $D_n$ and their reflected
copies. Thus the tilt retains these symmetries.

\begin{proposition}\label{prop:slab-finite-fold}
For the preceding data and every $t>0$,
\begin{equation}\label{eq:slab-finite-fold}
 \mu_{C(n)}(O)\le\mu_{D_n}^t(O)+\frac{|Q|\Delta(n)}{t}.
\end{equation}
The error is the fixed-parameter pressure error of
\eqref{eq:slab-pressure-error}.
\end{proposition}

\begin{proof}
Jensen's inequality and translation invariance on $C(n)$ give
\begin{equation}\label{eq:slab-fold-lower-tangent}
 \log\frac{Z_{C(n)}[e^{t\widetilde O_n}]}{Z_{C(n)}[1]}
                     \ge t\mu_{C(n)}(\widetilde O_n)
                     =tN_n\mu_{C(n)}(O).
\end{equation}
The diagonal chessboard inequality gives
\begin{equation}\label{eq:slab-fold-change}
 Z_{C(n)}[e^{t\widetilde O_n}]
               \le\|e^{tO}\|_{D_n}^{N_n}
               =Z_{D_n}[e^{tO_n^{\mathrm{fold}}}].
\end{equation}
To check that its $N_n$ factors have the same seminorm, translate
the base observable by $v\in\Lambda_h$. Translation permutes the cells
and conjugates their reflected maps, so the fold of the translated
observable is the translate of the original full fold. The law on
$D_n$ is translation invariant, hence its partition function and
the seminorm are unchanged. All observables here are bounded scalar
functions supported in distinct open cells, so the chessboard
inequality applies directly.

Set $L(t)=\log\mu_{D_n}(e^{tO_n^{\mathrm{fold}}})$. Boundedness
permits differentiation, giving
$L'(t)=\mu_{D_n}^t(O_n^{\mathrm{fold}})$. Convexity and $L(0)=0$
give $L(t)\le tL'(t)$. Wall-group invariance of the tilted law makes
the expectation of every reflected copy equal to that of $O$; hence
\begin{equation}\label{eq:slab-fold-upper-tangent}
 \log\frac{Z_{D_n}[e^{tO_n^{\mathrm{fold}}}]}{Z_{D_n}[1]}
                          \le tN_n\mu_{D_n}^t(O).
\end{equation}
Combining \eqref{eq:slab-fold-lower-tangent}--\eqref{eq:slab-fold-upper-tangent}
leaves, in addition, the normalization difference
\[
 \log Z_{D_n}[1]-\log Z_{C(n)}[1]\le n^3\Delta(n).
\]
Division by $tN_n$ proves \eqref{eq:slab-finite-fold}.
\end{proof}

For fixed cell geometry and $t>0$, the pressure error tends to
zero along admissible sides. To remove the tilt we need uniform local
bounds and enough translation symmetry to place an arbitrary observable
next to a prescribed diagonal wall. We prove the bounds in
\cref{sec:local-amplification}. The layer reduction in
\cref{sec:coordinate-symmetry} then recovers coordinate link reflections
and hence even translations in each comparison state. These are the
symmetries used in the final convex separation.

Fix $q>0$, $0<\rho<1$, positive weights $w_i$ and a real mass $M$.
All volume limits in this section keep these parameters fixed. For the comparison construction we
use sides $m_i\to\infty$ such that eight divides all sufficiently
large sides and the set
\[
 \mathcal D=\{m\in\N:m\mid m_i\text{ for all sufficiently large }i\}
\]
is infinite. The set $\mathcal D$ is closed under divisors and least
common multiples. The comparison sequence may be chosen so that every positive integer
divides all sufficiently large sides; \cref{lem:companions-exist}
constructs such sides below any prescribed original sequence.

\subsection{The tilted law at its walls}\label{subsec:diag-wall-law}

Fix diagonal cells with offsets $(c_0,c_+,c_3)$, even width $h$, and
$c_0,c_+\in\Z$, $c_3\in\Z+\tfrac12$, and $4h\in\mathcal D$.
Let $O$ be the real trigonometric polynomial used in the finite fold
comparison. Its support is contained in a box of radius $r_O$, and
we require $h>8r_O+8$. Such widths can be chosen because
$\mathcal D$ contains arbitrarily large multiples of eight. For
large divisible sides, put $D_i=D(m_i/2,m_i,m_i)$ and
$\mu_i^t=\mu_{D_i}^t$, as in \eqref{eq:slab-tilted-law}. Write
\[
 O(F)=\sum_\beta a_\beta e^{i\langle\beta,F\rangle},
 \qquad |O|_1=\sum_\beta|a_\beta|.
\]
The sum is over the distinct exponents in the chosen finite expansion.
In particular $|O|\le|O|_1$. Let $V_i=O_{m_i}^{\mathrm{fold}}$ and
$N_i=m_i^3/|Q|$. For an admissible Grassmann integrand we use
the same notation for the normalized functional
$\mu_i^t(G)=Z_{D_i}[e^{tV_i}G]/Z_{D_i}[e^{tV_i}]$.
On bosonic integrands it agrees with integration against the
tilted probability law. The normalizing denominator satisfies
\begin{equation}\label{eq:diag-tilt-denominator}
 Z_{D_i}[e^{tV_i}]\ge e^{-tN_i|O|_1}Z_{D_i}[1]>0,
 \qquad t\ge0.
\end{equation}
This uses only the positive bosonic measure; it makes no positivity
assertion about a Grassmann integrand.

At a wall $\Pi=(\ell^a,o_a+jh)$, the adapted torus has height period
$p_a$, with $(p_1,p_2,p_3)=(m_i,2m_i,m_i)$. The number $p_a/h$ is
even. Split the cells into the two half-slabs adjacent to this wall,
and define
\[
 \Phi_+(F)=\exp\left(t\sum_{Q_\kappa\subset
              \{o_a+jh<\ell^a<o_a+jh+p_a/2\}}O(g_\kappa^*F)\right),
\]
with representatives in one half-period. Reflection in the wall
maps these cells onto the other half-period modulo the normal closing
vector. Because the reflected copies of $O$ form one wall-invariant
array, on periodic curvatures
\begin{equation}\label{eq:diag-wall-factorization}
 e^{tV_i}=\Phi_+\Theta_\Pi\Phi_+.
\end{equation}
The function $\Phi_+$ is bounded, real, and supported strictly within
the upper half-slab. The support lies there because each plaquette
of a copy of $O$ is contained in its open cell.

Apply the adapted reflected representation to $\Phi_+G$ and
$\Phi_+G'$. For local even integrands $G,G'$ in the upper half-slab,
\begin{equation}\label{eq:diag-tilted-positive-form}
 \mu_i^t(\Theta_\Pi G'\,G)=
 \frac{\operatorname{Tr}
          (\mathfrak S[\Phi_+G']^*\mathfrak S[\Phi_+G])}
      {Z_{D_i}[e^{tV_i}]}.
\end{equation}
The numerator is the Hilbert--Schmidt inner product of the two slab
operators. Thus, on $D_i$, the tilted functional is reflection positive
at each cell wall for local even integrands supported in the corresponding
open positive half-period slab. Wall reflection also gives
$\mu_i^t(\Theta_\Pi G)=\overline{\mu_i^t(G)}$: the unperturbed
functional is reflection covariant, the scalar weight $e^{tV_i}$
is invariant, and its denominator is real and positive.

The bosonic law is invariant under the full group generated by the
cell walls. Indeed, that group preserves $D_i$ and permutes the
reflected copies in $V_i$, while the unperturbed bosonic law is
invariant under its lattice isometries. In particular the tilt is
invariant under translations by $he_-$, $he_+$ and $2he_3$.
These are the periods of the reflected cell array, rather than all
lattice translations.

We use the Polish closed-cochain space $\Omega$ with its product
topology. The pullback by a lattice reflection or
translation is continuous, so invariance passes to weak limits.
Positivity passes for each bounded continuous local positive-side test
once its support lies in the finite positive half-period slab for all
large volumes. Positivity against bounded Borel local tests follows once reflection invariance
has passed: on their finite-dimensional marginal approximate the
test in $L^2$ by compactly supported continuous functions. The
reflected-form difference is at most
\[
 \|\Phi-\Phi_k\|_2(\|\Phi\|_2+\|\Phi_k\|_2),
\]
because reflection preserves the $L^2$ norm. The Gaussian domination
bound also passes to weak limits, by truncating
$e^{\langle\beta,F\rangle}$ and then using monotone convergence.
The conditional-expectation approximation in
\cref{subsec:coordinate-limits} then extends positivity to the whole
open-side $L^2$ space. These facts will be applied after obtaining tightness of the tilted
laws; their extensive weight alone gives no uniform moment estimate.

\section{Local operators and amplification}\label{sec:local-amplification}

The reflected perturbation contains a number of terms proportional to
the volume. To control one local insertion uniformly, we repeat it
in large diagonal cells and then fold those cells in coordinate
directions. The coordinate trace bounds each repeated operator by
its local norm. Taking the root corresponding to the number of
copies leaves an error that tends to one when volume tends to
infinity before the perturbation is removed. The local norm,
insertion bound and amplification estimate are proved in
\cref{prop:diag-local-norm,prop:diag-coordinate-insertions,thm:diag-amplification}.

The estimate to be proved is
\[
 |\mu_i^t(G)|\le C\gamma_{t,i}(k),\qquad
 \gamma_{t,i}(k)=
 \exp\bigl(2t|O|_1k^3+|Q|k^3\Delta(m_i)\bigr).
\]
Here $C$ is the coordinate bound of the local insertion $G$, defined
in \eqref{eq:diag-coordinate-bound}; for the insertion $G_X$ of
\eqref{eq:diag-operator-insertions} it is $\|X\|$.
The integer $k$ depends on the fixed support
radius, cell width and side sequence, but not on $t$, $i$ or the
location of the support. We choose it before taking limits. The
pressure error $\Delta(m_i)$ tends to zero, so
$\lim_{t\downarrow0}\lim_{i\to\infty}\gamma_{t,i}(k)=1$.
\Cref{thm:diag-amplification} gives the precise assertion and proof.
We first construct the local norm and the insertions to which it applies.

\subsection{Local operators at a coordinate link plane}\label{subsec:diag-local-operators}

An insertion must carry an operator norm that is uniform in the
cross-section and in electric characters. Fix a finite set $S$ of
links of $\Z^2$ and a nonempty finite set $\Lambda$ of sites
containing their endpoints. Let $\mathcal F_\Lambda$ be the Fock
space of crossing modes over $\Lambda$, and put
\[
 H_S=L^2(\R^S)\otimes\mathcal F_\Lambda,
 \qquad E_b=-i\partial_{A(b)}.
\]
The diagonal unitary $u(a)$ on $\mathcal F_\Lambda$ has integer mode
charges. An operator $\Upsilon$ has charge $z\in\Z^\Lambda$ when
$u(a)\Upsilon u(a)^{-1}=e^{iq\langle a,z\rangle}\Upsilon$.

Let $\mathfrak A(S)$ be the span of the functions
$e^{-Q(\xi_{S'})+\langle\lambda,\xi_{S'}\rangle}$, where
$S'\subset S$, $Q$ is real and positive definite on $\R^{S'}$,
and $\lambda\in\C^{S'}$. The empty subset gives the constant one.
These functions are bounded and uniformly continuous. Their span
is closed under products, conjugation and real translations: products
add the positive quadratic forms on the union of their supports,
while translation completes the same quadratic exponent.

For $\delta\in\R^S$, extended by zero to the full slice, require
$d_s^*\delta=qz(\delta)$ with $z(\delta)\in\Z^\Lambda$. The local
operator algebra consists of the finite sums
\begin{equation}\label{eq:diag-operator-normal-form}
 X=\sum_\delta e^{i\langle\delta,A\rangle}\Upsilon_\delta(E_S),
 \qquad\Upsilon_\delta\in\mathfrak A(S)\otimes
                      \operatorname{End}_{z(\delta)}(\mathcal F_\Lambda).
\end{equation}
We denote it by $\mathfrak W(S,\Lambda)$. The Weyl identity gives
\begin{align}\label{eq:diag-operator-rules}
 (e^{i\delta A}\Upsilon(E_S))(e^{i\delta'A}\Upsilon'(E_S))
       &=e^{i(\delta+\delta')A}\Upsilon(E_S+\delta')\Upsilon'(E_S),
          \notag\\
 (e^{i\delta A}\Upsilon(E_S))^*
       &=e^{-i\delta A}\Upsilon^*(E_S-\delta).
\end{align}
Here $\delta A$ abbreviates the real pairing. These formulas make
$\mathfrak W(S,\Lambda)$ a unital $*$-algebra. The normal form is
unique. In Fourier variables its terms send a function supported
near $\xi_0$ to functions supported near $\xi_0+\delta$; small
enough neighborhoods separate the finitely many distinct exponents.
If the sum vanishes as an operator, every continuous symbol vanishes
on each such neighborhood and hence everywhere.

Every $X$ commutes with the local gauge transformations
$\Phi(A)\mapsto u(a)\Phi(A-d_sa)$. The phase change
$e^{-i\langle\delta,d_sa\rangle}$ cancels the Fock charge phase.
The total charge of $z(\delta)$ is zero, since it is the divergence
of a finitely supported cochain. Each Fock term is therefore even.
For $\beta\in\R^S$, define
\begin{equation}\label{eq:diag-alpha}
 \alpha_\beta X=e^{i\langle\beta,A\rangle}Xe^{-i\langle\beta,A\rangle}
    =\sum_\delta e^{i\langle\delta,A\rangle}
                                      \Upsilon_\delta(E_S-\beta).
\end{equation}
This is an isometric $*$-automorphism on $H_S$.

For a rectangular cross-section $\Xi$ containing the smallest integer
rectangle spanned by $\Lambda$,
the same formulas define $X_\Xi$ on the real slice quotient space
$H^l$ of \cref{subsec:slab-coordinate-trace}. The Fock operators
are localized into the full crossing Fock space by $\iota$.
The localization is the tensor inclusion for the exterior space, with the reordering
sign and cross-section character attached to each representative;
it is a unital injective $*$-homomorphism. The phase
$e^{i\langle\delta,A\rangle}$ is well defined on the quotient.
Indeed, on the support hull an integer closed one-cochain is an
integer coboundary, by the box Poincare argument. Thus for
$\gamma\in\Gamma_s$,
$\langle\delta,\gamma\rangle\in2\pi\Z$, using
$d_s^*\delta=qz(\delta)$. The local Weyl formulas therefore remain
valid on the quotient. Uniqueness of the normal form makes the
representation independent of its presentation. It commutes with
$P_G$ and with functions of the electric variables outside $S$.

Define the norm needed for characters by
\begin{equation}\label{eq:diag-sharp-norm}
 \|X\|_\#=\sup_{\Xi,\beta}\| (\alpha_\beta X)_\Xi\|.
\end{equation}
The support hull must embed in the cross-section; the supremum is
over such cross-sections and all $\beta\in\R^S$. The electric spectrum
of a quotient can be smaller than $\R^S$, so its unshifted operator
norm need not equal the local norm.

\begin{proposition}\label{prop:diag-local-norm}
For $X\in\mathfrak W(S,\Lambda)$,
\[
 \|X\|_\#=\|X\|_{\mathcal L(H_S)}.
\]
The norm is unchanged by enlarging $S,\Lambda$. In particular it is
a $C^*$-norm, $\|1\|_\#=1$, and every quotient realization satisfies
$\|(\alpha_\beta X)_\Xi\|\le\|X\|_\#$.
\end{proposition}

\begin{proof}
First remove the Fock localization. Its tensor extension is an
injective $*$-homomorphism and hence isometric; we may work on
$L^2(\R^{\Sigma^{(1)}}/\Gamma_s)\otimes\mathcal F_\Lambda$ with the
original finite Fock operators. Call this realization $X_\Xi^0$.
Introduce an additional variable $z\in\R^S$, with electric operators
$E'=-i\partial_z$, on
\[
 L^2(\R^S_z)\otimes
 L^2(\R^{\Sigma^{(1)}}/\Gamma_s)\otimes\mathcal F_\Lambda.
\]
The operator $X(z,E')$ has norm $\|X\|$, since it is $X$ tensored
with the identity. Choose a Borel representative $A^\flat$ of the
quotient and use the two unitaries
\[
 (U_1\Phi)(z,[A])=\Phi(z+A^\flat|_S,[A]),\qquad
 (U_2\Phi)(z,[A])=\Phi(z,[A-z]),
\]
where $z$ is extended by zero outside $S$. The first conjugation
multiplies $e^{i\delta z}$ by $e^{i\delta A}$. The second replaces
$E'$ by $E'+E_S$ and cancels the factor $e^{i\delta z}$. Hence
\[
 U_2U_1X(z,E')(U_2U_1)^{-1}
       =\sum_\delta e^{i\delta A}\Upsilon_\delta(E_S+E').
\]
Fourier transform in $z$. This becomes a direct integral of
$T_p=(\alpha_{-p}X)_\Xi^0$, $p\in\R^S$. Uniform continuity of the
symbols makes $p\mapsto T_p$ norm continuous. Its direct-integral
norm is therefore $\sup_p\|T_p\|$: the reverse inequality follows
by using a unit vector nearly maximizing $T_{p_0}$ on a small ball
about $p_0$. It follows that
\[
 \|X\|=\sup_p\| (\alpha_{-p}X)_\Xi^0\|
        =\sup_p\| (\alpha_{-p}X)_\Xi\|.
\]
This holds for each admissible cross-section and proves the asserted
supremum identity. Enlarging $S$ tensors with an identity on the
additional $L^2$ factor; enlarging $\Lambda$ uses the same isometric
Fock inclusion. Both preserve the norm.
\end{proof}

\subsection{Insertions and a coordinate trace bound}\label{subsec:diag-insertions}

Let $c=j+\tfrac12$. Write $\operatorname{Ins}_c$ for the crossing
insertion map, and $\omega$ for its diagonal mode weights. For an
even matrix kernel $k(A',A)$ on $\R^S$, set
\begin{equation}\label{eq:diag-kernel-insertion}
 \begin{split}
 G[k](\widehat A)={}&\operatorname{Ins}_c\!\left[
   \omega^{1/2}k(\widehat A_{j+1}|_S,
           (\widehat A_j+d_s\widehat a_j)|_S)
             u(\widehat a_j)\omega^{1/2}\right]\\
 &\quad\cdot\prod_{b\in S}g_1(F_j^{\mathrm{el}}(b))^{-1}
                      \exp(S_c^\times(\Lambda)).
 \end{split}
\end{equation}
Here $S_c^\times(\Lambda)$ is the part of the fermion action crossing
the plane over $\Lambda$. The two factors on the second line remove
the original electric and fermion crossing weights on the support
of $k$. Thus $G[k]$ inserts $k$ in their place. Its definition uses
the infinite lattice and is independent of a torus.

We require
\begin{align}\label{eq:diag-kernel-conditions}
 u(a)^{-1}k(A'+d_sa,A+d_sa)u(a)&=k(A',A),\notag\\
 \|k(A',A)\|&\le C\exp(-\epsilon|A'-A|^2),\qquad\epsilon>0.
\end{align}
Choose a box containing the two copies of $S$, the temporal links
over $\Lambda$, and the unit cubes centered at $(j,y)$ and $(j+1,y)$,
$y\in\Lambda$. In what follows we use the cube $B^\#(c,S,\Lambda)$
obtained by taking twice the largest half-width of the smallest
rectangular box containing these sets, with the same center. This
choice is covariant under coordinate translations and reflections.
The insertion reads only these links and the crossing fields, so
it is supported in $B^\#$.

The insertion is an even, gauge invariant local integrand. Indeed,
an even kernel has entries only between occupation vectors of the
same parity, so every monomial
$\Theta_c\zeta_J\,\zeta_I$ in its insertion is even. If the gauge
parameter is $\lambda$, its upper and lower slice values change
the matrix in \eqref{eq:diag-kernel-insertion} to
\[
 u(\lambda_{j+1})^{-1}
   k(\widehat A_{j+1},\widehat A_j+d_s\widehat a_j)
       u(\widehat a_j)u(\lambda_j).
\]
The matrix entry acquires precisely the inverse of the phase of
$\Theta_c\zeta_J\,\zeta_I$. The electric curvature and the crossing
action are gauge invariant. Finally, the kernel estimate and
$g_1(x)^{-1}=\sqrt{2\pi}e^{x^2/2}$ bound every Grassmann coefficient
by
\begin{equation}\label{eq:diag-insertion-growth}
 C'\exp\left((\tfrac12-\epsilon')
                 \sum_{b\in S}F_j^{\mathrm{el}}(b)^2\right),
 \qquad\epsilon'=\min(\epsilon,\tfrac12).
\end{equation}
The other factors have bounded coefficients. This is the local
integrand bound of \cref{subsec:slab-local}, since its electric
plaquettes lie in the box.

For $X\in\mathfrak W(S,\Lambda)$ let $k_X^{(s)}$ be the continuous
kernel of $e^{-s|E_S|^2}Xe^{-s|E_S|^2}$, with $s=\tfrac14$ or
$\tfrac18$, and write
\begin{equation}\label{eq:diag-operator-insertions}
 G_X=G[k_X^{(1/4)}],\qquad G'_X=G[k_X^{(1/8)}].
\end{equation}
Both kernels satisfy \eqref{eq:diag-kernel-conditions}. To see their
decay, consider one normal-form term $e^{i\delta A}\Upsilon(E)$ with
$\Upsilon(\xi)=e^{-Q(\xi)+\lambda\xi}\Upsilon_0$. Identify $Q$
with its symmetric matrix, so $Q(\xi)=\langle\xi,Q\xi\rangle$,
extending both $Q$ and $\lambda$ by zero to the other links.
The brackets in the following exponent are the complex-bilinear
extension of the real Euclidean pairing. With $Q_s=4sI+2Q$ and
$\lambda_s=\lambda-2s\delta$, its kernel is
\begin{equation}\label{eq:diag-regularized-kernel}
 \begin{split}
 k_X^{(s)}(A',A)={}&e^{i\delta A'}(2\pi)^{-|S|/2}
          (\det Q_s)^{-1/2}e^{-s|\delta|^2}\\
 &\cdot\exp\left(\tfrac12
  \langle\lambda_s+i(A'-A),Q_s^{-1}
                  (\lambda_s+i(A'-A))\rangle\right)\Upsilon_0.
 \end{split}
\end{equation}
The real part of this exponent has a strictly negative quadratic
form in $A'-A$, with a fixed linear term. It gives the required
Gaussian decay after decreasing its exponent. Gauge covariance
follows by conjugating the regularized operator by the gauge
unitary; equality of their continuous kernels gives the first
identity in \eqref{eq:diag-kernel-conditions}. Finite sums preserve
these properties. The identity operator has kernel
$\prod_{b\in S}g_1(A'(b)-A(b))$ when $s=\tfrac14$, so
$G_1=1$; unused links and modes cancel in the same way.

Reflection acts on these insertions by adjoints:
\begin{equation}\label{eq:diag-insertion-reflection}
 \Theta_cG_X=G_{X^*},\qquad\Theta_cG'_X=G'_{X^*}.
\end{equation}
The kernel is changed to $k(A,A')^*$. At the same time the two
slices exchange, $\widehat a_j$ changes sign, and covariance moves
$u(\widehat a_j)$ to the required side of the kernel. This gives
\eqref{eq:diag-insertion-reflection}, including its field signs.
We also need coordinate translations and spatial reflections to
preserve the operator norm, since each coordinate fold contains
such images.

For a spatial reflection $\theta$ in the second or third coordinate,
the spin matrix exchanges the two eigenspaces of $\gamma_1$.
Its unitary matrix $v$ on crossing modes preserves their weights
and reverses their charges. On exterior products the reordering
sign is $r_I=(-1)^{|I|(|I|-1)/2}$; hence the Fock unitary is
$U=(\bigwedge v)r$. The anticommutation of the two field reflection
maps contributes the sign $(-1)^{\deg}$, and moving the two mode
products past one another contributes $(-1)^{|I||J|}$. For the
parity-matched entries of an even kernel these signs leave exactly
$r_Ir_J$. Consequently
\[
 \Theta_\theta\operatorname{Ins}_c[\Upsilon(\widehat A)]
 =\operatorname{Ins}_c[
       U\overline{\Upsilon(\theta^*\widehat A)}U^*].
\]
The antiunitary
$(J\Phi)(A)=U\overline{\Phi(\theta_s^*A)}$ gives this kernel
transformation. It sends an oriented electric operator to the
negative of the reflected electric operator, and sends
$e^{i\delta A}$ to $e^{-i\langle\theta_s^*\delta,A\rangle}$.
Thus it preserves the charge constraint, the local algebra and
the heat factors $e^{-s|E|^2}$. Translations act by the corresponding
unitary relabeling. Together with adjoints at the link plane, these
facts show that each translate or signed coordinate reflection
of $G_X$ is an insertion $G_{X^g}$ with
$\|X^g\|=\|X\|$. The same holds for $G'_X$. The box $B^\#$
transforms with the insertion.

An even local integrand $G$ in a box $B$ satisfies the coordinate
bound $C\ge0$ if, on every rectangular coordinate torus
$T=(\Gamma,\chi)=C(n_1,n_2,n_3)$ with $n_\mu\ge4$,
for every nonempty finite family of images $g_aB$ under lattice
translations and signed coordinate reflections that embeds
disjointly in $\Gamma$, and every $\beta\in C^2(\Gamma)$,
\begin{equation}\label{eq:diag-coordinate-bound}
 \left|Z_T\left[\varphi_\beta\prod_{a=1}^rG^{g_a}\right]\right|
       \le C^rZ_T[1],\qquad
       \varphi_\beta(F)=e^{i\langle\beta,F\rangle},
\end{equation}
provided that, in each coordinate direction $\mu$, there is a
plaquette $\sigma$ orthogonal to $e_\mu$ such that none of the
boxes or their period translates contains any plaquette
$\sigma+k e_\mu$, $k\in\Z$. This condition places the flux
representative used in the slice trace on plaquettes read by
none of the insertions.
When the coordinate cell structure cuts $T$ and its walls are at
half-integers, boxes contained in its open cells satisfy the
free-column condition. Choose a plaquette orthogonal to the chosen
direction that meets a wall normal to another direction.
Period translations permute the open cells, while every plaquette
in this column meets that wall. Thus none of these plaquettes
is contained in a box or any of its period translates.

\begin{proposition}\label{prop:diag-coordinate-insertions}
The insertion $G_X$ satisfies \eqref{eq:diag-coordinate-bound} with
$C=\|X\|$. The insertion $G'_X$ satisfies it with
$C=2^{|S|/2}\|X\|$. A bounded bosonic integrand satisfies it with
$C=\|G\|_\infty$. These bounds survive box enlargement,
multiplication by a curvature character, and transport by a
signed permutation of the coordinates.
\end{proposition}

\begin{proof}
We give the trace calculation, including the support condition.
Choose the time direction to be the first coordinate. By the
covariance just proved, the images of the insertions have the same
operator norms and are at link planes $j+\tfrac12$. The free column
lets us choose the magnetic flux plaquette $\sigma_0$ so that no
plaquette in its column is contained in an insertion box.
At a fixed step $j$ their spatial supports and modes
are disjoint. Let $S_j$ be the union of the supports of the primed
insertions at that step.

The unfolded path variables are genuine local lifts on each
$[j,j+1]\times\operatorname{hull}(\Lambda_a)$. On electric
plaquettes their coboundary is $A_{j+1}-A_j-d_sa_j$. On magnetic
plaquettes it is $d_sA_j$; the flux correction is supported at
$\sigma_0$, which is not a plaquette of the local hull.
Two local lifts with
these curvatures and phases differ by an integer exact cochain
$\frac{2\pi}{q}d\lambda$. Gauge invariance therefore allows the
insertion to be read directly at the path variables. This is why
the free-column hypothesis is needed.

At each crossing, the factors $e^{S_c^\times(\Lambda_a)}$ remove
the crossing weights on the local modes. The remaining crossing
factors, together with their insertions, form
\[
 \operatorname{Ins}_c\left[
   \omega^{1/2}\left(\prod_a\iota(k_a(A_{j+1},A_j+d_sa_j))\right)
                   u(a_j)\omega^{1/2}\right].
\]
The even localized matrices commute because their mode sets are
disjoint. The factors $g_1^{-1}$ similarly remove the electric
Gaussian densities on the local links. The slice determinant
factorization of \cref{subsec:slab-coordinate-trace} now gives
the path product with these modified crossing kernels. Its
Gaussian majorant is integrable: every modified kernel has
strict decay in its electric difference, all remaining links
retain $g_1$, and the magnetic multiplier has its Gaussian
weight. Hence unfolding, periodization and the path trace are
valid for this product.

Separate the curvature character into electric factors
$e^{i\langle\beta_j,F_j^{\mathrm{el}}\rangle}$ and magnetic
unitaries $u_j$. Multiplying a regularized kernel by
$e^{i\langle\beta_j,A'-A\rangle}$ changes the corresponding
operator to
\[
 e^{-s|E-\beta_j|^2}(\alpha_{\beta_j}X)e^{-s|E-\beta_j|^2}.
\]
Fourier transformation of the disjoint link kernels therefore
gives
\begin{equation}\label{eq:diag-mixed-trace}
 Z_T\left[\varphi_\beta\prod_aG_a\right]
 =(2\pi)^{|T|}\operatorname{Tr}_{H^l}
  \prod_{j=n_1-1,\ldots,0}
  \left(K_{\beta_j,S_j}^{1/2}B_jK_{\beta_j,S_j}^{1/2}
                         M_0u_jM_0^*\right),
\end{equation}
where
\[
 B_j=\prod_{a:j_a=j}(\alpha_{\beta_j}X_a)_\Xi,\qquad
 K_{\beta,S}=r_q^v
     e^{-|E-\beta|^2/2+|E-\beta|_S^2/4}P_G.
\]
Here $r_q^v=(2\pi/q)^{|\Sigma|}$ is the gauge-volume factor in
\eqref{eq:slab-electric-operators}; the primed kernels change
the electric Gaussian variance only on $S$. Both the kernels and the local operators
commute with $P_G$. In particular
$\|B_j\|\le\prod_{a:j_a=j}\|X_a\|$ by the norm identity above.

Set $T_j=K_{\beta_j,S_j}^{1/2}M_0$. These are
Hilbert--Schmidt by the same quotient Gaussian calculation as
\eqref{eq:slab-shifted-hs}. Move the first square root cyclically
to the end of the trace. The product then has alternating factors
$B_jT_ju_jT_{j-1}^*$. Schatten H\"older, with $2n_1$ equal
exponents, bounds its absolute trace by
\[
 \prod_j\|B_j\|\prod_j\|T_j\|_{2n_1}^2.
\]
The slice trace applied to $(T_j^*T_j)^{n_1}$ gives exactly
\begin{equation}\label{eq:diag-variance-trace}
 \operatorname{Tr}(T_j^*T_j)^{n_1}
 =(2\pi)^{-|T|}Z_T\left[
    \prod_{l=0}^{n_1-1}e^{i\langle\beta_j,F_l^{\mathrm{el}}\rangle}
    \prod_{l=0}^{n_1-1}\prod_{b\in S_j}
                     \sqrt2e^{-F_l^{\mathrm{el}}(b)^2/2}\right].
\end{equation}
The second product changes $g_1$ into $g_{1/2}$ on those links.
The positive bosonic law bounds the absolute value by
$2^{|S_j|n_1/2}(2\pi)^{-|T|}Z_T[1]$. Consequently
\[
 \prod_j\|T_j\|_{2n_1}^2
 \le 2^{\frac12\sum_j|S_j|}(2\pi)^{-|T|}Z_T[1].
\]
Together with \eqref{eq:diag-mixed-trace} this proves the two
operator bounds. For insertions in another time direction use
an orientation-preserving coordinate permutation. The scalar
bound follows directly from the positive law. Multiplication by
a character merely changes $\beta$, and enlargement restricts
the allowed placements. A signed permutation transports the
torus, the free columns and the placements; conjugation of a
character still has modulus one. These observations prove the
remaining assertions.
\end{proof}

\subsection{Amplification in two cell structures}\label{subsec:diag-amplification}

A diagonal fold exchanges the first and second time directions.
Its copies of a coordinate insertion therefore cannot all be put
in one coordinate trace. We will place these copies in different
coordinate cells and apply a second chessboard estimate. Each
resulting coordinate fold contains copies of only one insertion,
all with the same time direction, so
\eqref{eq:diag-coordinate-bound} applies.

The placement requires two scales depending on the support radius
$r$, but not its center or the tilt. Choose $h''\in h\Z$ and
$k\in\N$ such that
\begin{equation}\label{eq:diag-coarse-widths}
 h''\ge2r+6h,\qquad kh\ge16h'',\qquad
 kh\in h''\Z,\qquad4h'',2kh\in\mathcal D.
\end{equation}
These choices exist: take $4h''$ to be the least common multiple
of $4h$ and an element of $\mathcal D$ exceeding $4(2r+6h)$,
then take $2kh$ to be the least common multiple of $4h''$ and an
element exceeding $32h''$. Fix choices depending only on $h$, $r$
and the side sequence. For all sufficiently large $i$, the needed
divisibility conditions then hold simultaneously.

Here is the geometry of the placement. Let $B=y+[-r,r]^3$. Choose
coarse diagonal offsets $c'_\pm\in c_\pm+h\Z$ (with
$c_-=c_0$) and $c'_3\in c_3+h\Z$ so that
\begin{equation}\label{eq:diag-coarse-center}
 \ell_-(y)-c'_-\in[2h'',2h''+h),\quad
 \ell_+(y)-c'_+\in[4h'',4h''+h),\quad
 y_3-c'_3\in[h'',h''+h).
\end{equation}
The box lies strictly in the coarse cell $Q'_0$ of width $kh$.
Every element of its wall group has the form
\[
 g=\tau_{kh(z_-e_-+z_+e_++2z_3e_3)}
            \theta_-^{a_-}\theta_+^{a_+}\theta_3^{a_3},
 \qquad z\in\Z^3,\quad a\in\{0,1\}^3,
\]
where the reflections are at the lower walls of $Q'_0$. From
\eqref{eq:diag-coarse-center} one obtains
\begin{equation}\label{eq:diag-coordinate-cell-index}
 gy=y+h''\kappa(g)+s(a),\qquad |s(a)|_\infty<2h,
\end{equation}
with
\[
 \kappa(g)=(-2a_--4a_+,\,2a_--4a_+,\,-2a_3)
       +\frac{kh}{h''}(z_-+z_+,\,z_+-z_-,\,2z_3).
\]
The eight first vectors are distinct, and their pairwise
differences have size at most six. Since $kh/h''\ge16$, equality
of two indices first forces equality of their $a$'s and then of
their $z$'s. Thus $g\mapsto\kappa(g)$ is injective.

Choose half-integer coordinate offsets $o_\mu$ in
\[
 I_\mu=(y_\mu+r+2h-h'',\ y_\mu-r-2h).
\]
These intervals have length at least $2h$. Equation
\eqref{eq:diag-coordinate-cell-index} places every $gB$ strictly
inside coordinate cell $\kappa(g)$ of width $h''$. We can also
choose the offsets so that all these coordinate walls avoid
every reflected copy of $\supp O$. Under the original diagonal
wall group, the centers of the copies of a box containing $O$
have at most four residue classes modulo $h$ in either of the
first two coordinates, and at most two in the third. Its radius
$r_O$ therefore forbids at most $4(2r_O+1)=8r_O+4$ half-integer
classes. Since $h>8r_O+8$, an allowed class exists. Each $I_\mu$
contains a representative, and $h''\in h\Z$ keeps every wall
in the allowed class.

For large $i$, the adapted torus $D_i$ carries the original and
coarse diagonal structures; its period lattice lies in the
translation subgroup of the coarse wall group. The coordinate
structure of width $h''$ cuts $D_i$ and is carried by
$T_i=C(m_i/2,2m_i,m_i)$. Formula
\eqref{eq:diag-coordinate-cell-index} also shows that coarse
diagonal cells distinct modulo the periods have their boxes in
coordinate cells distinct modulo those periods. Indeed a period
translation changes $\kappa$ by that translation divided by
$h''$; injectivity then gives the converse. Thus the counts are
\begin{equation}\label{eq:diag-coarse-cell-counts}
 N'_i=\frac{m_i^3}{|Q|k^3},\qquad
 N''_i=\frac{m_i^3}{(h'')^3}.
\end{equation}
The first is the number of coarse diagonal cells in $D_i$; the
second is the number of coordinate cells in either torus. The
shortest periods are at least $m_i/2\ge kh\ge16h''$, so the
boxes, with the margins used for local lifts, embed. Folding
one such box in the coordinate structure gives a disjoint
placement with free columns, since every image stays inside
its open cell.

\begin{proposition}\label{thm:diag-amplification}
Put $\Delta_i=\Delta(m_i)$, with $\Delta$ defined in
\eqref{eq:slab-pressure-error}. It tends to zero as $i\to\infty$,
with the charge, mass and weights fixed.
If $G$ satisfies the coordinate bound $C$ in a box of radius $r$,
then, for $k$ as above and all sufficiently large $i$,
\begin{equation}\label{eq:diag-amplification}
 |\mu_i^t(G)|\le C\gamma_{t,i}(k),\qquad
 \gamma_{t,i}(k)=
      \exp(2t|O|_1k^3+|Q|k^3\Delta_i),\qquad t\ge0.
\end{equation}
The threshold and $k$ depend only on the side sequence, $h$ and
$r$. In particular they are uniform in the location of $G$.
\end{proposition}

\begin{proof}
Each coarse diagonal cell contains $k^3$ original cells. Put
$P=\exp(t\sum_{Q_\kappa\subset Q'_0}O^{[\kappa]})$. The tilt
$e^{tV_i}$ is exactly the product of the reflected copies of $P$
over the coarse cells: their walls are original walls, so each
reflection carries the original array to the same array. Apply
the adapted chessboard estimate to the integrands $PG$ in one
coarse cell and $P$ in the others. The fold of $P$ is $e^{tV_i}$.
After dividing its powers out, this gives
\begin{equation}\label{eq:diag-amplified-power}
 |\mu_i^t(G)|^{N'_i}
 \le\frac{Z_{D_i}[e^{tV_i}\prod_{a=1}^{N'_i}G^{g_a}]}
              {Z_{D_i}[e^{tV_i}]}.
\end{equation}
The right side is nonnegative by the reflected fold form. This
step uses reflection positivity of even integrands and does not
require $G$ itself to be a positive scalar function.

Expand $e^{tV_i}$ in curvature characters. The coefficient sum
is at most $\exp(tN_i|O|_1)$: the character algebra is closed
under multiplication, and the exponential series is absolutely
summable in its coefficient norm. Every plaquette in a character
of this expansion lies strictly within a coordinate cell,
because the chosen walls avoid the support of every copy of $O$.
For each character term, collect its factors by coordinate cells.
Each cell then contains a character alone or a character times
one $G^{g_a}$; precisely $N'_i$ cells have the latter form. Period
translations can be used to place the integrands in a chosen
coordinate brick without changing their value on $D_i$.

Use the chessboard estimate from the adapted to the coordinate
frame, \eqref{eq:slab-chessboard-coordinate}. A cell with a
character alone has fold bounded by $Z_{T_i}[1]$. A cell with
a copy of $G$ has a fold consisting of a character and its
coordinate images; the geometry above supplies the free columns
and disjoint embedding, so its fold is bounded by
$C^{N''_i}Z_{T_i}[1]$. Taking the chessboard roots and multiplying
the cell contributions gives, for every character term,
\[
 \left|Z_{D_i}\left[\varphi_\beta
                    \prod_{a=1}^{N'_i}G^{g_a}\right]\right|
                \le C^{N'_i}Z_{T_i}[1].
\]
The coefficient sum therefore bounds the numerator of
\eqref{eq:diag-amplified-power} by
$e^{tN_i|O|_1}C^{N'_i}Z_{T_i}[1]$. The denominator bound
\eqref{eq:diag-tilt-denominator} contributes a second
$e^{tN_i|O|_1}$. Finally, the pressure comparison gives
$Z_{T_i}[1]/Z_{D_i}[1]\le e^{m_i^3\Delta_i}$.
The root of order $N'_i$ gives
\eqref{eq:diag-amplification}.
\end{proof}

We also need exponential moments, for which a coordinate operator
bound is unnecessary. Repeat \eqref{eq:diag-amplified-power} with
$G=e^{\langle\beta,F\rangle}$ supported in $B$. It is an admissible
local integrand, since
$\langle\beta,F\rangle\le\|\beta\|^2+
\frac14\sum_{\supp\beta}F_\sigma^2$. Its diagonal copies have
disjoint supports. Remove the tilt from the numerator by its
pointwise bound, and apply the second chessboard estimate to
their coordinate cell integrands. The fold of one exponential
has $N''_i$ disjoint copies of $\beta$, and hence squared norm
$N''_i\|\beta\|^2$. Gaussian domination on $T_i$ bounds its
chessboard root by
$e^{\|\beta\|^2/2}Z_{T_i}[1]^{1/N''_i}$. There are $N'_i$
nonempty cells. The same denominator and pressure bounds
therefore yield
\begin{equation}\label{eq:diag-tilted-domination}
 \mu_i^t(e^{\langle\beta,F\rangle})
      \le e^{\|\beta\|^2/2}\gamma_{t,i}(k).
\end{equation}
For each fixed $t$ this proves tightness on $\Omega$. Indeed,
Markov's inequality applied to $\beta=\pm s1_\sigma$ gives a
uniform coordinate tail $2e^{-sR+s^2/2}\gamma_{t,i}(k(h,1))$.
Choose radii $R_\sigma$ with these tails summable over plaquettes;
the intersection of $\Omega$ with their product of compact
intervals is compact. The finitely many initial laws are tight
individually. Any weak limit $\nu_t$ consequently satisfies
\begin{equation}\label{eq:diag-limit-domination}
 \nu_t(e^{\langle\beta,F\rangle})
      \le e^{\|\beta\|^2/2+2t|O|_1k^3}.
\end{equation}
This uses lower semicontinuity of the exponential integral,
as explained after \eqref{eq:diag-tilted-positive-form}.

\section{Recovering coordinate reflection symmetry}\label{sec:coordinate-symmetry}

The local operator estimate first gives reflection symmetry for
observables on the two layers adjacent to a coordinate link plane.
To use this estimate for an arbitrary curvature character, we remove
its outer fermion fields by exact integration identities and its
outer photon exponents by a finite contour shift. Only finitely many
terms of the reflected perturbation change in each shift; their
contribution is of order its strength, uniformly in volume. The
two-layer regularization in \cref{subsec:diag-two-layers}, the exact
field identities in \cref{subsec:diag-fermion-recursion}, and the
contour shift in \cref{subsec:diag-contour} give the reduction in
\cref{thm:diag-polynomial-reduction}.

\subsection{The two adjacent fermion layers}\label{subsec:diag-two-layers}

The remaining comparison argument uses gauge invariant polynomials
\[
 P(\widehat A)=\sum_a b_a e^{i\langle\delta_a,\widehat A\rangle}
                                     \Psi_a,\qquad
 d^*\delta_a=q\mathfrak q(\Psi_a),
\]
where the cochains have finite support and the $\Psi_a$ are ordered
Grassmann monomials. Denote their class by $\mathcal P$. The charge
condition is exactly gauge invariance. It also makes each term
even: the total divergence is zero, so its numbers of $\psi$ and
$\bar\psi$ factors agree. The expansion is unique after equal
pairs $(\delta_a,\Psi_a)$ have been collected, because the monomials
and the distinct characters are linearly independent. Each such
polynomial is an admissible local integrand with bounded coefficients.

Write $v_0=(v_1,v_2,v_3)$, with $0<v_3<v_2<v_1<1/2$, and set
\[
 (s_1,\ldots,s_6)=(-v_1,-v_2,-v_3,v_3,v_2,v_1),\qquad
 \mathcal X_{j,a}=\{x\in\mathcal X:x_1=j+s_a\}.
\]
These are the six thin layers of the thick layer centered at $x_1=j$.
At a coordinate link plane $c=j+\tfrac12$, a two-layer polynomial
reads photon links in the planes $j$ and $j+1$ and between them,
and fermion fields only at the two ends of the bonds crossing
$c$. These are the thin layers $\mathcal X_{j,6}$ and
$\mathcal X_{j+1,1}$; the other thin layers of the two neighboring thick layers do
not occur. Call this subclass $\mathcal P_1(c)$.

Two estimates connect this class to coordinate reflection symmetry.
First, for a monomial $Y\in\mathcal P_1(c)$, heat regularization constructs a
bounded operator $X_{Y,\epsilon}$ and an error
$\xi_Y(\epsilon)\to0$ such that
\[
 |\mu_i^t(\Theta_cY)-\overline{\mu_i^t(Y)}|
 \le4\sqrt{\gamma_{t,i}(k)^2-1}\,\|X_{Y,\epsilon}\|
       +C_Y\gamma_{t,i}(k)\xi_Y(\epsilon).
\]
The fixed support scale $k$ is independent of $\epsilon$.
The operator norm can grow exponentially in $1/\epsilon$; thus the
volume limit must precede $t\downarrow0$, and only then may
$\epsilon\downarrow0$. The construction and the monomial form of
this estimate appear in \eqref{eq:diag-two-layer-finite-error};
finite sums give the symmetry assertion for every two-layer polynomial.
Second, \cref{thm:diag-polynomial-reduction} assigns to every
$P\in\mathcal P$ a $Y_P\in\mathcal P_1(c)$, independent of the
torus and tilt, with
\[
 |\mu_i^t(P^g)-\mu_i^t((Y_P)^g)|\le C_{P,g}t
 \qquad(0\le t\le1,\ i\ge i_0(P,g))
\]
for each fixed lattice isometry $g$. Its proof removes the outer
fermion layers and then shifts the photon contour. This estimate
extends the two-layer symmetry to the curvature characters used
to identify the limiting law.

For a monomial
$Y=e^{i\delta\widehat A}\Psi$, write its exponent as
$(\delta_0,\delta_1,\delta_\perp)$ on the lower spatial, upper
spatial and temporal links. If $\mathfrak q_\pm$ are the lower
and upper charge densities, gauge invariance gives
\begin{equation}\label{eq:diag-two-layer-gauss}
 d_s^*\delta_0-\delta_\perp=q\mathfrak q_-,\qquad
 d_s^*\delta_1+\delta_\perp=q\mathfrak q_+.
\end{equation}
Let $S=\supp\delta_0\cup\supp\delta_1$ and choose a nonempty $\Lambda$
containing their endpoints, the support of $\delta_\perp$ and
all sites over which the fields of $\Psi$ occur.

Every such $Y$ is a kernel insertion, but the operator between
the usual quarter-heat factors need not be bounded. We give the
different regularization that permits approximation by bounded
operators. The crossing modes and their reflected modes form a
basis of the two-layer generators. Since all their weights are
positive when $0<\rho<1$, there is a unique finite matrix
$\Upsilon_Y$ with
\begin{equation}\label{eq:diag-two-layer-matrix}
 \left.\Psi e^{-S_c^\times(\Lambda)}\right|_{\widehat a_j=0}
       =\operatorname{Ins}_c[
                      \omega^{1/2}\Upsilon_Y\omega^{1/2}].
\end{equation}
To obtain it, expand the left side in the basis
$\Theta_c\zeta_J\zeta_I$ and multiply the resulting matrix by
$\omega^{-1/2}$ on both sides. Its nonzero entries have
$\mathfrak q(I)-\mathfrak q(J)=\mathfrak q_-+\mathfrak q_+$.
Thus $\Upsilon_Y$ has charge
$q^{-1}d_s^*(\delta_0+\delta_1)$ and is even.

Set
\begin{equation}\label{eq:diag-two-layer-kernel}
 k_Y^0(A',A)=e^{i\delta_1A'}
             \prod_{b\in S}g_1(A'(b)-A(b))e^{i\delta_0A}
                         \Upsilon_Y.
\end{equation}
It is the continuous kernel of
$e^{i(\delta_0+\delta_1)A}e^{-|E_S+\delta_0|^2/2}\Upsilon_Y$.
We also write $k_Y^0$ for this integral operator.
The charge identity and unchanged Gaussian difference prove
its covariance. Moreover
\begin{equation}\label{eq:diag-two-layer-is-insertion}
 Y=G[k_Y^0].
\end{equation}
Indeed, the Gaussian factors in the insertion cancel. Applying
the automorphism that multiplies the lower fields by their
charge phases to \eqref{eq:diag-two-layer-matrix} replaces its
matrix by $\omega^{1/2}\Upsilon_Yu(\widehat a_j)\omega^{1/2}$
and its left side by
$e^{-iq\langle\mathfrak q_-,\widehat a_j\rangle}
 \Psi e^{-S_c^\times(\Lambda)}$. The remaining temporal exponent
in $G[k_Y^0]$ is
$d_s^*\delta_0-q\mathfrak q_-=\delta_\perp$ by
\eqref{eq:diag-two-layer-gauss}, proving the identity.

Put $\widetilde E=E_S+\delta_0$, $D=\delta_1-\delta_0$ and
$a=(|\delta_0|^2+|\delta_1|^2)/8$. Define the bounded operator
\begin{equation}\label{eq:diag-two-layer-prime}
 X'_Y=e^{i(\delta_0+\delta_1)A}
       e^{-\frac14|\widetilde E|^2+
                         \frac14\langle\widetilde E,D\rangle+a}
                                                      \Upsilon_Y.
\end{equation}
It belongs to $\mathfrak W(S,\Lambda)$, and completion of the
quadratic exponent gives
\[
 e^{-|E_S|^2/8}X'_Ye^{-|E_S|^2/8}=k_Y^0,
 \qquad
 \|X'_Y\|=\|\Upsilon_Y\|e^{|D|^2/16+a}.
\]
For the first identity move the left heat factor through the
phase using the Weyl rule. Its resulting squares are
$\frac18|\widetilde E+\delta_1|^2+
 \frac18|\widetilde E-\delta_0|^2$; their linear and constant
terms cancel those in \eqref{eq:diag-two-layer-prime}. Therefore
$Y=G'_{X'_Y}$ and it satisfies the coordinate bound
$2^{|S|/2}\|X'_Y\|$. This bound will also control the approximation
error uniformly in the torus.

For $\epsilon>0$ let $P_\epsilon=e^{-\epsilon|E_S|^2}$ and
\begin{equation}\label{eq:diag-two-layer-epsilon}
 X_{Y,\epsilon}=e^{i(\delta_0+\delta_1)A}
  \exp\left(-\epsilon|\widetilde E+\delta_1|^2
       -\epsilon|\widetilde E-\delta_0|^2
       +\tfrac12\langle\widetilde E,D\rangle+2a\right)\Upsilon_Y.
\end{equation}
The quadratic part is strictly negative, so this too is in
$\mathfrak W(S,\Lambda)$. The same Weyl calculation gives
\begin{align}\label{eq:diag-two-layer-approximation}
 e^{-|E_S|^2/4}X_{Y,\epsilon}e^{-|E_S|^2/4}
                      &=P_\epsilon k_Y^0P_\epsilon,\notag\\
 Y-G_{X_{Y,\epsilon}}&=
               G'_{X'_Y-P_\epsilon X'_YP_\epsilon}.
\end{align}
The second equality follows from the first and the primed
representation, since their continuous kernels agree. Its
coordinate bound is $2^{|S|/2}\xi_Y(\epsilon)$, where
\begin{equation}\label{eq:diag-two-layer-error}
 \begin{split}
 \xi_Y(\epsilon)={}&\|\Upsilon_Y\|
       \sup_{x\in\R^S}e^{-|x|^2/4+\langle x,D\rangle/4+a}\\
 &\quad\cdot(1-e^{-\epsilon|x+\delta_1|^2
                         -\epsilon|x-\delta_0|^2})
       \le C_Y\epsilon.
 \end{split}
\end{equation}
The last inequality uses $1-e^{-u}\le u$ and the boundedness
of a Gaussian times a quadratic polynomial. By contrast,
\begin{equation}\label{eq:diag-two-layer-operator-size}
 \|X_{Y,\epsilon}\|
 =\|\Upsilon_Y\|
   \exp\left(\frac{(1/2-2\epsilon)^2|D|^2}{8\epsilon}
                  +(1/4-\epsilon)(|\delta_0|^2+|\delta_1|^2)\right).
\end{equation}
It can grow exponentially in $1/\epsilon$. This forces the
order of the limits in the next result.

Reflection changes $(\delta_0,\delta_1,\delta_\perp)$ to
$(-\delta_1,-\delta_0,\delta_\perp)$ and
$\Upsilon_Y$ to $\Upsilon_Y^*$. The latter follows by reflecting
\eqref{eq:diag-two-layer-matrix}; $\operatorname{Ins}_c[\omega]$
is even and invariant. Taking adjoints in the two heat identities
then gives
$X'_{\Theta_cY}=(X'_Y)^*$ and
$X_{\Theta_cY,\epsilon}=X_{Y,\epsilon}^*$; the heat operator
is injective with dense range, so these identities determine
the bounded middle operators uniquely. All constructions extend
linearly to two-layer polynomials. Coordinate permutations
give the same assertions at every coordinate link plane.

\begin{proposition}\label{prop:diag-two-layer-principle}
For every $Y\in\mathcal P_1(c)$ at a coordinate link plane,
\begin{equation}\label{eq:diag-two-layer-principle}
 \lim_{t\downarrow0}\limsup_{i\to\infty}
    |\mu_i^t(\Theta_cY)-\overline{\mu_i^t(Y)}|=0.
\end{equation}
The assertion concerns the finite-volume functionals on fermion
polynomials; it does not presume that they converge to fermion
laws.
\end{proposition}

\begin{proof}
For a fixed support, $o_i^t(X)=\mu_i^t(G_X)$ is a unital linear
functional on the local operator algebra with
$|o_i^t(X)|\le\gamma\|X\|$, where
$\gamma=\gamma_{t,i}(k)$. A unital functional this close in norm
to one is nearly Hermitian. If $X_0=X_0^*$ and $\|X_0\|\le1$,
write $o_i^t(X_0)=u+iv$. For each real $s$,
\[
 (v+s)^2\le |o_i^t(X_0+is)|^2
            \le\gamma^2\|X_0+is\|^2
            \le\gamma^2(1+s^2).
\]
For $\gamma>1$ take $s=v/(\gamma^2-1)$; this gives
$|v|\le\sqrt{\gamma^2-1}$. For $\gamma=1$ the same conclusion
follows by letting $s$ tend to either infinity. Decomposing a
general $X$ into its real and imaginary self-adjoint parts yields
\begin{equation}\label{eq:diag-near-hermitian}
 |o_i^t(X^*)-\overline{o_i^t(X)}|
           \le4\sqrt{\gamma^2-1}\|X\|.
\end{equation}
Use \eqref{eq:diag-insertion-reflection} to interpret this as
reflection covariance of the insertion expectation.

Apply \eqref{eq:diag-near-hermitian} to $X_{Y,\epsilon}$ and
use \eqref{eq:diag-two-layer-approximation} for $Y$ and its
reflection. The amplification theorem controls the two errors
and gives
\begin{equation}\label{eq:diag-two-layer-finite-error}
 \begin{split}
 |\mu_i^t(\Theta_cY)-\overline{\mu_i^t(Y)}|
 \le{}&4\sqrt{\gamma^2-1}\|X_{Y,\epsilon}\|\\
      &+2\cdot2^{|S|/2}\gamma\xi_Y(\epsilon).
 \end{split}
\end{equation}
The support box and its coarse scales do not depend on
$\epsilon$. First take $i\to\infty$, and then $t\downarrow0$
with $\epsilon$ fixed. Since $\Delta_i\to0$, the first term
vanishes and the remaining upper bound is
$2\cdot2^{|S|/2}\xi_Y(\epsilon)$. Let $\epsilon\downarrow0$
using \eqref{eq:diag-two-layer-error}. Finite sums and coordinate
transport give the stated result.
\end{proof}

\subsection{Removing fields from a thin layer}\label{subsec:diag-fermion-recursion}

Consider the six thin layers $\mathcal X_{j,a}$ defined above.
Each site has one neighbor in the next thin layer, one in the
preceding one and two in its own layer. The forward matching
bonds have types $2,1,0,1,2,3$ in succession. This follows
directly from the chamber reflections: type zero reverses
the coordinate of modulus $v_3$, types one and two interchange
the adjacent moduli, and type three reverses the coordinate of
modulus $v_1$ and shifts its nearest photon site by the corresponding
unit vector. The bonds of types zero, one and two carry no photon
phase. Type three bonds within thin layers $2,\ldots,5$ carry
spatial phases in plane $j$, while those between layer six and
the next layer one carry temporal phases between $j$ and $j+1$.
These facts control which planes the recursion can read.

For a photon site $y$, let $\mathcal P^\pm(y)$ consist of the
same exponential monomials as $\mathcal P$, with
$d^*\delta=q(\mathfrak q(\Psi)\mp1_y)$. They transform with
gauge phase $e^{\pm iq\lambda(y)}$ and are odd. Removing a
$\bar\psi(x)$ from an invariant monomial leaves an element of
$\mathcal P^+(\pi(x))$; removing a $\psi(x)$ leaves one in
$\mathcal P^-(\pi(x))$. Multiplication by a hopping phase transports
the specified charge to the other endpoint. These assertions
are the identities
\[
 d^*1_{(y,y')}=1_{y'}-1_y,\qquad
 U_A(x,x')=e^{iq\langle1_{(\pi(x),\pi(x'))},A\rangle}.
\]
They show that every term of the following replacements stays
in $\mathcal P$.

Berezin integration by parts, with left derivatives, gives
\begin{align}\label{eq:diag-equations-of-motion}
 Z_T[\Phi(\bar\psi \mathcal D_A)_\alpha(x')G]
       &=-Z_T[\Phi\partial_{\psi_\alpha(x')}G],\notag\\
 Z_T[\Phi(\mathcal D_A\psi)_\alpha(x')G]
       &= Z_T[\Phi\partial_{\bar\psi_\alpha(x')}G].
\end{align}
The appropriate charge of $G$ makes each integrand invariant.
These identities first hold at fixed complex $A$: the
derivatives of $e^{-S_T(A)}$ are respectively
$(\bar\psi \mathcal D_A)_\alpha e^{-S_T(A)}$ and
$-(\mathcal D_A\psi)_\alpha e^{-S_T(A)}$, and the Berezin integral
annihilates a derivative. Integrating against any bounded
bosonic factor $\Phi(F)$ then proves
\eqref{eq:diag-equations-of-motion}. The enlarged local box
must embed so that all generators involved are distinct torus
generators and the phases can be read on the local lift.

Let $x'$ be the neighbor of $x$ in the next thin layer. Its
matching hopping matrix $w(x,x')$ has eigenvalues
$\frac{w_i}{2}(\rho\mp1)$ and is invertible for our fixed
$0<\rho<1$. Solve the equation at $x'$ for the field at $x$.
With the replaced field moved to the left of the monomial,
the two rules are
\begin{equation}\label{eq:diag-field-replacement}
 \begin{split}
 \bar\psi(x)\longmapsto{}&
 \left(M\bar\psi(x')-
   \sum_{x''\sim x',\,x''\ne x}\bar\psi(x'')w(x'',x')U_A(x'',x')
                  +\partial_{\psi(x')}\right)
                     w(x,x')^{-1}U_A(x',x),\\
 \psi(x)\longmapsto{}&U_A(x,x')w(x',x)^{-1}
 \left(M\psi(x')-
   \sum_{x''\sim x',\,x''\ne x}w(x',x'')U_A(x',x'')\psi(x'')
                  -\partial_{\bar\psi(x')}\right).
 \end{split}
\end{equation}
The derivatives act on the monomial to their right. Multiplying
by the inverse matching matrix and its inverse phase in
\eqref{eq:diag-equations-of-motion} proves equality of the
expectations before and after either replacement, for every
bosonic factor $\Phi$. This equality is independent of a choice
of time direction or a symmetry of the torus.

For a thin layer $\mathbb L$ define a linear map $R^+_{\mathbb L}$
by recursion on the number of its fields in an ordered monomial.
Select the first such field, move it to the front with its
Grassmann sign, and apply \eqref{eq:diag-field-replacement}
through its matching neighbor in $\mathbb L^+$. That neighbor
has only one neighbor in $\mathbb L$, the field just removed.
All other fields created by the rule lie in $\mathbb L^+$ or
$\mathbb L^{++}$; its derivatives remove fields in
$\mathbb L^+$. Thus the number in $\mathbb L$ falls by exactly
one at every recursive call. The map is well-defined, preserves
$\mathcal P$, removes all fields in $\mathbb L$, and creates
no new fields outside its next two layers. A downward map
$R^-_{\mathbb L}$ is obtained using the preceding matching
neighbor.

If the initial polynomial is read in a box $B$, the entire
one-layer recursion is read in $B+[-1,1]^3$. The fields being
removed stay in the original box, and their next neighbors
and the other neighbors of those sites have photon projections
equal or nearest to the original projection: among two successive
different bonds at most one has type three. Hence these sites
and every added phase are in the same enlarged box, irrespective
of the number of recursive calls. Induction therefore gives
\begin{equation}\label{eq:diag-thin-layer-identity}
 Z_T[\Phi P]=Z_T[\Phi R^\pm_{\mathbb L}P]
\end{equation}
whenever this enlarged box embeds. The new photon links are
only those of type three bonds with an endpoint in
$\mathbb L^\pm$. The recursion commutes with multiplication
by a curvature character, since it differentiates only fields
and its choices depend only on the monomial.

We first use this recursion to bound an arbitrary polynomial;
here its photon links need not yet be reduced. Move its fields
upwards through all thin layers below $\mathcal X_{j,6}$ and
downwards through those above $\mathcal X_{j+1,1}$. The last
downward replacement can reach $\mathcal X_{j,6}$, and cannot
reach a lower layer. We obtain, on every sufficiently large
torus and with every bounded bosonic factor, an equality
\begin{equation}\label{eq:diag-polynomial-normal-form}
 Z_T[\Phi P]=\sum_a b_a Z_T[\Phi\varphi_{\beta_a}Y_a],
 \qquad Y_a\in\mathcal P_1(c).
\end{equation}
Here the $Y_a$ can be chosen with integer exponents $qj_a$.
Indeed each remaining monomial has integer charge of total
zero on the two adjacent photon planes. In their connected
rectangular support graph choose paths from a fixed vertex
to each charged vertex; the charge-weighted sum of these
paths is an integer cochain $j_a$ with the required divergence.
Then $u_a=\delta_a-qj_a$ is divergence free in the same
integer box and can be written $d^*\beta_a$ there.

The filling must stay inside the support box for the polynomial
bound below. In
$\prod_\mu[a_\mu,a'_\mu]$, set electric components
$\beta_j^{\mathrm{el}}=-\sum_{l=a_1}^j(u_a)_l$ for
$a_1\le j<a'_1$. Subtracting $d^*\beta$ removes all spatial
links below the top plane; divergence zero then removes
the temporal links. On that top plane a second partial-sum
sweep in the second coordinate removes the third-direction
links below its top edge. Divergence zero removes the
second-direction links and leaves only a terminal segment
in direction three. Its divergence equation and vanishing
boundary value make it zero. The two sweeps use only
plaquettes of the box. This proves the asserted filling
and thus the splitting $\delta_a=d^*\beta_a+qj_a$ in
\eqref{eq:diag-polynomial-normal-form}.

Apply the coordinate bound to each $Y_a$, absorb the extra
character and use amplification. For any fixed box $B_0$
there are constants $C_P$, $k_P$ and a volume threshold with
\begin{equation}\label{eq:diag-polynomial-apriori}
 \sup_{\supp\beta\subset B_0}
       |\mu_i^t(\varphi_\beta P)|
             \le C_P\gamma_{t,i}(k_P),\qquad t\ge0.
\end{equation}
One may take
$C_P=\sum_a|b_a|2^{|S_a|/2}\|X'_{Y_a}\|$, where the factor
$2^{|S_a|/2}\|X'_{Y_a}\|$ is one for $Y_a=1$. Such a term
therefore contributes $|b_a|$. All support boxes in the finite normal form
and $B_0$ fit in one fixed box; its radius determines $k_P$.
The size of $\beta$ does not enter the bound. Reflected
polynomials have the same kind of bound, since the normal
form identity can be reflected on the reflected torus and
the coordinate bounds are preserved by reflection.

\subsection{The contour shift and an outer plane}\label{subsec:diag-contour}

A fermion recursion removes fields from a plane but leaves its
photon exponent. An imaginary shift of links in the next plane
removes that exponent together with the outgoing temporal one.
The all-flux measure is essential for this shift: summation over
exact integer fluxes unfolds the compact link intervals to
real lines on which the contour can move.

Let $u$ be a real cochain supported on finitely many links $S$
of one coordinate plane, and let $P\in\mathcal P$. Choose a
box $B$ reading $P$ and containing $b+[-1,1]^3$ for every
$b\in S$. Put $P^{(u)}(\widehat A)=P(\widehat A-iu)$ by
its entire exponential expansion. Define
\begin{equation}\label{eq:diag-shift-fermion-polynomial}
 \begin{split}
 Q_u=\sum_{b\in S}\sum_{(x,x')\text{ over }b}\big[&
 (e^{qu_b}-1)\bar\psi(x)w(x,x')U_A(x,x')\psi(x')\\
 &+(e^{-qu_b}-1)\bar\psi(x')w(x',x)U_A(x',x)\psi(x)\big],
 \qquad V_u=e^{Q_u}.
 \end{split}
\end{equation}
The bonds are oriented with $b$. Each term is invariant and
even, and the exponential is a finite polynomial by nilpotence.
It reads only those links of $S$ with $u_b\ne0$ and their
overlying bonds. For a plane normal to $e_1$, those fields
are in thin layers $2,\ldots,5$ of its thick layer. The change
in the action is $S_T(A-iu)-S_T(A)=-Q_u$, so
$e^{-S_T(A-iu)}=e^{-S_T(A)}V_u$.

If $\Phi(F)$ is entire and bounded on each tube of bounded
imaginary curvature, then on every torus in which $B$ embeds
we have
\begin{equation}\label{eq:diag-contour-identity}
 Z_T[\Phi P]=e^{\|du\|^2/2}
       Z_T[\Phi(F-idu)\varphi_{du}(F)V_uP^{(u)}].
\end{equation}
The norm here is the lattice norm, equal to the torus norm
because its support embeds. We supply the contour justification.
For each $b$ choose its electric plaquette leaving the plane
in the normal direction. Its only side in $S$ is $b$, with
coefficient $\pm1$. Consequently $d$ is injective on cochains
supported in $S$, and the action
$z\mapsto z+dk$, $k\in\Z^S$, on closed integer fluxes is free.
Choose orbit representatives $z_0$. The compact $A_S$ intervals
and the sum over $k$ become the integral over $\R^S$, since
$(A,z_0+dk)$ represents the same configuration as
$(A+\frac{2\pi}{q}k,z_0)$.

For fixed $z_0$ and the other real links, choose an integer
one-cochain $k'$ with $dk'=z_0$ in the local box. Then
$A+\frac{2\pi}{q}k'$ is a local lift affine in the complex
variables $\zeta=A_S$.
The Berezin coefficients of $e^{-S_T(A)}$ are polynomials
in the link phases, so the full integrand extends to an
entire function $f(\zeta)$. On the strip
$|\Im\zeta_b|\le\|u\|_\infty$ the selected plaquettes give
\begin{equation}\label{eq:diag-contour-strip}
 |f(\zeta)|\le C
          \prod_{b\in S}e^{-(\Re\zeta_b-a_b)^2/2}.
\end{equation}
Indeed each of their real curvatures is $\pm\Re\zeta_b+a'_b$,
all imaginary curvatures are bounded in this finite strip,
the factor $\Phi$ is bounded there, and each remaining
exponential phase is bounded there. The constant can be
chosen independent of the remaining real links and the
orbit representative; only the real Gaussian centers vary.

Move one variable at a time along the rectangle with real
endpoints $\pm R$ and imaginary displacement $-iu_b$. The
vertical sides tend to zero by \eqref{eq:diag-contour-strip}.
Fubini then gives $\int f(\xi)\,\dd\xi=
\int f(\xi-iu)\,\dd\xi$. At the shifted arguments the
Maxwell exponent, using the complex-bilinear cochain pairing, becomes
\[
 -\tfrac12\langle F-idu,F-idu\rangle
       =-\tfrac12\|F\|^2+i\langle du,F\rangle+\tfrac12\|du\|^2.
\]
The change in the fermion action is the polynomial $V_u$,
and the observable becomes $P^{(u)}$. Both the initial
integrand and the resulting real integrand are absolutely
integrable by their local bounds. Hence we may sum the
identities over $z_0$ and integrate the remaining links,
which proves \eqref{eq:diag-contour-identity}.

We now eliminate the lowest plane $j$ read by $P$. This means
that its fields have photon projection height at least $j$
and its links have no endpoint below plane $j$. Apply the
six upward thin-layer maps of this thick layer to obtain
$P^\uparrow=\sum_a b_a e^{i\delta'_a\widehat A}\Psi_a$.
It reads no fields of thick layer $j$ and no links below it;
the only new links are spatial links in plane $j$ and temporal
links from $j$ to $j+1$. At each term, gauge invariance in
plane $j$ therefore says
\begin{equation}\label{eq:diag-lowest-plane-charge}
 (\delta'_a)^\perp_j=d_s^*(\delta'_a)_j.
\end{equation}
For this term let $u_a$ be supported on plane $j+1$ with
$(u_a)_{j+1}=(\delta'_a)_j$. Define
\begin{equation}\label{eq:diag-outer-polynomial}
 P_a=b_a e^{\|du_a\|^2/2+\langle\delta'_a,u_a\rangle}
       e^{i\langle\delta'_a+d^*du_a,\widehat A\rangle}
                                    V_{u_a}\Psi_a.
\end{equation}
Then \eqref{eq:diag-thin-layer-identity} and
\eqref{eq:diag-contour-identity} give the exact relation
\begin{equation}\label{eq:diag-outer-identity}
 Z_T[\Phi P]=\sum_a Z_T[\Phi(F-idu_a)P_a].
\end{equation}
The construction is on the infinite lattice; the same relation
holds for every $\Phi$ entire in the plaquette variables and
bounded on each tube $|\Im F_\sigma|\le r$ for all $\sigma$,
$r>0$, and every torus containing the fixed enlarged support box.

Here is the cancellation of the eliminated exponent. Write
$v=(\delta'_a)_j$. From the slice formulas
$(du)^{\mathrm{el}}_l=u_{l+1}-u_l-d_su_l^\perp$ and
$(d^*\beta)_l=\beta^{\mathrm{el}}_{l-1}
-\beta^{\mathrm{el}}_l+d_s^*\beta^{\mathrm{mag}}_l$, the only
nonzero components of $d^*du_a$ are
\begin{align}\label{eq:diag-outer-cancellation}
 (d^*du_a)_j&=-v,&(d^*du_a)^\perp_j&=-d_s^*v,\notag\\
 (d^*du_a)_{j+1}&=(2+d_s^*d_s)v,
       &(d^*du_a)^\perp_{j+1}&=d_s^*v,\notag\\
 (d^*du_a)_{j+2}&=-v.
\end{align}
Thus the spatial exponent in plane $j$ vanishes, and
\eqref{eq:diag-lowest-plane-charge} cancels its outgoing
temporal exponent as well. The polynomial $V_{u_a}$ reads
only thick layer $j+1$. Hence $P_a$ has lowest plane at
least $j+1$ and highest plane at most the maximum of $j+2$
and the previous highest plane. The six field recursions
enlarge the box by six, and the shifted links with their
plaquettes require two more. The whole step is read in
$B+[-8,8]^3$. Elimination from the highest plane is obtained
by the downward maps and a shift in plane $j-1$. The corresponding
incoming temporal exponent is
$-d_s^*(\delta'_a)_j$, and the same two components cancel.
The highest plane falls by one, while the lowest can fall
only to $j-2$.

The source of error in the tilted law is precisely the shifted
bosonic factor. For a finite cochain $\eta$, only finitely
many copies in the array $V_i$ meet its support. Their change
is the fixed trigonometric polynomial
\begin{equation}\label{eq:diag-array-shift}
 H_\eta(F)=\sum_{\kappa,b}a_b
       (e^{\langle(g_\kappa^{-1})^*\beta_b,\eta\rangle}-1)
                        \varphi_{(g_\kappa^{-1})^*\beta_b}(F),
\end{equation}
where $O=\sum_ba_b\varphi_{\beta_b}$ and only contributing
cells occur. At most $|\supp\eta|$ cells contribute; each
lies in the support box of $\eta$ enlarged by $h$. Its
coefficient norm has the bound
\begin{equation}\label{eq:diag-array-shift-size}
 |H_\eta|_1\le C(\eta):=|\supp\eta|\,|O|_1
        (e^{\max_b\|\beta_b\|_1\|\eta\|_\infty}-1).
\end{equation}
If that enlarged box embeds, no two contributing cells are
period equivalent and the same polynomial is the torus
difference $V_i(F-i\eta)-V_i(F)$. Expand its exponential
in characters:
\begin{equation}\label{eq:diag-shifted-tilt}
 e^{tV_i(F-i\eta)}=e^{tV_i(F)}
        \left(1+\sum_\alpha b_\alpha(t)\varphi_{\beta_\alpha}(F)\right),
 \qquad\sum_\alpha|b_\alpha(t)|\le e^{tC(\eta)}-1.
\end{equation}
The coefficient norm is submultiplicative, so the exponential
series proves absolute uniform convergence. All character
supports lie in one fixed enlarged box, independent of $i$
and $t$; sums of character exponents do not enlarge this support.

Set $P^\flat=\sum_aP_a$ in \eqref{eq:diag-outer-identity}.
Insert \eqref{eq:diag-shifted-tilt} with $\eta=du_a$ and divide
by the positive denominator. Termwise integration is valid:
at a fixed torus the functional on bounded bosonic factors
multiplying $P_a$ is continuous for their supremum norm, by
the absolute integrability of its Berezin coefficients.
We get
\begin{equation}\label{eq:diag-outer-error}
 |\mu_i^t(P)-\mu_i^t(P^\flat)|
 \le(e^{Ct}-1)\sum_a
       \sup_{\supp\beta\subset B'}|\mu_i^t(\varphi_\beta P_a)|,
\end{equation}
where $C=\max_aC(du_a)$ and $B'$ is a fixed box enlarged
by $h$. The a priori bound
\eqref{eq:diag-polynomial-apriori} makes this $O(t)$ uniformly
for all large $i$ and $0\le t\le1$: each of the finitely many
amplification factors is bounded once $\Delta_i\le1$, and
$e^{Ct}-1\le(e^C-1)t$ on this interval.

The same estimate applies after transporting the identity.
This needs care because $D_i$ need not be invariant under
a coordinate permutation. Reflect the equality on the
reflected torus and then reflect it back, using functional
covariance. The conjugations change a shift $\eta$ to
$-\theta^*\eta$. Indeed, if
$\Phi^\theta(F)=\overline{\Phi(\theta^*\overline F)}$, then
$(\Phi^\theta(\,\cdot-i\eta))^\theta=
\Phi(\,\cdot+i\theta^*\eta)$. Thus
\eqref{eq:diag-outer-identity} becomes an exact identity
for $\Theta P$ with shifts $-\theta^*du_a$ and polynomials
$\Theta P_a$ on the original torus. Iterating reflections
transports it by any $g$, with shifts
$\det(\vec g)(g^{-1})^*du_a$, where $\vec g$ is the linear
part of $g$. Its support box is transported at the same time. Apply \eqref{eq:diag-shifted-tilt} and
\eqref{eq:diag-polynomial-apriori} on $D_i$ itself to obtain
the transported $O(t)$ estimate. No invariance of $D_i$
or of its tilted array under $g$ is used.

\begin{proposition}\label{thm:diag-polynomial-reduction}
For each $P\in\mathcal P$ and coordinate link plane $c$ there
is $Y_P\in\mathcal P_1(c)$, independent of the torus and tilt,
such that, for every fixed lattice isometry $g$,
\begin{equation}\label{eq:diag-polynomial-reduction}
 |\mu_i^t(P^g)-\mu_i^t((Y_P)^g)|\le C_{P,g}t
 \quad(0\le t\le1, i\ge i_0(P,g)).
\end{equation}
At a plane normal to either other coordinate the corresponding
reduction and its reflected estimate hold as well.
\end{proposition}

\begin{proof}
For $c=j_0+\tfrac12$, eliminate lowest planes until the
lowest is at least $j_0$. Every step raises it by at least
one, so this takes finitely many steps. Then eliminate
highest planes until the highest is at most $j_0+1$.
During this phase an eliminated plane $j\ge j_0+2$ can
create links only as low as $j-2\ge j_0$, preserving the
first reduction. The resulting polynomial reads just
the two thick layers and their connecting links. Apply
the upward recursions in thin layers $1,\ldots,5$ of
the lower thick layer, followed by the downward recursions
in thin layers $6,\ldots,2$ of the upper one. New fields
remain at or above the lower layer six and at or below
the upper layer one. The new type three phases are
spatial phases in these two planes or the temporal
phases between them. Thus the result $Y_P$ is in
$\mathcal P_1(c)$.

There are finitely many outer-plane steps, each with
\eqref{eq:diag-outer-error} and its transported version;
the final ten thin-layer recursions are exact with the
unchanged bosonic factor and contribute no error. Take
one volume threshold for all their fixed support boxes
and a priori bounds, and sum the finite $O(t)$ errors.
This proves \eqref{eq:diag-polynomial-reduction}.
For another direction first apply the construction to
$P^{g_\mu^{-1}}$ and then transport its exact identities
by the cyclic coordinate permutation $g_\mu$. The general
transported estimate, evaluated on $D_i$, proves the
claim there and its reflected version.
\end{proof}

\section{Comparison states}\label{sec:comparison-states}

\subsection{Comparison states and translation invariance}\label{subsec:diag-comparison-states}

A comparison state for the chosen cell structure and $O$ is
a law $\nu_0$ obtained by the ordered limits
\begin{equation}\label{eq:diag-comparison-state}
 \mu_{i_m}^{t_k}\longrightarrow\nu_{t_k}
       \quad(m\to\infty)\text{ for each }k,
 \qquad
 \nu_{t_k}\longrightarrow\nu_0\quad(k\to\infty),
 \qquad t_k>0,\quad t_k\to0.
\end{equation}
All convergences here are weak convergences of bosonic laws on
$\Omega$, and the volume subsequence is common to all strengths.
Given a prescribed subsequence of volumes, a further subsequence
produces such a state. To see this, take strengths $t_k=1/k$,
extract nested convergent volume subsequences by the tightness
proved above, and choose their diagonal subsequence. The limiting
laws $\nu_{t_k}$ are themselves uniformly tight: applying
\eqref{eq:diag-limit-domination} to $\beta=\pm1_\sigma$ and
using $t_k\le1$ bounds their exponential coordinate moments
by a common constant. The same product-of-intervals argument
gives a weakly convergent strength subsequence. This proves
existence of \eqref{eq:diag-comparison-state} subordinate to
any prescribed volume subsequence.

Every comparison state has Gaussian domination
\begin{equation}\label{eq:diag-comparison-domination}
 \nu_0(e^{\langle\beta,F\rangle})\le e^{\|\beta\|^2/2},
                         \qquad\beta\text{ of finite support}.
\end{equation}
This follows from \eqref{eq:diag-limit-domination} by letting
$t_k\to0$ and truncating the exponential in the second weak
limit. Wall invariance and wall reflection positivity also
pass through both limits. For each local positive-side test,
its support is contained in the positive half-period slab
once $i_m$ is large, so \eqref{eq:diag-tilted-positive-form}
applies. The continuity and $L^2$ extension described in
\cref{subsec:diag-wall-law} give positivity for bounded Borel
local tests. In particular $\nu_0$ is invariant and reflection
positive at $(x_1-x_2,c_0)$.

The comparison with a prescribed cubic limit will be proved in
\cref{thm:comparison-any-side}, using nearby divisible sides.

We now obtain the extra invariance required to use
\cref{thm:comparison-any-side} for observables at arbitrary
positions. Fix a coordinate link reflection $\theta_c$
and a curvature character $P=\varphi_\beta$. Its polynomial
representation is $e^{i\langle d^*\beta,\widehat A\rangle}$;
on every sufficiently large torus this equals
$e^{i\langle\beta,\widetilde F\rangle}$ on the local lift.
Apply \Cref{thm:diag-polynomial-reduction} to $P$
and its reflection, and then
\Cref{prop:diag-two-layer-principle} to its two-layer
reduction. The finite errors are at most two $O(t)$ terms
plus a finite sum of the terms in
\eqref{eq:diag-two-layer-finite-error}. Along the common
volume subsequence in \eqref{eq:diag-comparison-state},
take volume to infinity, then $t_k\to0$ with the regularization
fixed, and finally remove the regularization. Only the
character expectations are passed to weak limits. The
conclusion is
\begin{equation}\label{eq:diag-character-reflection}
 \nu_0(\Theta_c\varphi_\beta)
                           =\overline{\nu_0(\varphi_\beta)}.
\end{equation}
There is no need to extract limiting fermion expectations
in this argument.

For a bosonic character
$\Theta_c\varphi_\beta=\varphi_{-\theta_c^*\beta}$, while
positivity of the probability law gives
$\overline{\nu_0(\varphi_\beta)}=\nu_0(\varphi_{-\beta})$.
Thus \eqref{eq:diag-character-reflection}, with $-\beta$
in place of $\beta$, is invariance under the tensor pullback
$F\mapsto\theta_c^*F$. Characters determine all finite
plaquette marginals by uniqueness of their Fourier
transforms \cite[Theorem~6.2]{Potzelberger}; these marginals determine the law because
cylinder sets generate the Borel sigma algebra of $\Omega$.
Consequently every comparison state is invariant under
all coordinate link reflections, in all three directions.
The product of reflections at $c$ and $c+1$ is translation
by $2e_\mu$, so
\begin{equation}\label{eq:diag-even-translation}
 (\tau_v^*)_*\nu_0=\nu_0,\qquad v\in2\Z^3.
\end{equation}
We need only these even translations. The original wall
periods alone would not have permitted the next placement.

\section{Comparison for arbitrary cubic sides}\label{sec:different-sides}

The comparison states constructed above require increasingly divisible
periods. The original cubic sequence need not have this property.
We compare its side $n$ with a divisible side $m\le n$, using an
empty slab to absorb the difference in each of the three folding
directions. The empty slabs contribute partition functions without
insertions. Their logarithms are controlled by a common pressure
bound, and their total contribution is measured by their volumes
multiplied by the corresponding H\"older exponents.
\Cref{prop:fold-any-side,thm:comparison-any-side} establish the
comparison for the original cubic laws used in \cref{thm:main}.

Throughout this section the parameters retain their fixed values
$q>0$, $0<\rho<1$, $M\in\R$ and $w_i>0$. Write
$D_m=D(m/2,m,m)$ when $m$ is even. The original torus is always
$C(n)$; all tilted laws and comparison states in this section are
formed on the auxiliary tori $D_m$.

\subsection{Unequal slabs}\label{sec:unequal-heights}

Let $\ell(x)=n_\ell\cdot x$ be a mirror family and let
$\Xi=(\Gamma_\parallel,\chi_\parallel)$ be a cross-section. Put
\[
 |\Xi|=[\Z^3\cap\{\ell=0\}:\Gamma_\parallel].
\]
If $P>0$ belongs to $|n_\ell|^2\Z$, set
\[
 u_P=\frac{P}{|n_\ell|^2}n_\ell,\qquad
 \mathbb T(\Xi;P)=
   (\Gamma_\parallel\oplus\Z u_P,\chi_P),
 \qquad \chi_P|_{\Gamma_\parallel}=\chi_\parallel,
 \quad\chi_P(u_P)=-1.
\]
We use this notation only when the displayed period lattice has systole
at least four. Then $u_P$ is an integer vector, the torus is adapted
to $\ell$, and its volume is $P|\Xi|$. More generally, a torus with
this cross-section and height period $p$ has volume $p|\Xi|$.
Indeed, $\ell(\Z^3)=\Z$, and the quotient is counted by its $p$
height classes and the $|\Xi|$ classes in each height. The character
does not affect this count.

\begin{proposition}[Unequal slab folds]\label{prop:unequal-slabs}
Let $\mathbb T=(\Gamma,\chi)$ have cross-section $\Xi$ in the family
$\ell$, and let $g\in\Gamma$ be a closing vector of minimal positive
height $p=\ell(g)$. Choose an integer $k\ge1$ and admissible positions
\[
 c_0<c_1<\cdots<c_{2k}=c_0+p,\qquad h_j=c_j-c_{j-1}.
\]
Assume
\begin{equation}\label{eq:sides-systole}
 \min_{\gamma\in\Gamma_\parallel+
       \Z(2kh_j n_\ell/|n_\ell|^2),\ \gamma\ne0}
       |\gamma|_\infty\ge4,\qquad 1\le j\le2k.
\end{equation}
Let $G_j$ be a bounded scalar function of finitely many curvature
plaquettes contained in the open slab $c_{j-1}<\ell<c_j$.
For each $j$ put $\mathbb T_j=\mathbb T(\Xi;2kh_j)$ and form
$\operatorname{fold}_\ell G_j$ by successive reflections in the
planes $\ell=c_{j-1}+ih_j$. Then
\begin{align}
 Z_{\mathbb T_j}[\operatorname{fold}_\ell G_j]
   &=\|\mathfrak S[G_j]\|_{2k}^{2k}\ge0,
       \label{eq:sides-slab-norm}\\
 \left|Z_{\mathbb T}\left[\prod_{j=1}^{2k}G_j\right]\right|
   &\le\prod_{j=1}^{2k}
       Z_{\mathbb T_j}[\operatorname{fold}_\ell G_j]^{1/(2k)},
       \label{eq:sides-holder}\\
 |\mathbb T|&=\sum_{j=1}^{2k}\frac{|\mathbb T_j|}{2k}.
       \label{eq:sides-volumes}
\end{align}
In particular, the original torus need not admit a straightened torus
of height $p$.
\end{proposition}

\begin{proof}
Admissible positions form the group
$\frac12|n_\ell|^2\Z$. Thus $h_j$ belongs to this group,
$2kh_j\in|n_\ell|^2\Z$, and every plane used in the $j$th fold
is admissible. Condition \eqref{eq:sides-systole} therefore defines
each $\mathbb T_j$. Its volume is $2kh_j|\Xi|$, so summing
$h_j$ proves \eqref{eq:sides-volumes}.

The boundary spaces and the slab operator $\mathfrak S[G_j]$ depend
on the cross-section and the two endpoints, not on the torus used
to close a chain of slabs. Apply the reflected-pair identity
\eqref{eq:slab-fold-holder} on $\mathbb T_j$, whose normal character
is $-1$. Its closing translation cancels the translation of the
repeated pairs, giving
\[
 Z_{\mathbb T_j}[\operatorname{fold}_\ell G_j]
   =\operatorname{Tr}
       (\mathfrak S[G_j]^*\mathfrak S[G_j])^k.
\]
This proves \eqref{eq:sides-slab-norm}.

On the original torus the exact trace \eqref{eq:slab-trace} is
\[
 Z_{\mathbb T}\left[\prod_jG_j\right]
  =\operatorname{Tr}\left(
     V_g^{-1}U_{\mathrm f}^{(1+\chi(g))/2}
     \mathfrak S[G_{2k}]\cdots\mathfrak S[G_1]\right).
\]
Its closing factor is unitary. All slab operators are
Hilbert--Schmidt by \Cref{thm:slab-operators}, hence belong
to the Schatten class of exponent $2k$. Apply
\eqref{eq:slab-schatten-holder} with every exponent equal to $2k$,
and then use \eqref{eq:sides-slab-norm}. This gives
\eqref{eq:sides-holder}. The bounded scalar insertions are read on
periodically extended curvatures throughout; they require no choice
of a fermionic box or local lift.
\end{proof}

For example, when $n$ is odd the first diagonal height period of
$C(n)$ is $n$, whereas a straight normal closing vector would be
$(n/2)e_-$. This is not an integer vector. In the preceding proposition
each replacement normal period is $2kh_j$, which is even for integer
diagonal slab heights. This is the step that permits odd original sides.

\subsection{An empty slab and its pressure cost}\label{sec:empty-slab}

The exceptional torus in an unequal fold may be long. We use a bound
valid for every shape. Set
\begin{equation}\label{eq:sides-constants}
 H_M=|M|+\frac{1+\rho}{2}\sum_{i=0}^3w_i,
 \qquad \bar f=192\log H_M+3\log(r_qC_q).
\end{equation}
Both constants are finite, and $H_M>0$.

\begin{lemma}\label{lem:uniform-ceiling}
Every torus $\mathbb T$ used in the finite model satisfies
\[
 0<Z_{\mathbb T}[1]\le e^{\bar f|\mathbb T|}.
\]
Consequently $f_\infty\le f(a,b,c)\le\bar f$ for all integer
$a,b,c\ge4$.
\end{lemma}

\begin{proof}
The determinant bound is $\det\mathcal D_A^{\mathbb T}
\le H_M^{192|\mathbb T|}$, and the Maxwell integral is at most
$(r_qC_q)^{3|\mathbb T|}$. Together with the positive lower bound
in \eqref{eq:reg-partition-bounds}, these give the first assertion.
The second follows from the definition of pressure and
\Cref{prop:slab-pressure}.
\end{proof}

Use the diagonal cells \eqref{eq:slab-cells}, with heights
$(\ell^1,\ell^2,\ell^3)=(x_1-x_2,x_1+x_2,x_3)$, admissible
offsets $o_a$, and even width $h$. A finite family
$\mathfrak G=(G_\kappa)_{\kappa\in\mathcal K}$ will mean bounded
scalar curvature functions with $G_\kappa$ supported in the open
cell $Q_\kappa$. For a direction $a$ and a layer $j$, write
$\mathfrak R_a^{j,2k}\mathfrak G$ for the family obtained by retaining
the factors with $\kappa_a=j$ and reflecting that layer into the
$2k$ layers with indices $j,\ldots,j+2k-1$. The other two indices
are unchanged. Its product is exactly the slab fold of the original
layer product. This follows because the unique wall-group element
taking one layer to another acts on each of its cells by the
corresponding cell reflection.

\begin{lemma}[One empty slab]\label{lem:one-step}
Fix a direction $a$, a torus $\mathbb T$, and its cross-section $\Xi$
and height period $p$ for $\ell^a$. Suppose $(2k-1)h<p$, and put
\[
 h^{\mathrm{exc}}=p-(2k-1)h,\qquad
 \mathbb T^{\mathrm{reg}}=\mathbb T(\Xi;2kh),\qquad
 \mathbb T^{\mathrm{exc}}=\mathbb T(\Xi;2kh^{\mathrm{exc}}).
\]
Assume both replacement period lattices have systole at least four.
If every index of $\mathfrak G$ satisfies $0\le\kappa_a\le2k-2$,
then
\begin{align}
 \left|Z_{\mathbb T}[\prod\mathfrak G]\right|
 &\le Z_{\mathbb T^{\mathrm{exc}}}[1]^{1/(2k)}
    \prod_{j=0}^{2k-2}
      Z_{\mathbb T^{\mathrm{reg}}}
       [\prod\mathfrak R_a^{j,2k}\mathfrak G]^{1/(2k)},
       \label{eq:sides-step}\\
 |\mathbb T|
 &=\frac{2k-1}{2k}|\mathbb T^{\mathrm{reg}}|
     +\frac1{2k}|\mathbb T^{\mathrm{exc}}|.
       \label{eq:sides-step-volumes}
\end{align}
Every folded partition function on the right is nonnegative.
\end{lemma}

\begin{proof}
Cut at $o_a+jh$ for $0\le j\le2k-1$ and at $o_a+p$.
The first $2k-1$ slabs have height $h$; the final slab has height
$h^{\mathrm{exc}}>0$ and contains no factor. The positions are
admissible, since $p$ and $h$ are integers and $o_a$ is admissible.
Apply \Cref{prop:unequal-slabs} to their layer products
and to the constant $1$ in the final slab. The regular folds are
the products of $\mathfrak R_a^{j,2k}\mathfrak G$, while the empty
slab folds to $1$. Its three conclusions give the assertions.
\end{proof}

The empty slab has no local observable to transport. Thus its reflected
copies may have a different width, while every regular factor remains
in the original cell structure. Folding in direction $a$ leaves the
other two cell indices unchanged, so the same lemma can be applied
successively in the remaining directions.

\subsection{The comparison inequality}

Fix a real local trigonometric polynomial $O$ supported in a box
$y+[-r,r]^3\subset Q_0$, with $r\ge1$. Choose the cell width even,
$h\ge18$, and $h>8r+8$. The offsets are
$o_1=c_0,o_2=c_+,o_3=c_3$, with $c_0,c_+\in\Z$ and
$c_3\in\Z+1/2$. For $4h\mid m$, let $O_m^{\mathrm{fold}}$ be
the full reflected array on $D_m$, and let $\mu_{D_m}^t$ be the
tilt \eqref{eq:slab-tilted-law}. Its number of cells is
$N=m^3/|Q|$, where $|Q|=h^3/2$.

\begin{proposition}\label{prop:fold-any-side}
Suppose $4h\mid m$, $m\ge4h$, and $n\ge m$ is an integer.
Put
\begin{align*}
 V^{\mathrm{reg}}&=\frac{(m-h)^2(2m-h)}2,
 &V^{\mathrm{exc}}&=n^3-V^{\mathrm{reg}},\\
 \Delta'(n,m,h)&=f(m/2,2m,m)-f_\infty
       +\frac{V^{\mathrm{exc}}}{V^{\mathrm{reg}}}
                     (\bar f-f_\infty).
\end{align*}
Then $\Delta'(n,m,h)\ge0$ and, for every $t>0$,
\begin{equation}\label{eq:sides-expectation-comparison}
 \mu_{C(n)}(O)\le\mu_{D_m}^t(O)
                   +\frac{|Q|\Delta'(n,m,h)}t.
\end{equation}
There is no parity restriction on $n$.
\end{proposition}

\begin{proof}
Set $s=n-m\ge0$ and
\[
 k_1=k_3=\frac{m}{2h},\qquad k_2=\frac mh,\qquad
 \mathcal K^{\mathrm{reg}}=
   \{\kappa:0\le\kappa_a\le2k_a-2,\ a=1,2,3\}.
\]
These are integer ranges, with every $k_a\ge2$. Write
\[
 R=\prod_{a=1}^3(2k_a-1),\qquad
 N=8k_1k_2k_3=\frac{m^3}{|Q|}.
\]
Then $R|Q|=V^{\mathrm{reg}}$, and
$0<V^{\mathrm{reg}}<m^3\le n^3$. Since
$f(m/2,2m,m)\ge f_\infty$ and $\bar f\ge f_\infty$,
the claimed sign of $\Delta'$ follows.

\emph{The array on the original torus.}
For $\kappa\in\Z^3$ let
\[
 v(\kappa)=\frac h2(\kappa_1e_-+\kappa_2e_+)
                         +h\kappa_3e_3,\qquad
 G_\kappa=e^{tO^{\tau_{v(\kappa)}}}.
\]
The vector $v(\kappa)$ is integral because $h$ is even, and its
three height coordinates are $h\kappa$. Thus $G_\kappa$ is
supported in $Q_\kappa$. Define
\[
 \widetilde O_{\mathrm{reg}}
     =\sum_{\kappa\in\mathcal K^{\mathrm{reg}}}
                               O^{\tau_{v(\kappa)}},\qquad
 \mathfrak G=(G_\kappa)_{\kappa\in\mathcal K^{\mathrm{reg}}}.
\]
Then $\prod\mathfrak G=e^{t\widetilde O_{\mathrm{reg}}}$.
Jensen's inequality and translation invariance give
\begin{equation}\label{eq:sides-jensen}
 tR\mu_{C(n)}(O)
 \le\log Z_{C(n)}[e^{t\widetilde O_{\mathrm{reg}}}]
                            -\log Z_{C(n)}[1].
\end{equation}

\emph{The three replacement geometries.}
Apply \Cref{lem:one-step} in the order
$(\ell_-,\ell_+,x_3)$. The regular tori are
\[
 \mathbb T_0=C(n),\quad
 \mathbb T_1=D'(m/2,n,n),\quad
 \mathbb T_2=D(m/2,m,n),\quad
 \mathbb T_3=D_m.
\]
For the first step the cross-section of $C(n)$ is generated by
$ne_+,ne_3$, with characters $+1,-1$, respectively. Its height
period is $n$, with closing vector $ne_1$. The regular normal
vector is $k_1he_-=(m/2)e_-$, and the exceptional slab has height
$n-(2k_1-1)h=h+s$. Hence
\[
 \mathbb T_1^{\mathrm{exc}}=D'(k_1(h+s),n,n).
\]
This calculation uses only that $n$ is an integer.

In the second direction, $\mathbb T_1$ has cross-section generated
by $(m/2)e_-,ne_3$, both with character $-1$. Its height period is
$2n$, and the original closing vector $ne_+$ has character $+1$.
The trace therefore contains the parity operator, which is unitary
and is retained in \Cref{prop:unequal-slabs}. The regular
normal vector is $k_2he_+=me_+$, now with character $-1$.
The empty slab has height $h+2s$, giving
\[
 \mathbb T_2^{\mathrm{exc}}=D(m/2,k_2(h+2s),n).
\]
In the third direction, the cross-section of $\mathbb T_2$ is
generated by $(m/2)e_-,me_+$, both with character $-1$, and the
height period is $n$. The regular closing vector is $2k_3he_3=me_3$.
The empty slab has height $h+s$, so
\[
 \mathbb T_3^{\mathrm{exc}}=D(m/2,m,2k_3(h+s)).
\]

All these tori are admissible. For positive integers $a,b,c$,
the lattice generated by $ae_-,be_+,ce_3$ has systole
$\min(a,b,c)$: a vector has coordinates $(aA+bB,-aA+bB,cC)$,
and the maximum of the first two absolute values is
$|aA|+|bB|$. The stated minimum follows, and is attained by a
generator. Here every displayed coefficient is at least
$\min(m/2,2(h+s),n)\ge36$. This checks the regular and exceptional
systoles at each application of \Cref{lem:one-step}.

\emph{The reflected arrays.}
The first fold can be applied because the first indices of
$\mathfrak G$ lie in $[0,2k_1-2]$. After retaining a first layer
$j_1$ and reflecting it, the second and third index ranges are
unchanged. The second fold may therefore be applied to each
resulting family; afterwards the third range is still
$[0,2k_3-2]$, so the third fold applies as well.

For $j=(j_1,j_2,j_3)\in\mathcal K^{\mathrm{reg}}$, the final
family is the full reflected array of $G_j$ in the cells
\[
 j+\mathcal K_0,\qquad
 \mathcal K_0=\{\kappa:0\le\kappa_a\le2k_a-1\}.
\]
Indeed the three wall maps commute, and their product is the
unique wall-group element carrying $Q_j$ to the prescribed cell.
The periods of $D_m$ shift the cell indices by
$(2k_1,0,0)$, $(0,2k_2,0)$ and $(0,0,2k_3)$, so
$j+\mathcal K_0$ is a complete set of cell representatives.
Changing those representatives does not change a scalar fold on
periodic curvatures.

The full fold of $G_j$ has the same partition function as the full
fold of $e^{tO}$. Translation by $v(j)$ permutes the cell walls,
conjugates their reflected maps, and preserves $D_m$. Thus it
translates the full array of $e^{tO}$ into that of $G_j$, and
translation invariance of the finite functional gives
\[
 Z_{D_m}[\operatorname{Fold}_{D_m}G_j]
       =Z_{D_m}[e^{tO_m^{\mathrm{fold}}}].
\]

\emph{The weighted exceptional volume.}
The three H\"older steps give $R$ regular factors, each with exponent
$1/N$, and the exceptional exponents
\[
 \vartheta_1=\frac1{2k_1},\qquad
 \vartheta_2=\frac{2k_1-1}{4k_1k_2},\qquad
 \vartheta_3=\frac{(2k_1-1)(2k_2-1)}{8k_1k_2k_3}.
\]
All folded partition functions are nonnegative. Successive
substitution in \eqref{eq:sides-step} is therefore valid and yields
\begin{equation}\label{eq:sides-three-steps}
 Z_{C(n)}[e^{t\widetilde O_{\mathrm{reg}}}]
 \le\prod_{a=1}^3 Z_{\mathbb T_a^{\mathrm{exc}}}[1]^{\vartheta_a}
       Z_{D_m}[e^{tO_m^{\mathrm{fold}}}]^{R/N}.
\end{equation}
The three volume identities \eqref{eq:sides-step-volumes} are
\[
 |\mathbb T_{a-1}|=
       \frac{|\mathbb T_a^{\mathrm{exc}}|}{2k_a}
        +\frac{2k_a-1}{2k_a}|\mathbb T_a|,
        \qquad a=1,2,3.
\]
Substitute the third into the second, and then the result into the
first. Since $|\mathbb T_3|=m^3$ and $Rm^3/N=R|Q|$, this gives
\begin{equation}\label{eq:sides-exceptional}
 \sum_{a=1}^3\vartheta_a|\mathbb T_a^{\mathrm{exc}}|
       =n^3-R|Q|=V^{\mathrm{exc}}.
\end{equation}
Apply \Cref{lem:uniform-ceiling} separately to the exceptional
factors in \eqref{eq:sides-three-steps}. The resulting bound is
\begin{equation}\label{eq:sides-partition-comparison}
 Z_{C(n)}[e^{t\widetilde O_{\mathrm{reg}}}]
 \le e^{\bar f V^{\mathrm{exc}}}
       Z_{D_m}[e^{tO_m^{\mathrm{fold}}}]^{V^{\mathrm{reg}}/m^3}.
\end{equation}
The weighted identity \eqref{eq:sides-exceptional} is essential here.
The exceptional volumes themselves may exceed $n^3$.

\emph{Normalization and convexity.}
All partition functions in \eqref{eq:sides-partition-comparison}
are positive. For the inserted ones this also follows directly
from $|\widetilde O_{\mathrm{reg}}|\le R|O|_1$ and
$|O_m^{\mathrm{fold}}|\le N|O|_1$.
As in \eqref{eq:slab-fold-upper-tangent}, convexity of the logarithm
of the tilted partition function gives
\[
 \log Z_{D_m}[e^{tO_m^{\mathrm{fold}}}]
      \le\log Z_{D_m}[1]+tN\mu_{D_m}^t(O).
\]
The pressure bounds give
\[
 \log Z_{C(n)}[1]\ge n^3f_\infty,\qquad
 \log Z_{D_m}[1]\le m^3f(m/2,2m,m),
\]
where the second inequality is \eqref{eq:slab-partition-chain}
at the even side $m$. Combining these estimates with
\eqref{eq:sides-jensen} and
\eqref{eq:sides-partition-comparison} yields
\begin{align*}
 tR\mu_{C(n)}(O)
 &\le tR\mu_{D_m}^t(O)
       +V^{\mathrm{reg}}\bigl(f(m/2,2m,m)-f_\infty\bigr)\\
 &\hspace{35mm}+V^{\mathrm{exc}}(\bar f-f_\infty).
\end{align*}
Divide by $tR=tV^{\mathrm{reg}}/|Q|$ to obtain
\eqref{eq:sides-expectation-comparison}.
\end{proof}

\subsection{Divisible companion sides}\label{sec:companion-sequences}

The comparison side must accommodate every fixed cell width used in
the local estimates, including the coarser widths used for
amplification. Rounding to multiples of a single fixed integer would
not supply these widths. The following choice is eventually divisible
by every positive integer and changes the side by less than its square
root.

\begin{lemma}\label{lem:companions-exist}
Let $n_i\ge4$ be integers with $n_i\to\infty$. There are positive
integers $m_i$ such that $m_i\le n_i$ eventually, $m_i/n_i\to1$,
and every positive integer divides $m_i$ for all sufficiently large
$i$. One choice, for $n_i\ge64$, is
\[
 k(n)=\max\{k\in\Z:k\ge1,\ 8k!\le\sqrt n\},\qquad
 m_i=8k(n_i)!\left\lfloor\frac{n_i}{8k(n_i)!}\right\rfloor.
\]
The finitely many remaining terms may be set equal to $8$.
\end{lemma}

\begin{proof}
For $n\ge64$ the set defining $k(n)$ is nonempty and finite.
The remainder after division by $8k(n)!$ satisfies
\[
 0\le n-8k(n)!\left\lfloor\frac n{8k(n)!}\right\rfloor
       <8k(n)!\le\sqrt n.
\]
Thus $n_i-\sqrt{n_i}<m_i\le n_i$ for all large $i$, which gives
$m_i/n_i\to1$ and $m_i\to\infty$. Given an integer $d\ge1$,
the inequality $n_i\ge64(d!)^2$ holds eventually. It implies
$k(n_i)\ge d$, so $d$ divides $k(n_i)!$ and hence $m_i$.
\end{proof}

\begin{proposition}[Comparison for arbitrary original sides]
\label{thm:comparison-any-side}
Let $n_i\ge4$ tend to infinity, and choose $m_i$ as in
\Cref{lem:companions-exist}. Fix the cell data and observable
$O$ of \Cref{prop:fold-any-side}. Form the tilted laws
$\mu_i^t=\mu_{D_{m_i}}^t$ and their ordered comparison states using
these companion sides. Then
\[
 \Delta'_i=\Delta'(n_i,m_i,h)\longrightarrow0.
\]
If $\mu_{C(n_i)}\Rightarrow\nu$ on a prescribed subsequence and
$\nu_0$ is any comparison state subordinate to that subsequence,
then
\begin{equation}\label{eq:sides-limit-comparison}
 \nu(O)\le\nu_0(O).
\end{equation}
\end{proposition}

\begin{proof}
For all sufficiently large $i$, one has $4h\mid m_i$, $m_i\ge4h$
and $m_i\le n_i$. \Cref{prop:fold-any-side} therefore
applies. The pressure term $f(m_i/2,2m_i,m_i)-f_\infty$ tends
to zero by \Cref{prop:slab-pressure}. Moreover
\begin{equation}\label{eq:sides-volume-ratio}
 \frac{V_i^{\mathrm{reg}}}{n_i^3}
 =\left(\frac{m_i}{n_i}\right)^3
    \left(1-\frac h{m_i}\right)^2
    \left(1-\frac h{2m_i}\right)\longrightarrow1.
\end{equation}
Hence $V_i^{\mathrm{exc}}/V_i^{\mathrm{reg}}\to0$, proving
$\Delta'_i\to0$.

All divisibility requirements used to construct comparison states
hold for $(m_i)$. Choose the common volume subsequence $i_j$ and
strengths $t_k\downarrow0$ in \eqref{eq:diag-comparison-state},
subordinate to the prescribed original subsequence. For fixed $k$,
\eqref{eq:sides-expectation-comparison} gives
\[
 \mu_{C(n_{i_j})}(O)\le\mu_{i_j}^{t_k}(O)
                         +\frac{|Q|\Delta'_{i_j}}{t_k}.
\]
The observable $O$ is bounded and continuous. Letting $j\to\infty$
at fixed $t_k>0$ therefore gives $\nu(O)\le\nu_{t_k}(O)$.
Now let $k\to\infty$ to obtain \eqref{eq:sides-limit-comparison}.
The factor $1/t_k$ in the error is the reason for this order of limits.
\end{proof}

The law on the left of \eqref{eq:sides-limit-comparison} remains the
limit of the original cubic laws. The construction of $\nu_0$ uses
the companion volumes and retains the wall positivity and even
translation invariance proved for ordered comparison states.
\section{Reflection positivity of the original limit}\label{sec:completion}

\subsection{Convex separation and the nine mirror families}\label{subsec:diag-separation}

Fix an integer $c_0$. Let $\mathcal R(c_0)$ be the set of
laws on $\Omega$ satisfying Gaussian domination and both
invariance and reflection positivity at
$(x_1-x_2,c_0)$. This is convex and weakly compact. Gaussian
domination gives a common tail
$\lambda(|F_\sigma|\ge R)\le2e^{-R^2/2}$ and hence uniform
tightness. It is weakly closed by truncation of its exponential
integrals. Mirror invariance and positivity against continuous
local tests are also closed, and the approximation in \cref{subsec:coordinate-limits}
extends the latter to the full open-side $L^2$ space. Each condition is affine
or an affine inequality in the law, proving convexity.

We claim that, for every real local trigonometric polynomial
$O$ and every cubic limit $\nu$ along an arbitrary sequence of sides tending to infinity, some $\nu_0\in\mathcal R(c_0)$ satisfies
\eqref{eq:diag-dominating-state}. The polynomial need not
initially fit in a cell having its lower wall at $c_0$.
Let $\supp O\subset y_0+[-r,r]^3$, with $r\ge1$, and choose
an admissible even width $h>8r+8$. Choose an even integer
$a$ in
\[
 (c_0+2r-\ell_-(y_0),\ c_0+h-2r-\ell_-(y_0)).
\]
The interval has length greater than two. For $v=ae_1$ and
$y=y_0+v$ its defining inequalities place the translated
box between the two $\ell_-$ walls. Choose integer $c_+$
and half-integer $c_3$ so that
\[
 c_+<\ell_+(y)-2r<\ell_+(y)+2r<c_++h,
 \qquad c_3<y_3-r<y_3+r<c_3+h.
\]
The available intervals have length greater than one.
The box therefore lies in the open diagonal cell, since
$|\ell_\pm(x)-\ell_\pm(y)|\le2r$ on it. Form the tilted
law with the translated polynomial $O^{\tau_v}$.

Choose companion sides by \cref{lem:companions-exist} and form a
comparison state on those sides, subordinate to the original volume
subsequence defining $\nu$. It belongs to $\mathcal R(c_0)$,
and \cref{thm:comparison-any-side} gives
$\nu(O^{\tau_v})\le\nu_0(O^{\tau_v})$. The original cubic
laws are invariant under every lattice translation, and
this invariance passes to $\nu$. Equation
\eqref{eq:diag-even-translation} gives the same invariance
under this particular $v$ for $\nu_0$. Hence
\begin{equation}\label{eq:diag-dominating-state}
 \nu(O)=\nu(O^{\tau_v})
       \le\nu_0(O^{\tau_v})=\nu_0(O).
\end{equation}
The comparison state may depend on $O$. Convex separation
is what turns these individual inequalities into reflection
positivity of the original law.

\begin{lemma}\label{lem:diag-separation}
Let $\mathcal R$ be a convex weakly compact set of laws on
$\Omega$. If for every real local trigonometric polynomial
$O$ some $\lambda\in\mathcal R$ satisfies
$\nu(O)\le\lambda(O)$, then $\nu\in\mathcal R$.
\end{lemma}

\begin{proof}
Give the real vector space of finite signed measures on
$\Omega$ the weak topology defined by integration against
these polynomials. It is locally convex and Hausdorff:
if every such integral of a signed measure vanishes, its
finite-coordinate marginals have zero Fourier transform,
so its positive and negative parts agree on cylinder sets
and hence on every Borel set. Its continuous linear
functionals are precisely integration against the test
space; a continuous linear functional for a topology
defined by these seminorms depends on finitely many of
them, and their finite real span is still that space.
Since the tests are bounded and continuous, the inclusion
from the usual weak topology of laws is continuous.
Thus $\mathcal R$ is a compact convex subset of this
locally convex Hausdorff vector space. It is nonempty
by the hypothesis at $O=0$.

If $\nu$ were outside $\mathcal R$, strict separation
\cite[Theorem~6.2.9]{Peterson2017} of
a compact convex set and a disjoint closed point would
give a continuous linear functional with
$\max_{\lambda\in\mathcal R}\lambda(O)<\nu(O)$ for some
test $O$. Its maximum exists by compactness. This contradicts
the assumed inequality for that $O$.
\end{proof}

\begin{proof}[Proof of \Cref{thm:main}]

The coordinate part of the theorem, including tightness, Gaussian
domination, all lattice symmetries and coordinate reflection positivity,
was proved in \cref{subsec:coordinate-limits} for arbitrary integer
sides. Let $\nu$ be any limit along a prescribed original sequence.
For the diagonal mirror $(x_1-x_2,c_0)$,
\eqref{eq:diag-dominating-state} holds for every real
trigonometric polynomial. \Cref{lem:diag-separation}
therefore gives $\nu\in\mathcal R(c_0)$, proving its
diagonal positivity. This is valid for every integer
$c_0$.

A signed permutation transports this diagonal mirror
and its positive half-space to any plane
$x_\mu\pm x_\nu=c$. The corresponding reflection is the
conjugate of the original one. Pull a positive-side test
back by this isometry; invariance of $\nu$ identifies
its reflected form with the already nonnegative one.
This gives the six diagonal families and completes
the proof.
\end{proof}

\printbibliography
\end{document}
